% ============================================================
%  Hans Brøchner, "Lectures on Søren Kierkegaard"
%  [Forelæsninger om Søren Kierkegaard]
%  University of Copenhagen, autumn semester 1858 ("Forelæsninger over den
%  nyere Philosophies Forhold til Christendommen") and spring semester 1859.
%  TRANSLATION (English).
%
%  ---- SOURCE ----
%  Translated from the reading text edited by Carl Henrik Koch, Hans
%  Brøchners bidrag til Kierkegaard-receptionen i Danmark (Scientia Danica,
%  Series H, Humanistica 8, vol. 13; Copenhagen: Royal Danish Academy of
%  Sciences and Letters, 2016), pp. 72--188 -- via a Danish working text
%  made after it (koch-text.tex in this directory, excluded from the
%  repository, since Koch's edition is in copyright; its header records how
%  it was established against the page images).
%  The translation is of BRØCHNER'S TEXT ONLY.  Koch's introduction (pp.
%  9--44), manuscript description and editorial principles (pp. 67--71) and
%  commentary (pp. 189 ff.) are not translated.  Koch's editorial additions
%  within the lectures appear only as bracketed references marked "ed.":
%  \ed{...} = his references (the manuscript's gathering.page, e.g. [ed. 1.2];
%  SKS locations; the doubtful-reading mark ?); \ednote{} = the place of one
%  of his editorial notes (descriptions of the manuscript and of inserted
%  leaves, summaries of lectures he does not print), which is not
%  translated.  Letters and words Koch supplies inside Brøchner's sentences
%  (expanded abbreviations, K{ierkegaard}; commas "for læselighedens skyld")
%  are rendered silently as part of the text.  Koch's footnotes keep his
%  numbers (4--149, continuous from his introduction); they carry what
%  Brøchner wrote in the right-hand column of his manuscript, except for the
%  few that are Koch's (bracketed "ed.").  \ins{...} = a right-column addition
%  that Koch prints in the text between slashes, printed / ... /.
%  THE MANUSCRIPT (Det Kgl. Bibliotek, Add. 414, Brøchner's lecture
%  manuscripts 1853--75) HAS NOT BEEN COLLATED.  Every reading here is
%  Koch's.
%  Method: ../../../TRANSLATION-PLAYBOOK.md.  No published translation of
%  Kierkegaard or of Brøchner was consulted, and SKS was not consulted.
%
%  ---- REGISTER ----
%  Scholarly, moderately literal, readable English; American spelling.
%  These are lecture notes: Brøchner's long, nominalized periods, his note
%  form (fragments, lists, recaps, "cf." asides) and his emphasis are kept;
%  fragments are not smoothed into sentences.  Long periods are broken only
%  where English will not carry them, never across a page marker or a
%  footnote anchor.
%
%  ---- QUOTATIONS ----
%  Brøchner quotes and paraphrases Kierkegaard throughout.  Every quotation
%  is translated from Brøchner's Danish as Koch prints it -- not from, or
%  towards, Kierkegaard's text or any published English translation (Hong,
%  Lowrie, Swenson, Hannay): Brøchner abbreviates and alters, and the
%  alterations are his.
%
%  ---- TERMINOLOGY ----
%  Aligned with this repo's other Brøchner translations:
%    Tro / Viden = faith / knowledge;  Troesvidenskab = science of faith;
%    Erkjendelse = cognition;  Videnskab = science (Wissenschaft);
%    Aand = spirit;  Existents = existence (Tilværelse = existence too,
%    commented where both stand together);  Inderlighed = inwardness,
%    Inderliggjørelse = inward deepening, Indvorteshed = interiority;
%    det Religieuse / Religiøse = the religious, Religiøsitet = religiousness
%    (Religiousness A / B);  den Enkelte = the single individual;
%    Forargelse = offense;  Opbyggelse = upbuilding;  Livsanskuelse = view
%    of life;  Christendom / Christenhed = Christianity / Christendom;
%    det Humane = the humane (the sphere of the universally human, set
%    against the Christian; det Menneskelige = the human).
%  Terms harmonized across batches after translation: Virkelighed =
%  actuality;  Bestemmelse = determination;  Tilegnelse = appropriation;
%  Afgjørelse = decision, Beslutning = resolution;  Opfattelse = conception;
%  Gebet = domain;  Guden = the god (Gud = God);  Salighed = happiness, evig
%  Salighed = eternal happiness (salig = blessed);  Lykke = good fortune;
%  Forbillede = prototype;  Efterfølgelse = imitation;  Arvesynd =
%  hereditary sin;  Anfægtelse = spiritual trial;  Sandhedsvidne = witness
%  to the truth;  Tilblivelse / Tilbliven = coming into existence, Vorden =
%  becoming;  Spændkraft = tensile force;  Misforhold = misrelation for the
%  self's relation to itself, otherwise disproportion;  ophæve = sublate in
%  Hegel's sense, otherwise annul;  Forestilling = representation (Hegel's
%  Vorstellung) or idea, by context.  Further terms of art are recorded in
%  "% TR:" comments at their first occurrence.
%  Kierkegaard's works take their standard English titles (Either/Or, Fear
%  and Trembling, Philosophical Fragments, The Concept of Anxiety, Stages on
%  Life's Way, Concluding Unscientific Postscript, Two Ethical-Religious
%  Essays, The Sickness unto Death, Practice in Christianity, The Moment,
%  What Christ Judges of Official Christianity ...), in the Danish's markup.
%  All other titles stay in their original language -- including
%  Bladartikler, the 1857 volume of Kierkegaard's newspaper articles
%  (S. Kierkegaards Bladartikler), which Brøchner cites by page.
%
%  ---- CONVENTIONS ----
%  - \kpage{N} (margin [N]) marks where Koch's page N begins, at the same
%    joint as in the Danish working text.  Koch's p. 92 begins inside his
%    summary note, so here [92] follows that note's [ed. note].
%  - The same paragraphs, heads, rules, \emph spans (Koch's italics =
%    Brøchner's underlining), \ins spans, footnotes (same numbers, same
%    anchors) and \ed references as the Danish.  Lecture heads: "1ste
%    Forelæsning" = "Lecture 1", etc.
%  - »...« -> ``...''.  Latin, German and French stay as printed (roman
%    unless underlined); Greek as Koch prints it (τελοσ with its medial
%    sigma).  dvs. -> i.e.; jfr. / cfr. -> cf.; s. -> p.
%  - Koch's misprints in Brøchner's text are rendered by sense and
%    commented at the site; misprints in Koch's own references (e.g.
%    "SKS 12, s. 28f.", p. 108; "74,4", p. 164) are kept as printed.  Where
%    the Danish opens a quotation mark it never closes (pp. 105, 168) the
%    English supplies it, so the English marks balance (162 / 162).
%
%  ---- HOW THE TRANSLATION WAS CHECKED (2026-09-23) ----
%  (1) Twelve translators, one per batch of <= 10 Koch pages (72-81, ...,
%  172-181, 182-188), each re-reading the English against the Danish
%  sentence by sentence for omissions, negations, modals, quantifiers and
%  terminology before returning.
%  (2) Mechanical audit (tr_audit.py): per Koch page, the same \kpage
%  joints, \emph, \ed, \ednote, \ins, heads, rules, flushright blocks,
%  footnote numbers and paragraph breaks as koch-text.tex, and an
%  English/Danish word ratio within 0.95--1.60: 117 pages, 0 flagged
%  (ratios 1.03--1.29; 38,720 Danish / 44,628 English words).  Paragraph
%  breaks at all 117 page joints identical to the Danish; all 11 batch
%  seams read across (the text runs on mid-sentence at each, except p. 92,
%  which begins a paragraph in both).
%  (3) Independent review: twelve readers who had not translated the pages
%  read every sentence of the Danish against the English and returned
%  corrections as exact old -> new pairs.  No omission found.  73 proposed;
%  each checked against the Danish before applying: 60 applied, 13 rejected
%  (all would have translated "Bladartikler", a title cited by page).
%  Applied, besides terminology harmonization (39): p. 80 a clause that did
%  not parse; pp. 90-91 "som dens Princip" returned to p. 91; p. 108 italics
%  narrowed to "om Christendommens"; p. 123 "Gud" = God (not "the god"),
%  "først ... naar" = only when; p. 138 "vil stræbe" = wills to strive,
%  "Forestilling" = idea; p. 167 "uendelige lidenskabelige" kept as two
%  adjectives, as printed; p. 171 italics on "Bevidstheden", not on
%  "vorde"; p. 177 "et qualitativt Høiere" = something qualitatively
%  higher; p. 186 the plural "faae" construed; Koch's references restored
%  as printed at pp. 108, 157, 162, 164, 173.  Checked in the main pass:
%  p. 181 "altsaa" (dropped) restored; p. 168 closing quotation mark
%  supplied.
%  (4) Compile test (lmodern for libertinus): 101 pp., 0 errors, 0 missing
%  characters.
%
% ============================================================
\documentclass[12pt,a4paper,openany]{book}
\usepackage[utf8]{inputenc}
\usepackage[T1]{fontenc}
\usepackage{libertinus}
\usepackage{libertinust1math}
\usepackage{textalpha}   % Greek copied from Koch's text: τελοσ, μετρον, πνευματικῶς ...
\usepackage{graphicx}
\usepackage[english]{babel}
\usepackage[protrusion=true,expansion=true]{microtype}
\usepackage{setspace}
\usepackage[margin=1.2in]{geometry}
\usepackage{enumitem}
\usepackage{fancyhdr}
\setlength{\headheight}{28pt}
\addtolength{\topmargin}{-13.5pt}

\usepackage{hyperref}
\hypersetup{
  pdftitle={Lectures on Søren Kierkegaard (1858--59)},
  pdfauthor={Hans Brøchner},
  hidelinks
}

% ---- Koch's page numbers ----
% \kpage{N} marks where page N of Koch's edition begins, at the same joint as
% in the Danish working text (koch-text.tex, not published).
\usepackage{marginnote}
\newcommand{\kpage}[1]{\marginnote{\footnotesize\textit{[#1]}}}

% ---- Koch's editorial matter ----
% \ed{...}: a reference Koch adds in braces -- the manuscript's gathering.page
%   (1.2, 67.3), an SKS location, the doubtful-reading mark ? -- printed in
%   brackets and marked "ed.".
% \ednote{}: the place of one of Koch's editorial notes (description of the
%   manuscript, his summaries of unprinted lectures); not translated.
% \ins{...}: an addition Brøchner wrote in the right-hand column, which Koch
%   prints in the text between slashes.
\newcommand{\ed}[1]{{\small[\textit{ed.}\ #1]}}
\newcommand{\ednote}[1]{{\small[\textit{ed. note}]}}
\newcommand{\ins}[1]{\,/\,#1\,/\,}

\newcommand{\parttitle}[1]{\begin{center}\Large #1\end{center}\par}
\newcommand{\coursetitle}[1]{\par\medskip\noindent #1\par\medskip}
\newcommand{\lecture}[1]{\par\bigskip\noindent{\large #1}\par\nobreak\medskip}
\newcommand{\sepline}{\par\begin{center}\rule{2cm}{0.4pt}\end{center}\par}

\pagestyle{fancy}
\fancyhf{}
\fancyhead[LE,RO]{\thepage}
\fancyhead[RE]{\textit{H. Brøchner: Lectures on Søren Kierkegaard}}
\fancyhead[LO]{\textit{University of Copenhagen, 1858--59}}

\setstretch{1.3}

\title{\textbf{Lectures on Søren Kierkegaard}\\[8pt]
  \large [Forelæsninger om Søren Kierkegaard]\\[4pt]
  \normalsize University of Copenhagen, 1858--59}
\author{Hans Brøchner}
\date{From the text edited by Carl Henrik Koch, \textit{Hans Brøchners bidrag
  til Kierkegaard-receptionen i Danmark} (Copenhagen: Royal Danish Academy,
  2016), pp.~72--188\\[10pt]
  \small Translated from the Danish}

\begin{document}
\maketitle
\thispagestyle{empty}
\clearpage

\parttitle{\kpage{72}The Lectures}

\coursetitle{\ed{1.1}
Lectures on the Relation of Modern Philosophy to Christianity.
\textbf{(Autumn Semester 1858.)}}

\sepline

% Koch: only the gist of lecture 1 is recorded; the later lectures are written out in full
\ednote{}

\lecture{Lecture 1}

As the purpose of the lecture course it was stated: to arrive, through the
historical-critical investigation, at a \emph{present-day} problem: the relation between
faith and knowledge, religion and philosophy, as it is posed in our
day in the question of the possibility of a \emph{science of faith}.
But this problem, which coincides with the problem of the relation of the religious,
and specifically of the \emph{Christian}-religious, to reason and
thinking, leads on to a more comprehensive problem: the relation of the Christian
to those spheres of spiritual life which \emph{simultaneously with}
Christianity have developed into distinctive forms, which
as a rule are designated by the predicate ``\emph{Christian}'', (thus: Christian
\emph{state}, Christian \emph{ethics}, Christian \emph{art} etc.)

In other words, the problem to which we are led is to be
posed thus: is Christianity a power \ed{1.2} embracing and pervading the whole of
spiritual life, which takes up into itself all the forms of
spiritual life, or is Christianity something particular,
which stands only \emph{alongside} other forms of spiritual life without
being able to enter into an organic unity with them; --- and in that case: is
the task this: to make a choice between Christianity and
what stands outside Christianity; --- or: is it possible,
with the consciousness of the discrepancy and the heterogeneity, to
assign each of the two a separate sphere in the human spirit?
\kpage{73}The answering of this question leads, if not to a
reconciliation of faith and knowledge, yet to a clarification of the actual relation
between them, to an annihilation of illusions that prevent
a firm and self-conscious choice.

I emphasized, as the reason why I had chosen the historical-critical
form for the lecture course instead of contenting myself with
conceptual developments, \emph{partly} this, that the determination of the Christian had in
our day become so uncertain through its being mixed up with the
humane that it was necessary to go back to the \ed{1.3} times when
the factors that are now unclearly mixed appeared in their original
sharpness, and from there to follow the struggle between them in order to become
conscious of the distinctive features of both; --- \emph{partly} this, that I wished to seek
the objective testimony of history for what I set forth, and
to let the conceptual development come into being successively by way of the historical
presentation. ---

After these introductory remarks I went on to set out
\emph{the plan for the lecture course} (following the enclosed outline, which was carried out
more precisely orally \ed{see below}). Anticipating the result of a
critique of our age's philosophizing attempts at restoration, I showed
how their failure must lead to the consciousness that the reconciliation
between Christianity and philosophy had at least \emph{not
yet} succeeded, and through this consciousness to the setting up of
a formula that was to secure to each a separate domain, but in such a way
that the religious stood highest. (\emph{Kierkegaard}). With the critique of this
formula, then, the historical task would be accomplished, and a result
could then be drawn which, if not mediating, would at any rate be clarifying
and a safeguard against self-deception.

\sepline

% Koch: describes the enclosed outline (33 lectures autumn 1858, 30 spring 1859)
\ednote{}

\ed{1} Lectures on: \emph{the relation of modern philosophy to the Christian, and the attempts
to give a ``Christian philosophy''}

\sepline

\kpage{74}\emph{Introduction}.

1. The general essence of philosophy and of religion.

2. Philosophy and Christianity.

3. The various conceptions of the relation of the Christian to
concrete natural and moral actuality.

4. The relation between Christianity and philosophy as it appeared historically
in the period of patristics and scholasticism. ---
The ``Christian philosophy'' of that time.

\sepline

\emph{The Philosophy of the Modern Age}.

I. \emph{The transitional period}:

1. Polemic against scholasticism.

2. Attacks on Christianity.

3. Reaction. Catholic philosophy.

4. Christian mysticism and theosophy.

II. \emph{The period from Bacon and Descartes to our time}.

1. General characterization of the period. Its designation
as the period of the dissolution of the Middle Ages. --- The double
series of systems. The unlikeness and the likeness in their relation
to the principle of Christianity.

2. The \emph{realistic} line of development. --- Stages in its relation to
the Christian: defensive accommodation --- indifferentism\ed{2},
--- naturalism, --- conscious hypocrisy, --- open attack
(materialism --- atheism.)

3. \emph{the idealistic line of development}. --- The possibility of a double relation
to Christianity, given by the idealistic principle.

a) \emph{The relation of opposition} --- α) conscious opposition to historically
formed Christianity: Spinoza, the philosophy of identity,
Schopenhauer. --- β) Opposition to
Christianity, concealed under a reinterpretation of it:
the philosophy of the Enlightenment. --- Kant, --- Fichte; --- Hegel.

b) \emph{The relation of agreement}: α) formal agreement: Descartes,
β) limitation of the philosophical by the
theological interest: Geulincx, Malebranche (NB.
\kpage{75}\emph{Pascal}); --- γ) attempts at mediation: Leibnitz; --- δ) attempts at restoration:
the philosophy of faith, Schleiermacher, Steffens,
Fr. Baader, Schelling's philosophy of revelation,
post-Hegelian philosophy. --- Critique of these attempts.

4. General result. Consciousness of the discrepancy.
\emph{Feuerbach}. Kierkegaard's formula for the separation. Critique of it.
--- The task of philosophy.

\lecture{\ed{1.4} Lecture 2. Introduction}

% Koch's note 4: the parenthesis "(Væsen)" and the word "Forhold" are added in pencil
1. \emph{The general (essence) relation of philosophy and religion}.\footnote[4]{\ednote{}}

Since the more determinate concept of religion and of philosophy is to
emerge from the analysis of the historical phenomena, namely of the
phenomenon that is specifically designated ``\emph{Christian philosophy}'' (patristic
and scholastic philosophy), what remains for the introduction is only
to indicate what is generally characteristic of both, \emph{in connection with
the common idea}.
% Forestilling: "idea" where it means the ordinary notion (here and p. 76); "representation" (Hegel's Vorstellung) where it is set against thought and concept (from p. 76 on)

As the most abstract determination of religion it may be taken that
it is \emph{the human being's relation to the divine}. But this leads at once to
the necessity of a closer determination: is it a relation \emph{in
thought} or is it an \emph{existential} relation? At this point the relation
between thought and existence is still left undecided, and since the question
has been raised with particular regard to the determination of the mutual relation of religion
and philosophy, while the domain of philosophy
is indisputably that of thought, and moreover everything \ed{1.5} that is to exist
for the human being, for human consciousness, must be brought
into the forms of consciousness, hence into the forms of thought in the widest sense,
the answer is provisionally restricted to this: that the religious relation
is a relation of consciousness, without its being ruled out thereby that it can also
be something other and more.
% Koch prints "uden af det afskjæres" (for "at"); rendered as the sense

Religion is thus conceived provisionally as \emph{cognition} of God, and of
the definition of the old Protestant dogmatics: religio est modus
\emph{cognoscendi} colendique Deum, only the first half is provisionally taken,
so that religion and philosophy are accordingly set over against each other
as \emph{forms of the cognition of the divine}. The justification for such a
\kpage{76}provisional restriction lies not only in what was indicated above,
but also in the basic religious category: \emph{revelation}, which
expresses that religion, as a divinely communicated \emph{truth},
enters into relation with the receptive human \emph{consciousness}, and through
this into relation with human \emph{knowledge}.

When religion and philosophy thus meet in being a relation
to the divine through consciousness, and when it is assumed
to be possible that the divine is accessible to philosophy,
which, if the opposite\ed{1.6} were assumed, would cease to exist,
since philosophical knowledge would cease to deserve that name
if the highest and concluding, that which in the last instance grounds
all knowledge, were excluded from it, then the next task is to
determine what separates them within that likeness. We can find a
starting point for the solution in the common idea of philosophy.
To the name of philosophy consciousness most immediately attaches
the vaguer idea of a \emph{thinking consideration of the thing}.
But when what lies in this idea is developed more determinately,
we are carried a long way forward. In considering things thinkingly,
we apply the laws and forms of thought to them, and since these
prove to be universal and universally valid, thinking does not stop
at the partial, but seeks the all-embracing connection,
seeks to comprehend the wholeness of things, existence as totality. And
as it advances in this comprehending and finds the forms of thought
to be the essentialities of things, it finds \emph{itself} again in things, and in the
comprehended whole its own self-unfolding. Philosophy thus becomes
\emph{the systematic unfolding of logical thought as the ideal image of existence};
or in other words, it is \emph{the reason that through objectivity}\ed{1.7} \emph{has been
confirmed}. But in this determination lies a consequence that becomes
important for the relation of philosophy to religion. Philosophy
is thought unfolded \emph{by the very energy of thought}, a \emph{result}, then,
not something \emph{immediately} given. Immediate consciousness is consciousness held fast
in \emph{sense impressions} or in \emph{feeling}, which has not yet
freely acquired itself in the form of universality. The consciousness striving from the unfreedom
of immediacy toward the freedom of thought
is \emph{intuiting, representing} consciousness; in it thinking begins
to stir, but thought has not yet found its own form. And
since this form constitutes philosophy, philosophy becomes something
\kpage{77}\emph{acquired}, and as such the property only of those who themselves produce
it, --- the property of \emph{individuals}.

In contrast to this stands \emph{religion}, as addressing itself to \emph{all}
human beings at whatever level of consciousness: the \emph{form} of religion
--- of the content there is as yet no question at this point --- must therefore
be one that is accessible to all; it cannot be the acquired
--- and self-acquired --- form of thought, --- it must be that of feeling
and of \emph{representation}, the latter especially insofar as the religious content
becomes the \emph{object} of consciousness. Provisionally, then, \ed{1.8} the difference
between philosophy and religion --- within the domain of cognition
--- is determined by the difference between \emph{thought} and \emph{representation}.

Through this difference in form, however, the problem of the relation is carried
a step further. Since the form of thought is that \emph{of universality, of the concept},
% Koch prints "Abolute" (for "Absolute"); rendered as the sense
the absolute is conceived by philosophy as \emph{idea}; the unconditioned
content is conceived in a form that is adequate to it;
the inner movement of the concept becomes the expression of the movement in itself
of the absolutely concrete content. The life-process of the absolute idea is seen
sub specie æterni as an \emph{eternal} process, in which the basic forms of existence
become moments that have their ideal subsistence in the whole.
And the eternal ideal process, which is conceived in the thought-form
of universality, is eo ipso conceived in the form of \emph{necessity}, but a necessity
which, as one with the self-unfolding of the absolute, is eternally transfigured
into \emph{freedom}.

As distinct from this, the absolute, when it is conceived in the form of \emph{representation},
is determined --- provisionally at least \emph{formally} --- by the
peculiarity of that form. As the psychological middle term between
sensation and \ed{2.1} concept, representation has in it a moment of
both: it combines the receptive dependence of sensation with
the freedom of the concept, its singularity with the universality of the concept.
What I represent to myself is detached from the immediately given external,
transferred from outer space into the inner space of consciousness; it has begun
to be separated from the contingencies of the external particular, but as to its content
it is still determined by them. The consequence of this is that what is
taken up into representation is affected by the form of sensuousness, and since
the absolute is to be conceived through it, its universal essence must
be intuited in a \emph{sensuous image}, its eternal, ideal movement within itself,
through which it has its content, as a \emph{process in time}, whose individual
\kpage{78}moments fall apart \emph{before} and \emph{after} one another: the forms of the \emph{image}, of \emph{myth}
and of \emph{history} become inseparable from the absolute.

Here the question now arises: \emph{does this difference in form entail a similar
difference in content}? is it not the same absolute content
that passes through both forms, untouched by their difference?
This question of the mutual relation of form and content,
which, as a question of the relation of the two categories\ed{2.2} to each other,
must find its answer in logic, and in its more concrete form,
with particular regard to religion and philosophy, can find
it only through the development of the historically given, can nevertheless already
be answered here in general, by way of indication. It is already
evident here that a spiritual content can be clothed in a form
inadequate to the content without being transformed for consciousness only
when it \emph{has been} present for consciousness in the adequate form.
The form then becomes a clothing: \emph{symbol}, \emph{allegory}. But in the development
of humanity the form of religion always precedes
philosophy: the absolute is present as an object of \emph{faith} before it is
present as an object of \emph{knowledge}. When it takes on the representational form of religion,
it has therefore not been present in any other form
for consciousness; the form is therefore not merely external, inessential form,
clothing; but it is for consciousness the \emph{only} form for the content,
and so it counts as the true and \emph{adequate} form. The content
thus cannot remain untouched by the peculiarity of the form;
it is for consciousness only as \emph{formed} content, and a separation
from the form becomes impossible at the original \ed{2.3} religious standpoint.
A further separation takes place only through an awakening \emph{reflection},
and at the religious standpoint such a reflection comes forward only
in relation to \emph{other}, subordinate religious standpoints. In relation
to its own content the religious consciousness is by the nature of the case
uncritical: it finds itself again in the content \emph{so formed}.

Now, when this relation between content and form points to
the essential influence of the latter on the former, then by the same
token a \emph{real difference of content} between philosophy and religion
has been made at least probable. The object of both is indeed ``the absolute'',
but the absolute as such is as yet only an abstract
concept, which first receives its closer development through the form
by which it is determined.

% "Navnelighed" read as sameness of name (the object is called "the absolute" on both sides)
\kpage{79}In the difference of content despite the object's sameness of name lies
a \emph{double consequence}. \emph{Partly}, it points to the essentially
different standpoint of philosophy and religion; \emph{partly}, there lies
in it the possibility of a conflict. Since what is named in the same way,
the absolute, what is posited as the highest, is determined differently
in the two domains, the difference cannot confine itself merely
to the difference\footnote[5]{NB. 1.5. A contradiction with that must be avoided.}
between the two forms of \emph{cognition} of it, \ed{2.4} but
their difference, as a difference in the forms of cognition of the
highest, seems bound to stand in connection with a difference in \emph{the whole standpoint}
(view of life) and in the driving \emph{interest}. On this
point too the difference can only be indicated provisionally, insofar as it lies in
what has already been developed. Since the form of philosophy is thought, but
thought in its universality, its universal validity and objectivity,
lets the individual vanish, in philosophy the accent falls
on \emph{the matter}, the objective. The thinking individual immerses itself in its
object, gives itself up to it, and even if, in dissolving
the essence of the object into thought, it finds itself again in the object, it is
not itself as this single individual, but the objective thought
that at once constitutes the essence of the object and of the individual. In
philosophy I thus relate myself \emph{disinterestedly}, \emph{theoretically}; I relate
myself \emph{freely} to the object, letting it subsist in its
freedom. The relation must turn out otherwise where \emph{representation} is the form
of cognition. Since representation is burdened with the form of the sensuous,
it is burdened with \emph{unfreedom}\ins{contingency, subjectivity},\footnote[6]{this must be developed more precisely.}
but to theoretical unfreedom there must correspond a \emph{practical} one. In my
relation to the object, theoretical freedom is the condition for
my being free myself and letting the object subsist freely in its objectivity.
When I am theoretically unfree, \ed{2.5} I am in practical respects both
unfree myself and unable to treat the object according to its
objective peculiarity. To the theoretical standpoint which,
while striving toward the freedom of thought, still has the unfreedom of sensation
within it, there corresponds a practical one in which the will is determined by \emph{desire} and
\emph{passion}. What impels here becomes \emph{the individual's desire}, its \emph{need} and
wishes; and this desire influences the conception of the object.
\kpage{80}Not \emph{the nature of the matter} but \emph{the heart's need} becomes the decisive thing here,\footnote[7]{(NB. The miracle, prayer).}
and while the inadequate theoretical form conditioned an unfree practical
relation, this in turn conditions the incongruence of the theoretical conception
with the nature of the matter.

As the \emph{second} consequence of the difference of content,
\emph{the possibility of a conflict} between the two domains was indicated. This conflict
we shall follow in its various forms through the whole lecture course;
at this point only the essential possibilities are to be touched on
briefly.

When the difference of content becomes known, either religion can hold it fast
in its interest, or philosophy in its own. \emph{Religion} can, by
declaring its content to be \emph{divine revelation}, set it in
opposition to human knowledge. The opposition then becomes one
between the infinite and the finite, between divine wisdom
\ed{2.6} and human cleverness, between the Holy Spirit and the human spirit
darkened and confused by sin. Where the discrepancy
appears, religion uses its language of authority; it
denies the ability and the right of finite human thought to grasp
the infinite, it requires faith in what darkened reason
cannot comprehend, and at most lets reason keep some latitude
in the domain of the worldly and within the bounds of the formal.

Against this requirement of supremacy on the part of religion,
\emph{thinking} that knows of the discrepancy can now come forward with a
similar requirement: reducing the religious doctrines to
the forms of representation and of imagination, while as the form of truth
there steps forward \emph{incorruptible} reason, the organ of divine
truth, it vindicates for itself the judgment about the true and the possession
of the true, and that not only in the \emph{theoretical} domain, but
also, insofar as it maintains the impossibility of ethical action without
a rational cognition as its foundation, --- in the \emph{practical}. It requires
not only that it rule in the theoretical, but that true theoretical
cognition shall pervade the whole of human life and
become the principle of all the forms of moral life. \ed{2.7} Against religion's
separation of the domain of the divine and of the human, of the
\kpage{81}infinite and of the finite, it sets the dialectical proof
of the inseparability of the two domains.
% "og som Sandhedens Form træder den uforvanskelige Fornuft": construed as "while incorruptible reason steps in as the form of truth"

Through the struggle between these two opposites one may arrive
either at an acknowledgment of the hegemony of the one, or at the attempt
to delimit the domains of both against each other, so that the relation
becomes an apparently neutral one, yet with a supremacy
for the religious. These various possibilities will be developed
in what follows.

After this provisional development of the relation between philosophy
and religion in general, the next task is to
determine the relation between philosophy and the particular religion
with which these lectures are to be concerned: \emph{the Christian}.

\sepline

\ed{2.8} 2. \emph{Philosophy and the Christian religion}.

It might seem that what was developed in the preceding\footnote[8]{The development must be formulated more precisely.} section about
the relation of philosophy to religion must \emph{immediately} be able to find
its application to \emph{Christianity}, since this after all falls together
with the other religions under the \emph{species} concept religion. But here
there appears the possibility, which has been asserted by \emph{Hegel},
that Christianity, although it is united with the other religions
in a common genus concept, can be essentially different from them, and
thus come to stand in a different relation to reason and philosophy
than they do. The ``historical'' religions preceding Christianity
would then express the essence of religion only
\emph{partially}, would develop only one-sidedly particular moments of its concept, and
precisely in this would lie the fact that they were in disagreement with philosophy,
did not fulfill its requirements. Christianity, by contrast,
as the absolute ``revealed'' (manifest) religion, would realize
the concept of religion, be the adequate expression of it, and
eo ipso agree in content with philosophy.\footnote[9]{since religion has its root in the thinking nature of the human being.}
% "aabenbarede" (aabenbare): Hegel's geoffenbarte / offenbare Religion; the play is lost in English

In what has been developed in the foregoing about the form of religion,
and about the relation of this form to the content, this possibility has \ed{3.1}
\kpage{82}already been judged.
% Forestilling = representation (Hegel's Vorstellung) in this passage
The form of representation is posited as the characteristic one,
not for a single imperfect form of religion, but for
\emph{religion itself}, and it was indicated how the form must affect
the content. This implies that, even if a single historically emergent
form of religion exhausted the essence of religion, it would nonetheless
fall under that determination which holds precisely of religion \emph{as}
religion.

And Hegel will not be able to escape this conclusion either,
since he precisely holds fast to that distinction between representation and thought,
and since, with regard to the relation between form and content,
he holds fast in his Logic to their \emph{essential} relation, and where in
his philosophy of religion he lets the content be indifferent to
changes of form, he does this only by taking the content in a similar
indeterminacy as when he lets the name of \emph{religion} cover
the \emph{whole} sphere of absolute spirit, and thus also \emph{philosophy}.

Were one to take the identity of philosophy with Christianity seriously,
one would not escape \emph{consequences} either that the Christian
consciousness would not acknowledge. For were
that identity carried through, then the \emph{historical}, which is the
foundation of Christianity and the point of departure of the Christian consciousness, \ed{3.2} would become
a merely vanishing garb for an \emph{eternal} process; it would become
image, \emph{myth}, not self-valid reality. And at the same time the \emph{transcendence}
of the divine as \emph{personal} would vanish: nature and
finite spirit would become moments in the very life-process of the absolute;
moments that were eternally posited and eternally sublated as ideal.
% TR: ophæve = sublate (Hegelian sense)

The premises of philosophy itself and the religious consciousness thus
lead to the same result: the rejection of an identification of
Christianity with philosophy, and the applicability of what holds
for religion to Christianity as well.

But this has here been reached essentially by way of \emph{reasoning}; it must
be corroborated by investigating \emph{historically} \emph{what Christianity} is; which,
however, can take place in the introduction only in general outline,
and therefore in such a way that the result stands as a hypothesis
until the investigation of the concrete historical relations corroborates
it. ---

\sepline

\lecture{\kpage{83}\ed{3.3} Lecture 3}

% TR: Virkelighed = actuality
\emph{The Christian view of life and the relation of the Christian to concrete natural}
\emph{and spiritual actuality}. ---

If, in answering the question \emph{what Christianity is}?
we were to keep to the conception of our time and to the phenomena that now
appear in the domain of religious life, then the latter would present
a difficulty for analysis on account of the unclarity of their
composition, and the time's conception of the Christian would
prove so little certain that \emph{two fundamentally different views} stand opposed to
each other.

% TR: Uensartethed / Ensartethed = heterogeneity / homogeneity
On the \emph{one} side Christianity is held fast in its heterogeneity
with everything worldly and with what is accessible to reason. Its
center, \emph{the God-man}, is the absolutely incomprehensible, the \emph{paradoxical}, to
whose acknowledgment in faith it is not the inferences of reason but the despair
of the consciousness of sin that brings one. One does not arrive there \emph{by} thought,
but \emph{in spite of} thought and through a break with it: faith grasps
the incomprehensible, and in holding it fast with the passion of inwardness,
it guards precisely its incomprehensibility against every
attempt of thought \ed{3.4} to appropriate it. And just as Christian
\emph{faith} is posited only as having issued from the break with thought,
so the Christian \emph{life}, which springs forth from this faith, is posited as
having issued from a break with worldliness. Christianity is the religion of \emph{suffering};
its spirit is ``\emph{to die away}'' from the world; only through this
death does it show the way to life. And this \emph{negative} relation to the world
is not confined to a single time, as e.g. to Christianity's
first historical appearance; as little as the absolute paradox of faith
ever becomes more comprehensible with time, so little does
the life that rests in it ever become homogeneous with the life of the world. Christianity,
which, in demanding death, the death of selfhood and the renunciation
of everything worldly outside and within the individual, arouses the resistance of selfishness
and worldliness, must constantly stand equally
hostile toward ``\emph{the world}''; it must at all times be reviled,
scorned, hated, persecuted. And where history shows a peaceful relation,
that is precisely a sign that Christianity has lost its
distinctive character, that it is corrupted by a mixture
with the heterogeneous, that it has lowered the ideal requirement, denied\kpage{84}
itself. In the ``\emph{Christendom}'' of our time this
conception therefore does not recognize \ed{3.5} \emph{Christianity}; it sees in it only a deception
continued through many centuries and carefully concealed; in the
\emph{Church} protected by the state only a distorted image of Christianity,
whose preachers live \emph{off} Christianity instead of living \emph{for} it.
It doubtingly raises the question: \emph{whether Christianity is still to be found}
\emph{on earth}, and in going through the history of the Church, with its ideal before its eyes,
right back to its earliest period, it carries this doubt with
it right back to that period, and asks doubtingly: have there really \emph{ever} been
Christians? has Christianity's ideal requirement \emph{ever} been fully
met? And this question is consistent; for can the religion
whose principle is \emph{not} to live in the world, to be dead to the world,
really find its perfect expression \emph{in} the world?

% TR: det Humane = the humane (kept distinct from det Menneskelige = the human):
% the sphere of the universally human, Humanität, set against the Christian -- not "kind"
In opposition to this conception of Christianity stands
\emph{another}, which extends the name of the Christian to all the forms
of spiritual life that have developed \emph{simultaneously with} and \emph{under the influence}
\emph{of} Christianity. It does not see Christianity in opposition
to our time, but as \emph{its animating principle}; it speaks of \emph{a}
\emph{Christian world}, a \emph{Christian state}, a \emph{Christian ethics}, a \emph{Christian art}, a
\emph{Christian science}, in a \ed{3.6} word, it applies the name of Christianity
to all forms of the \emph{humane}. The sharp opposition
between the principle of Christianity and that of the world it restricts,
even if it constantly points to the \emph{apostolic} as what orients,
to Christianity's first period, in which the struggle against
the forms of the old world had to lead to a \emph{one-sided} opposition.
From the moment when Christianity had become victorious,
its task was not to be to separate itself from the world, to die away from it,
but to \emph{permeate it}; its power was to show itself precisely in this, that it
could
% Koch: "formaaede" written above "kunde" as an alternative
\ednote{}
% Koch prints "gjennnemtrænge" (three n's); rendered as the sense
% TR: potensere / Potensation = potentiate / potentiation
do this; and the possibility of permeating, as a principle of life, the
heterogeneous is sought in this, that Christianity was a \emph{new creation}, which
acts \emph{potentiatingly} back upon nature and upon the powers of spiritual life
and its laws. Not flight from the world but immersion in the world, not
struggle against it but peace with it and dominion over it then becomes
Christianity's goal, and in the \emph{Christian state} it has thus been
reached, in such a way that precisely the state's acknowledgment of the Christian, its
favoring of it in the way of worldly goods, is a proof that
\kpage{85}the Christianity thus acknowledged is the true one. And this
acknowledgment is also granted by \emph{reason} \ed{3.7} to Christianity. Even if
the unregenerate, debased reason does not comprehend the mysteries of Christianity,
yet the \emph{believing} reason, potentiated by the new principle of life,
is able to penetrate into their depths, and, even if
it does not attain \emph{logical} and \emph{mathematical} clarity, it can nevertheless
explore the \emph{connection} of the articles of faith and \emph{thereby} their\footnote[10]{this presentation had best be kept somewhat ironic-pathetic.}
\emph{ground}, and
develop a \emph{science of faith}. Philosophy comprehends in Christ the unity
of the divine and the human and acknowledges in this unity
its own doctrine of the relation of absolute spirit to finite spirit.
Christianity is thus reconciled with thought and with the world,
for whose concrete forms it becomes the new animating principle of life.

If we now raise the question: \emph{which of these conceptions is the}
\emph{true one}, then at this point it cannot be answered exhaustively, but
the answer can only be \emph{indicated} by an indirect proof, and this indication
can be made \emph{probable} by \emph{historical} considerations.

If one assumed without further question the \emph{latter} view to be the true one, so
that the question of what Christianity was were to be
answered by analyzing the forms of spiritual life as
\ed{3.8} they have \emph{now} taken shape, since it would then be found as their
principle; then through a new question a \emph{new}
\emph{requirement} would at once present itself. Is the unity between Christianity and the forms of spiritual
life such that the former can issue from the purely human
by a continuous development? --- \emph{or}: is it something new which,
in entering history in a primitive manner, does not issue
from what was historically earlier, even if the latter can \emph{typologically} foreshadow it, but
acts potentiatingly back, so that the unity we now find has been
brought about through such a potentiation? For one who will not
deny Christianity a supernatural origin, the
\emph{latter} answer becomes the one that must be preferred. But if it is preferred, then
the requirement cannot be dismissed: \emph{first}, to demonstrate the actual
inner unity of the specifically Christian with the human, and \emph{second}:
to demonstrate in particulars the concrete, the so-called ``\emph{potentiation}''
of the spiritual faculties and powers that Christianity is supposed to
have brought about. Now if it does not succeed -- and at least hitherto it has not
\kpage{86}succeeded -- in demonstrating in determinate forms such a\footnote[11]{to be carried out ironically.} potentiation,
when one compares the individual Christian state, art, science,
etc.\ with those independent of Christianity; or if it succeeds
only in demonstrating something specifically different from antiquity \ed{4.1} as
something that does not combine organically with what has issued from
the human spirit, then in the latter case we are led back to the
view that sets the Christian at variance with the world; in the former
we are led to see the Christian as a \emph{natural development} of the ground of the human
spirit. But with this latter we would either immediately
abandon the religious standpoint, or, if we still wished, by a
vague conception, to hold fast to it, and let Christianity
be something supernatural although \emph{now}, after having been revealed, it joins
together without a break with the purely human, then it would surely have to
be asked, \emph{partly}: whether a revelation would not be superfluous
when it communicates only what the human spirit
could develop\footnote[12]{partly whether such a revelation, revealing nothing, could be maintained as
actual.}
% fn 12: "deels ved en saadan ... Aabenbaring kunde fastholdes": "ved" does not construe; rendered as "whether"
% "om der da vil saa, med en Analyse": the Danish does not construe (Brøchner marks the passage "forvirret!"); rendered by the evident sense
of itself; --- (\emph{partly}, whether then, with an analysis
of what is now called \emph{Christian}, a determination of the
essence of Christianity could be found that secured for it a distinctiveness actually
explaining its history?)\footnote[13]{\ednote{} confused! \ednote{} partly whether the history of Christianity can then be explained
on the presupposition of such a joining together with something preceding, and
whether it can secure for Christianity an essential distinctiveness.}

The \emph{first} of these questions the religious consciousness would
naturally have to affirm at once; the \emph{second} \ins{last} would undoubtedly\footnote[14]{(The development here is somewhat confused).} have to be answered
in the negative. For a determination of Christianity
that proceeded from its \emph{homogeneity with the humane} would have to let the \emph{historical}
side of\footnote[15]{would first have to disregard the historical starting point and fix
the dividing line solely on the content of ideas.}
Christianity recede, and essentially bring out
the \emph{ideal} side. But \ed{4.2} in this way it would, even if it
could demonstrate a stronger and more determinate appearance of the
\kpage{87}\emph{universalistic} and \emph{spiritualistic} element than in the\footnote[16]{\ednote{} Christianity
would then have to be brought together idealistically with the ideal content in the systems that conclude
the philosophy of antiquity.}
philosophical systems
that conclude the philosophy of antiquity, not be able to find a
view of life that was \emph{in principle} essentially different from what has
arisen outside Christianity. Christianity\footnote[17]{(This point must be determined somewhat more precisely orally; it can find its proper
development only in the analysis of the phenomenon of Christian philosophy.)}
would then
enter the series of the historical forms of development of the human spirit;
its appearance would not be \emph{a break} in continuity,
not a \emph{new creation}; but as such it does, after all,
indisputably proclaim itself, while at the same time declaring what historically preceded it
to be untruth and sinfulness; whereby the possibility of a
continuous development out of it is excluded. \ins{The necessity of
Christianity's struggles with the world would become inexplicable.}
If that homogeneity were held fast, then all the Christian phenomena
in which the break with the world appears most clearly would also have to
be posited as one-sidednesses and exaggerations, and \emph{that} would be brought out
as the truly Christian in which the agreement with the purely
humane appears most decisively.

When these considerations make it, provisionally at least, \emph{probable}
that the view which sets Christianity \emph{at variance with the world}
is the true one, then this will also find its corroboration by \ed{4.3}
considering the \emph{historical circumstances} under which Christianity appeared,
and the \emph{phenomena} through which it, with its appearance,
most decisively separated itself from the nature-bound view of life.

If we wish briefly to consider the main forms of spiritual life
that formed the presupposition of Christianity, then it is
above all the development of the \emph{Greek} and the \emph{Jewish} spirit that must
be taken into account. The task of the \emph{Greek} spirit had been to realize
the idea of \emph{beauty}, the task of its thought to find the concept of the \emph{idea}, and,
proceeding from this, to seek the beautiful, harmonious unity of form and
material, spirit and matter. With the nature-life and view of nature of the Orient
as its presupposition, it had raised itself from material as principle to
form, the idea, had grasped this as the true and only being, and
\kpage{88}now strove to find the formula for its relation to non-being,
matter. But in the course of this search it had to become clear that the
Greek spirit, precisely because it rested in the immediacy of beauty,
could not carry through the sole being of the idea; it had to end in
an unreconciled \emph{dualism}, which led either to skepticism or to a failed
attempt, by raising itself above the foundation of the Greek spirit,
nature, to find the solution of the problem in a spiritualistic sphere.
But once the idea \ed{4.4} had been posited as the highest, as
the true being, the development could only go in the direction
that it, when the consciousness of its incompatibility with matter became
clear, is posited as the \emph{only} being, matter as non-being,
untrue, \emph{evil}. The task, therefore, that was set by the problem left unsolved
by the Greek spirit had to be this: to comprehend \emph{spirit as the}
\emph{only being}, and thus to develop an \emph{abstract-spiritualistic} view.
Herein lies a probability that this task was Christianity's.

If we consider the other of the peoples of antiquity whose spiritual development
immediately prepared Christianity: \emph{the Jewish}, then we see
a similar direction. The earlier Hebraism's secure resting in the
present and this-worldly had gradually been transformed, \emph{first} into
faith in a \emph{future} worldly messianic kingdom under a ruler of David's
line, \emph{then}, under foreign influence, into the development of the doctrine
of a spirit-realm that is the prototype and truth of the earthly, and of
the idea of a \emph{supernatural} Messiah and a \emph{resurrection} from the
dead, and \emph{finally} into an ascetic flight from actuality and into the hope
of spirit's perfect \emph{liberation from earthly matter}. Here too, then,
there is a striving toward the \emph{supernatural} as the only true, \ed{4.5} and
existence divides dualistically into a kingdom of the world and of matter, the
kingdom of the dark powers, and a longingly awaited heavenly \ed{?}
kingdom in a nature-less actuality. From a starting point other than the
Greek, the same goal had thus been approached here, and the same
task set for a new religion.

When \emph{Christianity} appeared historically where the Greek and the
Palestinian views of life met, and at a time when both were
developed in the same direction, it was to be expected that its view of life
would be the complete expression of what those strove toward.
And when one considers the points in Christianity that most decisively\kpage{89}
characterize it as something new, then it is precisely such as
point to an \emph{abstract spiritualism}, which can relate itself only \emph{negatively}
to the natural. In the person of the \emph{God-man} faith beholds the living
unity of God and the human being, but the human in Christ
is detached \ins{, by the \emph{miracle}, which marks \emph{all} the main points of his
life,} from the natural conditions of human life; the unity becomes
a unity between God and \emph{the nature-less human}, and the \emph{believer's}
unity with God through Christ is conditioned by \emph{having died away} from
the world, the natural; only as \emph{nature-less spirit} does the human being have unity
with God. For \emph{the natural} is viewed as ``\emph{lying in evil}'', as that
which is corrupted \ed{4.6} by sin, as that to which spirit is to relate itself
only negatively. The natural (the natural faculties and
powers) is granted no significance for itself; it does not become a
material that is to be formed by the Christian; it becomes only something that
is to be cast away. The Christian is to \emph{flee} the relations of natural life,
which have no significance for him, but can only become defiling;
he is to become \emph{indifferent} to the moral relations resting on a natural
ground; for these cannot in their concretion be placed in
the positive relation to the abstractly spiritual (e.g. marriage,
\emph{political life, property}). The determination of the natural is: \emph{in itself to be a nothing};
this is expressed in the dogma of \emph{creation}, which rests on the view of the
absolutely mighty nature-less spirit as that which in truth is, and
on the view of nature as what has been called forth, by an incomprehensible act of the
almighty arbitrariness, to a lived momentary existence.
% "levet momentan Existents": "levet" is doubtful (perhaps for "blot"); rendered literally
Over against the dogma of creation the concept of natural \emph{laws}
and natural \emph{forces} vanishes, and in general that of a ``\emph{nature}'' as distinct from a
creatum. Nature has no being in itself; there can therefore be no
talk of laws, which presuppose something substantially abiding; at every moment
the divine power must call nature forth out of nothing
through \emph{preservation}, and at every moment the divine
arbitrary \emph{omnipotence}, which manifests itself through\ed{4.7} the \emph{miracle}, is free
to let nature vanish; for it is by its essence something \emph{superfluous}, and
not merely that, but something that the Christian longingly wishes
away, from whose limitation, which is felt only as oppressive
bonds, he wishes himself free, a wish which the hope of \emph{immortality}
\emph{expresses}; for in the Christian heaven every natural limitation is
removed, and the idea of immortality expresses precisely\kpage{90}
what for this life stands as an ideal, but an unattainable one.
What the Christian wishes for here, but despairs of being able to attain,
he transfers to a life beyond, where what here hinders and
weighs down shall have vanished. --- And in the same\footnote[18]{(NB. ``spiritual body''.)} way as the
Christian relates himself negatively to nature, he consistently relates himself
negatively to \emph{history}. It is not accidental that the first Christians
awaited the \emph{Parousia} as imminent: history was, after all, \emph{concluded}
with Christ's coming; its continuation was, after all, necessary only as long
as was needed for the peoples to gain knowledge of
it. The decision was therefore thought so near that the generation then living
would not yet have died out; it was still only a matter of the choice
between Christ's kingdom, which \emph{was not of this world}, and the kingdom of the world, between
life in the world and death from the world in baptism, through which
the one believing in Christ (Christ's resurrection) appropriated his
life, and now no longer lived in the world, \ed{4.8} but \emph{spiritually}. There was
no question of a transformation of the worldly; it is left standing
as something indifferent, which did not touch the inner; Christianity
was not to recreate the world, but to create a new life that was \emph{not}
of this world, and this new life was to be livable only through death
from the world and everything that belonged to the world that was soon to
cease. In his death Christ had died from the world; by his \emph{resurrection}
he had overcome death, vindicated his divine
nature and \emph{secured the victory of the kingdom of God} \ins{of which the pre-Christian consciousness
had to despair}; he had by it made it possible for
believers, by giving up their life in the flesh, in him who now lives
% Koch prints "πνευ ατικῶς" (the μ dropped); rendered as πνευματικῶς
πνευματικῶς, to live a new spiritual life by which they, in the midst of the world,
are raised above the world. --- And everything Christian is posited as \ins{\emph{incomprehensible}
to the world, as ``\emph{foolishness}'' to it; only to faith are the mysteries
opened, which were inaccessible to \emph{carnal reason}. The wisdom of the world
and the wisdom of Christ were absolute opposites.}

We are thus led, through these prominent tenets of
Christian doctrine, to side with the conception that characterizes
Christianity by the \emph{sharp opposition} between the divine
and the \emph{concretely} human, between spirit and nature;
which designates an \emph{abstract spiritualism}, which relates itself only negatively
\kpage{91}to the concrete forms of natural and spiritual life, as its principle.
But an \emph{objection} could at once be made here, which must at least
provisionally be met. One could point out that Christianity
has, after all, \emph{acknowledged} the forms of human life, that it
has \emph{entered into connections} with the concretely human\ed{5.1}, and not
merely \emph{made use of} the faculties and powers of the human spirit, but also
sought to \emph{appropriate} them. To this objection, which in what follows will be met
more fully on many points, a provisional answer can already
be given here. One can answer with the question:
how would an abstract-spiritualistic view of life have to position itself in
relation to the concrete forms of human life and to the powers of the
human spirit, when the hope of an imminent
final conclusion of history was disappointed, and the task now became:
with a world-denying principle to live \emph{in} the world, to seek an existence
\emph{in the midst of that} which it posited as the untrue? If one pursues this
relation, it will easily be seen that the phenomena appealed to become
explicable when one proceeds from the conception to which I have given preference,
whereas the phenomena that most decisively characterize
the classical period of Christianity stand as \emph{incomprehensible} when one
proceeds from the homogeneity of the Christian with the humane. ---
% "nærforstaaende" (cf. "nærforestaaende", p. 90) rendered "imminent"

\sepline

% Koch: summary of the unprinted lectures 4-33 (runs on to p. 92)
\ednote{}

\kpage{92}

\coursetitle{\ed{44.2} (\emph{Spring Semester 1859})}

\lecture{Lecture 1}

Survey of what was gone through in the previous semester, --- summing up
of the results. --- Preliminary orienting survey of what follows.
--- (The development is given according to the points indicated
on the enclosed leaf.)

% Koch: the enclosed "leaf" described (a folded sheet written on three sides)
\ednote{}

% TR: Momenter = points (the heads of an outline)
\ed{1} \emph{Introductory lecture}. (8/2 59.)
(Points.)
1.) \emph{Statement of the relation between what was developed in the previous semester and what
remains}. --- The latter comprises nothing but phenomena that reach directly
into the most recent time; it comprises at the same time representatives of
\kpage{93}the most essential still-existing conceptions of the religious, and
therefore forms a smaller, fairly self-contained whole of its own.

\sepline

% TR: Opfattelse = conception
2.) \emph{Retrospect on what has gone before; summing up of its preliminary result}. ---
The result has hitherto been essentially \emph{negative}; it has been shown that there
has \emph{hitherto not} been an agreement between the philosophical
systems and Christianity. What was put forward \emph{hypothetically} in the Introduction
concerning the relation of religion, and specifically of Christianity,
to philosophy as a relation of divergence has found its
confirmation negatively; but a \emph{final} result, a \emph{positive}, decisive
and exhaustive determination of the relation, has been postponed to what
follows. Yet there already exists a \emph{material} from which a
\emph{preliminary} result can be extracted. (Here reference is made especially to the account
of the character of Scholasticism, of the direction of the transitional period, of
the unfolding of the realistic and the idealistic series. What has been given in the
earlier lectures is concentrated: in the case of the realistic
series the significance of the \emph{nature-moment} is brought out, while the \emph{one-sidedness} in its
conception and the ground of the one-sidedness are at the same time pointed out; --- in the case of the
idealistic, there are brought out: the autonomy of \emph{self-consciousness}, the influence of the \emph{concept of nature},
the conception of the divine as a \emph{rational being} [with
preponderance of the theoretical or the practical standpoint], \emph{pantheism},
the abandonment of the \emph{concrete} dogmatic content or of the \emph{historical side}.)

By summing this up, the question will \ed{2} become answerable:
whether the task set at the beginning of the lectures
has been realized? Elements show themselves that dominate the spiritual life of modern times
but that are foreign to Christianity, and these
elements show themselves to be the proper life-principles of modern philosophy.
\emph{This} result, then, has been gained: that there is not merely a divergence
present between Christianity and the individual systems of modern philosophy,
but that what grounds this difference is
something essential to the whole of modern times, a positive life-principle. A
multitude of phenomena outside the domain of philosophy have been pointed out
that have their root in this, and thus fall outside the Christian.

% TR: Gebet = domain
% "Fordringen paa": the claim to (ordinary sense, not the term of art "requirement")
What has been developed is confirmed \emph{further} still by the critique of the systems
that come forward with the claim to agree with
\kpage{94}Christianity. It has been shown that their endeavors to bring about
an agreement have rested on an untenable
foundation, and that what has prevented their carrying it through has been precisely
the direction indicated above, characteristic of modern times.

\sepline

% TR: Restaurationsforsøgene = the attempts at restoration; Troesphilosophien = the philosophy of faith
3.) \emph{Statement of the course of what follows}. --- a) \emph{The attempts at restoration}. Their
relation to those already criticized (the philosophy of faith, Schleiermacher.
--- Steffens). --- The \emph{reactionary} character in them: \emph{Fr. Baader}
(reaction of the medieval direction in philosophy. --- The concept of nature.
--- Theosophy.). --- \emph{Schelling}. (The consciousness of the deficiency
in the preceding idealistic philosophy; --- attempt,
instead of taking up the empirical element demanded by this deficiency,
\ed{3} to take up an element of fantasy, by which philosophy,\footnote[19]{The historical connection of the ``positive'' philosophy with Schelling's earlier
philosophy.}
in spite of all endeavors toward a logical deduction, becomes theosophy.)
--- \emph{Speculative dogmatics}. (its slight scientific value. --- It
will be treated briefly, since expositions of Hegel have already
given the necessary foundation for conceiving it. --- The contradiction in its
elements. --- In demonstrating this contradiction,
the \emph{opposite} of speculative dogmatics, the left
side of the Hegelian school, will be discussed.). --- b). \emph{The consciousness of the divergence}. α) \emph{Feuerbach}
(problems that would come under treatment in the exposition\footnote[20]{cf. 30,5. 31,1.}
of his theory.). --- β) \emph{Kierkegaard} (relation to Feuerbach.
--- Demarcation of the domains of religion and philosophy. --- The formula
for the determination of the boundary. --- Problems.).

\emph{Final result}. --- The religious and humane task of the future. ---

NB: ``\emph{Mankind has never been without religion}.'' --- Application of
this thesis to the determination of the final task. This
becomes, on the presupposition that Christianity is no longer
the dominant spiritual power, --- not: to invent a new religion;
for religions are not \emph{invented}; --- but: to \emph{find} the religious elements in
\kpage{95}the spirit of modern times and to see them in their connection and significance.

% Koch: summary of spring lectures 2--20 (Baader, Schelling, speculative dogmatics, Feuerbach) and of lectures 21--30
\ednote{}

% Koch: lecture 20 closes Feuerbach and opens Kierkegaard
\ednote{}

% Koch: a loose leaf, written on both sides, perhaps notes for the introduction in lecture 20
\ednote{}

\ed{1} \emph{Kierkegaard}.
1. The relation to what precedes (the attempts at restoration --- Feuerbach.)
General characterization of Kierkegaard's standpoint and
task. --- Plan for the exposition of his theory.

Orienting remarks on his relation to speculative
philosophy of religion, to ordinary theology, ordinary
religiousness, and to Feuerbach.

2. Kierkegaard's personal development. --- Childhood influences.
--- The Christian. --- Dialectic. --- Study of theology.
\kpage{96}--- Philosophy. --- Aesthetics. --- Relation to the political. --- His
authorship --- His personality.

3. \emph{Kierkegaard's writings}. --- The unity in them. --- The task. --- Closer
determination of it by the situation. --- ``The category of the authorship''
--- The movement. --- The method. --- ``The single individual''. --- The pseudonymity.
--- Division of the writings into 3 groups. Characterization of
these.

\ed{2} \emph{The categories of existence}.

General statement of the task and the problem. --- The historical
presupposition of Kierkegaard's standpoint.

The \emph{objective} consideration. --- Holy Scripture. --- The Church. --- The proof from
world history. --- The speculative truth. (NB. The human being
as synthesis of the temporal and the eternal.)

\emph{Subjectivity, existing subjectivity in relation to the eternal}. --- The subjective
thinker's interestedness in his thinking. --- He places everything in
becoming. --- The double reflection. --- The negative in cognition.

% "Tilværelsens System": Tilværelse (not Existents) = existence
--- The subject is ``an \emph{existing infinite spirit}.'' \ed{Cf. \emph{SKS} 7, p. 81.}
(NB. that the eternal \emph{becomes}) --- The expression for this --- the negative
form. --- The infinite striving. --- Unity of the comic and the
pathetic.) --- The impossibility of a system of existence. --- The subject-object
% Koch: above "Subject-Objectet" the manuscript has "Den rene Tænken" (pure thinking)
\ednote{}. --- The
existing thinker cannot be absorbed into the metaphysical.

\emph{Subjectivity} as the task --- The opposition between subjectivity
and the objective determination of the matter. --- Becoming a
subject as the human being's highest task. --- (NB. the world-historical
and the ethical.)

\emph{Subjectivity as the truth}.

\lecture{\ed{Lecture 20 continued:}}

\ed{65.1} (\emph{The separation of philosophy and religion}.)
2. \emph{S. Kierkegaard}.
a) By the problems to which the critique of Feuerbach's conception of religion
has led us, we are led to the last formula for the determination
of the relation of philosophy and religion that we have to examine:
the formula that has been set up \emph{by Kierkegaard}. \ins{The critique of the \emph{attempts at restoration}\kpage{97}
had shown the untenability of their attempt to
bring the view of life of Christianity and of modern philosophy
into harmony. \emph{Feuerbach} had, in extending the investigation
to religion in general, separated religion's and philosophy's} domains
in \emph{the interest of philosophy}; Kierkegaard separates them in the interest of faith, in order
to assert it as the highest form of human existence,
its most intensive expression, as \emph{the form of truth in the actuality of existence}.
\emph{Objective} science is acknowledged in its significance; it is assigned its
place in the sphere of human life, but a subordinate and limited
place. \ins{But from science's \emph{impersonal} relation to the truth
is distinguished the existing individual's \emph{subjective relation} to the truth. The inwardness
of appropriation as the truth for the existing individual, and the existing individual's}
highest truth becomes that in which the \emph{existing} individual
has his \emph{essential} interest, which he expresses in his life, that to
which he relates himself in \emph{infinite passion} out of his immersion in
existence, and this truth is for Kierkegaard \emph{Christianity}.

% TR: Tilegnelse = appropriation
% TR: Virkelighed = actuality
\emph{The task} for Kierkegaard thus becomes this: first to assert the distinction
between the truth of \emph{life} and that of \emph{science}, the \emph{subjective} and
the \emph{objective} truth, and then, when it has been asserted, when the claim of \emph{existence}
upon the individual has been decisively made good in opposition to
thought's flight from existence, to order the \emph{spheres of existence} and the
corresponding\ed{65.2} \emph{views of life}, and to point to Christianity
as the highest, eternal, revealed \ins{inaccessible to thought}
truth, ``to bring it forth in its supernatural greatness.''
\ed{cf. \emph{SKS} 7, p. 66}.

By this task his relation to the contemporary phenomena is
determined. \emph{Speculative philosophy of religion} had let religion
be a step in the development of spirit that found its completion in philosophy,
in speculative thought, which unfolded itself with logical necessity
out of the lower form. Only in unity with thought did the content of
faith find its true form; this unity was the higher truth;
in it the content of faith was penetrated by thought, converted into
the concept \ins{and the true as \emph{thought} was the truth.} In \emph{Kierkegaard}
the spheres are separated in another way: science becomes one thing, life
another; the objective unfolding of thought in the form of
objectivity, in the sphere of possibility according to the law of necessity, is divided from
the \emph{appropriation} of the truth in the individual in the form of subjectivity, in the sphere of actuality\kpage{98},
according to the law of freedom. The personal existence of the truth
becomes the highest for the existing individual, and within the personal
forms of existence the \emph{religious} becomes the highest, and within the
religious in turn what is absolutely inaccessible to thought, the \emph{paradoxical},
Christianity's specific determination, the absolute. The task of
thought here becomes not to penetrate with thought, not \ed{65.3} to
comprehend, but to hold fast to the paradox \emph{as paradox}, to use thought to
discover the opposite of thought, the absolutely unthinkable, so that one may
then believe by \emph{virtue of the absurd}.

\emph{Ordinary, unphilosophical theology} immerses itself \emph{objectively} in learned
investigations; through these it wants to bring out what Christianity
is as objective doctrine; by way of history and criticism it wants
to secure certainty about this; and during this striving, which by
the nature of the case loses itself in an \emph{infinite approximation}, it lets it be
the presupposition that the investigator is a Christian, and that what
matters is the \emph{scientific result}, whose appropriation is supposed to follow
of itself. For Kierkegaard \emph{the infinite decision of existence} is the main thing:
the approximation, by which precisely the decision is prevented,
is the deception, \emph{appropriation} the task, and the essential thing therefore
not a \emph{what} but a \emph{how}. The question is not what Christianity
historically is; but the question is: \emph{how do I become a Christian}.
The truth is \emph{the inwardness of appropriation}, the truth is \emph{subjectivity}.

% the "and in that it lets ..." clause has no main clause in the Danish; kept incomplete
\ins{\emph{Ordinary, so-called religiousness}, which has unconsciously taken up
a new age's view of life under the name of Christianity, lets
Christianity melt together with the world, become commensurable
with it. It lives in the determination of immediacy or
in the relations of civic-moral life and calls this life a Christian
life, Christian in and for itself or at any rate with an addition
of a general religious \emph{mood}, or through participation in worship,
or through belonging to the Christian Church. It naturalizes
Christianity, and in that it lets having become a Christian
\ed{65.4} have taken place in the child without a personal decision. Against this
Kierkegaard sets the opposition between Christianity
and the world, their essential dissimilarity, the impossibility of
Christianity's being assimilated by the historical process,
the necessity of suffering as its distinguishing mark in the world,
the necessity of the personality's most painful decision in order
\kpage{99}to become a Christian, the necessity of positing an absolute relation
to the absolute that is not exhausted in the relative relations
of actuality.}

\ed{65.3 continued} In Feuerbach the spheres were separated, religion
determined as the irrational, the right of thought asserted against faith,
that of sensuousness, of immediate concrete actuality, against
the fantastic idealization of the beyond. The moral
necessity of atheism was stated, Christianity \ed{65.4} posited as the
untrue, which was to be abandoned. The opposition between Feuerbach
and Kierkegaard in the conception of the truth of Christianity is
therefore the greatest possible, and yet there is no one with whom Kierkegaard
has more points of contact in determining what Christianity
is. And this has not escaped him either, and he has,
precisely on account of his conception of Christianity, not been able
to wonder at it. If the individual's relation to Christianity
is a relation of subjectivity, and more precisely one in which subjectivity
is at its highest potentiation, in its most energetic \emph{passion},
then besides the believer, who in passion holds fast to it as
truth, only \emph{he} can grasp its essence who \emph{in passion rejects it}.
The latter, standing outside Christianity, will be able to say decisively and
sharply \emph{what Christianity is}; he will be able to depict it
in such a way that the believer can acknowledge his depiction in each and every point.
It is undoubtedly Feuerbach whom Kierkegaard
has had in mind when, in Concluding Unscientific Postscript p.
473 \ed{\emph{SKS} 7, p. 558}, after having spoken of the confounding of the Christian by theology and by revivalism,
he says, ``on the other hand: a
scoffer attacks Christianity and at the same time presents it
so soundly that it is a pleasure to read him, so that one who is at a loss
to get it rightly and definitely set forth must almost resort to him.''
At this point, ahead of the exposition of Kierkegaard's
conception, I can indicate only by way of intimation\ed{65.5} his points of contact
with Feuerbach; they will become clearer in what follows.
They agree in determining the content of religion as something \emph{subjective}\footnote[21]{Absolute subjectivity (Wesen des Christenthums p. 125) -- inwardness
(Wesen des Christenthums p. 97)}
even if the kind of subjectivity is conceived differently; in designating\kpage{100}
religion's\footnote[22]{Faith's opposition to reason} standpoint as the \emph{practical} one; in having the category
\emph{the single individual} for the religious person\footnote[23]{Appropriation.} (Feuerbach Sämmtliche Werke.
I, 92f.); in holding fast to the equal relation of the \emph{incarnation} to all subsequent
times (II, 190); in conceiving \emph{the relation between Christianity and the world}
as unchangeable (I, 62. note); in letting the Christian's sympathy
be \emph{restricted to Christians}; and in denying a \emph{world} with a being
higher than the\footnote[24]{The Christians' relation to nature -- suffering as Christianity's distinguishing mark
(suffering is an object of imitation, redemption not.) -- Dying away
(Wesen des Christenthums 242) as the believer's wish. -- God is only for
me when he is believed. -- God = the absolute τελοσ, eternal happiness. -- God as
susceptible of suffering. (Concluding Unscientific Postscript p. 306 \ed{\emph{SKS} 7, p.
364f.})}
% TR: Salighed = happiness (evig Salighed = eternal happiness; the adjective salig = blessed)
% "Liden" (verbal noun of lide), not "Lidelse": "susceptible of suffering"
human. These parallels could be increased still
further, but I restrict myself to the most
prominent. ---

\lecture{Lecture 21}

b) \emph{Kierkegaard's personal development, studies and circumstances of life}. --- This I touch on
only briefly, partly because I lack the necessary material
to give anything complete, partly because Kierkegaard's personality
must still be so fresh in memory. \ins{I restrict myself to
% Koch prints "indskænker" (for "indskrænker"); rendered as the sense
the points that become significant for the understanding of his standpoint
and the peculiarity of his direction.} --- His personal
development is decisively influenced by his \emph{childhood} impressions. From
the influence of his father, a highly marked, naturally gifted
and independently developed personality, who in living religious
interest had attached himself to the party in the Church from which, in this country,
the struggle against rationalism went out, --- he had received
the direction\footnote[25]{NB. early development of the dialectical skill.}
toward \emph{the Christian in a strict form}. \ed{65.6} His study
became theology, but \ins{not in a professional way; it was really
the forms of religious \emph{life} that interested him, and with passion
he occupied himself, especially in his youth, with} comprehensive \emph{philosophical}
and aesthetic studies. \emph{Hegel's} philosophy, then the dominant\kpage{101}
system, was made the object of a strenuous and passionate
study, and from it he retained, despite all deviation from it,
the idealistic conception, the strict dialectic, and the holding fast
of \emph{speculation} as the true philosophy. From Hegel he went back
to the study of other philosophers, and it was above all the \emph{German
idealistic} philosophy \emph{from Kant}, and the \emph{Greeks}, particularly \emph{Plato} and \emph{Aristotle},
that\footnote[26]{NB. Plutarch}
occupied him. With the more recent forms of philosophy
that have arisen outside Germany, and with the oppositions
that have appeared there against idealism, he seems to have been less
familiar, and in the domain of philosophy \emph{nature} became for him a
foreign sphere. The natural sciences remained constantly indifferent to him.

Beside the study of philosophy, the aesthetic studies were
those that filled the greater part of his time. \emph{Heiberg's} activity as aesthetician
had at that time awakened the liveliest interest in this science
in Denmark, an interest that in the patriarchal time of absolutism
filled everyone, as the political interest did later. For
Heiberg as aesthetician Kierkegaard always retained reverence;
from him, no doubt, came the incitement to \ed{65.7} an extensive
occupation with poets and theoretical aestheticians. There
was in Kierkegaard the stuff of a great poet; he made this talent\footnote[27]{(Reference to his writings, especially to the 1st part of Either/Or, Fear and Trembling, and
Stages.)}
\emph{subservient} to another interest, but through it he penetrated deep
into the secrets of poetry. The \emph{Greeks} and \emph{Shakespeare} became objects
of his admiration; the \emph{Romantics} of his liveliest interest.
In Kierkegaard's first larger work, the dissertation on
\emph{irony}, the influence of Schlegel's Lucinde is recognizable, and a restrained
echo of this influence is the \emph{1st part of Either/Or},
where the life-principle of Romanticism is brought forth in existing individuality.
% Koch prints "forbybede" (for "fordybede"); rendered as the sense
But however much Kierkegaard immersed himself in aesthetic
studies, the interest that stood out ever more definitely was
not the merely theoretical one, but this: through the conception of the aesthetic
to comprehend the peculiarity of a whole sphere of life, to pose \emph{the
problems by which the aesthetic is brought into relation with personal life}, a whole
\kpage{102}series of problems which it is to his credit to have posed and
clarified. --- It was characteristic of Kierkegaard that his aesthetic
interest was concentrated essentially on \emph{poetry} and in part on \emph{music},
and that \emph{painting}, \emph{sculpture} and \emph{architecture} stood further from him.

From \emph{political} life he kept himself aloof; his task was directed not
outward but inward, and consistently, therefore, he took a polemical stance
toward political endeavors \ed{65.8} and maintained a strictly conservative
tendency, grounded not only in his task but in the veneration for royal
dignity implanted in him early, a veneration he held fast to, as he held fast
in general to reverence for whatever had made a strong impression on him in his
youth, whether in the religious or in the poetic domain.

After completing his university studies he at once began his
\emph{authorship}, an authorship laid out according to a comprehensive plan,
with one sharply defined thought in view --- as will be developed in what
follows --- and he carried it through, living, favored by outward independence,
as a private man with only one task, to whose accomplishment he applied an
energy of will whose equal will seldom be found, and, in the otium of an
artistically ordered life, all the rich, manifold abilities with which he was
endowed. His writings bear witness to the degree in which his spirit possessed,
in unity, knowledge, dialectical power, poetic imagination, humor and pathos;
those who knew him personally will be able to recall \ins{the degree in which
his thought had gained actuality in his personal life, that the ethical and the
religious were \emph{alive} in him, and that his life was a life of spirit. It
was an earnestly struggling life, but the earnestness shielded itself with jest
and irony, but through the jest the earnestness sounded for the initiated,
through the irony the humanity. His life was a struggle for the ideal.}
% TR: Virkelighed = actuality

Just as in his younger years, bold and combative, with the never-failing sharp
weapons of dialectic, irony and wit, he had fought for \emph{the ideals of
poetry} against prose and \ed{66.1} mediocrity, so it became his task in his
older years to fight for \emph{the ideal of religiousness} against the
trivialization of religion. In this struggle, which will be described in what
follows, he spent his last strength; he died young, only 42 years old, but he
lived long enough to complete his work, to accomplish the task on which he
staked his life. ---

\kpage{103}c) \emph{Kierkegaard's Writings}. --- It has already been touched on in the
foregoing how one thought, one task \emph{carries} Kierkegaard's whole
authorship. Through the whole series of his numerous writings there runs one
dialectically laid-out and artistically executed plan, whose individual moments
the individual writings realize, and this in such a way that the final goal
stood definitely before his eyes from the very beginning, and that no single
writing came to light before the next in the series was already finished. The
goal at which the whole authorship aimed was \emph{the
Christian};\footnote[28]{Cf. the Preface to Two Discourses at the Communion on
Fridays, \ed{\emph{SKS} 12, p. 281}.} even where the other views of life were
brought forward, where the scene was another sphere of existence, there was a
hidden striving toward it, and \emph{therefore the authorship, viewed as a whole,
is religious from first to last}, and Kierkegaard is right when he says that
there corresponds to this authorship as its originator\footnote[29]{NB. ``One
unnamed, whose name will one day be named'' (the Dedication of Two Discourses at
the Communion on Fridays, \emph{SKS} 12, p. 279.)} \emph{one who qua author has
willed only one thing}. This one thing, then, is \emph{the religious}, and more
precisely \emph{the Christian-religious}, and this not \ed{66.2} as an object of
theory but as \emph{a task for existence}. ``\emph{Becoming a Christian}'' is the
problem, but since this problem is posed \emph{in Christendom}, where one thus
\emph{is} a Christian, the situation puts everything \emph{into reflection}. The
author's task is then more precisely determined thus: ``\emph{Without authority
to make}\footnote[30]{Cf. Two Ethical-Religious Essays, \ed{\emph{SKS} 11, p.
105ff.} --- On My Work as an Author p. 6, note \ed{\emph{SKS} 13, p. 12, cf. also
p. 19}.} \emph{aware of the religious, the Christian}''; \emph{the movement} is
not: to reflect oneself into Christianity, but: \emph{to reflect oneself out of
something else in order to become a Christian more and more simply}!; the
religious is put wholly and entirely into reflection, but in such a way that it
is wholly and entirely taken out of reflection again, so as to be back in
simplicity, so that the path traversed is: to \emph{reach}, to \emph{come} to
simplicity.
% TR: Enfold / Enfoldighed = simplicity

The watchword becomes: ``\emph{for Christianity against the illusions}'' \ed{see
e.g. \emph{SKS} 7, p. 550}. This makes the \emph{indirect method} necessary. If
there is a misdirection present in an age, such that what is to be a task for
personal existence has become an object of impersonal consideration, then the
task of the method is not that of presenting something anew as doctrine, of
correcting by lecturing; rather it becomes \kpage{104}this: dialectically to
place the opposites together in such a way that the recipient can take
possession of what is communicated only through the independent appropriation
of reproduction, and, in doing so, comes \emph{by} \ed{66.3} \emph{himself}
into a relation to the truth, not by a direct communication into a relation to
the communicator. And if an age's deficiency is this: to live in a deception,
then the task, by the indirect method, the author concealing himself, is
\emph{dialectically to disrupt the deception} and bring the truth forward so
that the reader by himself can at least \emph{become aware}, and then choose.
The method in Kierkegaard's writings, insofar as they do not communicate the
Christian in the form of the upbuilding discourse, is therefore the indirect
one, the form the \emph{doubly reflected} one, by which the author distances
the reader from himself in order, as it were, to force him into himself.
% TR: Sandsebedrag = illusion (lit. "deception of the senses")
% TR: Tilegnelse = appropriation; Betragtning = consideration

Already here lies what is decisively given with the \emph{religious
character} of the authorship: that its movement is, not to gather a crowd, a
public, but precisely ``to shake off the public,'' to address itself to ``\emph{the
single individual}'' religiously understood. ``For religiously there is no
public, but only single individuals, for the religious is earnestness, and
earnestness is: the single individual, yet in such a way that every human
being, unconditionally every human being, as indeed he is one, can be, indeed
shall be the single individual --- \emph{the single individual before God}.''
\ed{\emph{SKS} 13, p. 17.}

Connected with this, too, is the \emph{pseudonymity}\footnote[31]{NB. the lower
--- the higher pseudonymity. (On My Work as an Author. p. 6, note \ed{\emph{SKS}
13, p. 12}.)} of all the writings that do not communicate the Christian. In
order to determine the place of Christianity, to separate it categorically from
the other spheres of existence, these had to be brought forward, and brought
\ed{66.4} forward in the form of existence. ``The production itself required,
for the sake of the dialogue, for the sake of the difference of the
psychologically varied individuality, poetically that recklessness in good and
evil, in contrition and exuberance, in despair and presumption, in suffering and
jubilation, which is ideally limited only by psychological consistency, which no
factually actual person within the moral limitation of actuality dares allow
himself or can want to allow himself.'' \ed{\emph{SKS} 7, p. 569.} The different
views of life are represented not merely by different existing individualities
but by different\kpage{105} fictional authors. The pseudonymity becomes a
further means of distancing in the relation between reader and author, a means
of distancing that is in the interest of the ideality of understanding. --- (As
a condition it requires in the author a range of abilities that comprehends the
peculiarities of all the individual spheres of life.)

Kierkegaard's writings can be divided into 3 groups: the \emph{first} would
include the series of writings \emph{from the dissertation on irony to
Concluding Postscript}. It presents in existence the different stages of life
which Kierkegaard places before the Christian, and brings the latter forward in
the form of an experiment (in the \emph{Philosophical Fragments} and in
\emph{Concluding Unscientific Postscript}). This group can be designated the
\emph{propaedeutic} one; their character is \emph{maieutic-dialectical}. The
\emph{second} group includes the \emph{upbuilding}\ed{66.5} \emph{discourses}, of
which the first appeared at the same time as Either/Or, the last toward the end
of Kierkegaard's authorship. This group could be designated \emph{the religious
communication}; its character: \emph{to be upbuilding through the ideal as
requirement}, and in it the ideal advances from the generally religious to the
Christian, yet in such a way that the latter is not presented in its utmost
rigor. The \emph{third} group contains the
\emph{Christian-polemical}\footnote[32]{cf. the accompanying leaf to On My Work
as an Author. \ed{\emph{SKS} 13, p. 23ff.}} writings: first those in which
% Koch prints "s, 23ff." (comma for full stop) in fn 32; rendered "p. 23ff."
Christianity is lifted up ``to the utmost of ideality'' \ed{cf. \emph{SKS} 12,
p. 15}, its ideal requirement brought forward in all its rigor: \emph{The
Sickness unto Death} and \emph{Practice in Christianity}, --- but in such a way
that no judgment\footnote[33]{cf. On My Work as an Author p. 19 \ed{\emph{SKS}
13, p. 25.}} is pronounced on the established order; next ``\emph{the newspaper
articles},'' \ins{i.e., the articles in ``Fædrelandet'' from December 1854} and
% Koch prints an opening » before "Bladartiklerne" that is never closed; the English closes it after the word
``The Moment,'' in which, after the established order had made the attempt to
reject the requirement, the struggle is waged for ideal Christianity against
Christendom. This group can be designated the \emph{polemical-reformatory}
group.\footnote[34]{cf. ``For Self-Examination,'' the words on the back of the
title page \ed{\emph{SKS} 13, p. 32.}}

The content of the different groups and the difference of the situation
condition a nuancing in their \emph{form} and \emph{style}. Within the
\emph{first} group is found the greatest diversity: in the presentation of the
different spheres of life that it includes there is room for the richness of
imagination, for poetic pathos, for the play of wit, for irony's\kpage{106}
sting, \ins{for the power of the comic, for the penetrating sharpness of
psychological depiction,} for the calm \ed{66.6} earnestness of the ethical, for
the inflexibility of dialectic, and for the profound jest of humor, while the
ideal rigor of the religious is glimpsed as if from a distance. --- In the
\emph{upbuilding} writings, poetic depiction, psychological reflection,
dialectic and the comic have come to serve the religious-ideal; \emph{pathos}
predominates there, but not the poetic pathos that draws the individual away
from himself out into the ideality of imagination; rather \emph{the pathos of
personality}, which, \emph{concerned} for the highest, relates itself to
\emph{the earnestness of actuality}, and with \emph{the requirement of the
ideal} penetrates all the illusions, whose ways it knows. --- In the last group
% "I den i sidste Gruppe": the first "den" is stray (a correction left standing?); rendered as the sense
of writings, psychological depiction is taken up on a large scale in The
Sickness unto Death, but as a \emph{means} of illuminating the decisive
Christian determinations; in Practice in Christianity the dialectical is taken
up on a similar scale, but in order to guard the paradox. In both, \emph{the
highest pathos} is the character. In the newspaper articles and The Moment,
dialectic and the comic are the weapon: ``laughter, administered in fear and
trembling'' \ed{\emph{SKS} 13, p. 301} is the scourge used against the
established order, but dialectic and satire are everywhere borne by the deepest
pathos, which, sharpened by the pain of sympathy, expresses itself with
shattering truth in enthusiasm and longing for death.

% Koch: 66.7-8 is an inserted leaf, apparently a rewriting of 66.9
\ednote{}

\ed{66.7} Such is the picture of him that his activity gives. But those who
knew him while he lived know a richer and fuller picture, and for them may I be
permitted further to call up the memory of his personality, to recall the clear
fullness of his thought, the rich inwardness of his feelings, the inflexible
energy of his will; --- to recall him in jest, which on his lips always let the
earnestness shine through transparently; to recall him in earnestness, when a
loving word of his enlivened, a reconciling word calmed, a releasing word
guided; --- to recall him in his youth, when, bold and combative, with the sharp
weapons of wit\footnote[35]{irony} and dialectic that never\kpage{107} failed,
he fought for the ideals of poetry against prose and mediocrity; to recall him
when older, when, with the highest goal in view, he labored in the service of
the Deity with a steel strength of will that willed only one thing; to recall
him finally in the earnest struggles of the last year, when his ``wish was
death, his longing the grave, his desire that this wish and this longing might
soon be fulfilled'' \ed{\emph{SKS} 13, p. 124}; to recall how in the suffering
of the struggle he preserved his loving concern for others, even for the
smallest circumstances of life, preserved his gentleness, his kindliness, even
his jest, preserved his equanimity of mind and clarity of thought; preserved
above all the peace and rest in faith that did not fail him even amid the
sufferings of the deathbed. So stands his picture for us, in beautiful
agreement with the noble picture that those who know him only from his writings
receive from these.

\ed{66.8} What he accomplished as a religious author is what gives him his
essential significance; it is as religious writings that his works must
essentially be assessed. But beside their worth as such they have a particular
worth, in the author's eyes no doubt vanishing and inessential, but significant
nevertheless, namely on the side of language. For in no Danish prose writer,
surely, has the Danish language, which he ``with a filial devotion, and with an
almost feminine infatuation was proud to have the honor to write,'' which he has
so ravishingly depicted (in his Stages \ed{Two Discourses at the Communion on
Fridays, \ed{\emph{SKS} 12, 281, cf. \emph{SKS} 6, p. 449f.}}), attained such
range, such flexibility, such richness in forms. Through him it found expression
for every mood, for every feeling, for every thought. Always it was fresh,
always rejuvenated in the luxuriance of mood, equally beautiful in the playful
sport of form with itself and in the plain, moving, simple utterance of the
highest truth. Like the form of the crystal, which cannot be separated from the
crystal, the form of the language always came into being with the thought,
stamped by it, inspirited by its play.

\ed{66.9} Common to all Kierkegaard's writings is their high worth in
\emph{stylistic} respect, a worth that for the author was no doubt the
inessential and vanishing, but which is nevertheless significant.
\kpage{108}In no Danish prose writer, surely, has the\footnote[36]{(Preface to
Two Discourses at the Communion on Fridays)} Danish language, ``which he with
filial devotion and with an almost feminine infatuation was proud to have the
honor to write'' \ed{\emph{SKS} 12, p. 28f.}, and which he has ravishingly
% Koch prints "s. 28f." (kept as printed); the preface quoted is SKS 12, p. 281 (cf. fn 28 and p. 107)
depicted in his ``Stages'' \ed{\emph{SKS} 6, p. 449f.}, attained such range,
such flexibility, such diversity in forms. Through him it found expression for
every mood, for every feeling, for every thought. Always it was fresh, always
rejuvenated in the luxuriance of mood, equally beautiful in the playful sport of
form with itself, in the exuberance of laughter, in the wrestling of struggle,
in the outbreathing of passion, and in the plain, moving, simple utterance of
the highest truth. Like the form of the crystal, which cannot be separated from
the crystal, the form of the language with him always came into being with the
thought and the feeling, stamped by them, inspirited by their play. The mother
tongue was for him a material that he shaped artistically, his prose a work of
art whose equal perhaps only the Greek language in Plato's writings has.

\lecture{Lecture 22}

d.) \emph{The Categories of Existence}.

d.) In the foregoing\ed{66.10} \emph{Kierkegaard's task} has in general been
designated as this: \ins{to assert \emph{the relation of faith} as the highest
by} pointing to the individual's relation to the truth \emph{in existence}; his
\emph{problem} as this: how do I, \emph{existing}, express the truth \ins{and,
since the truth is determined as Christianity: \emph{how do I become a
Christian?}} The deficiency of the age he concentrates in the words: \emph{that
it has forgotten to exist, forgotten what inwardness} is \ed{cf. \emph{SKS} 7,
p. 220}. Herein lies his opposition to speculation and to all objective
occupation with Christianity; it is from this point that one must set out in
order to grasp his view in its coherence.

The \emph{objective} consideration asks purely objectively \emph{about Christianity's}
truth, either in the sense that truth means \emph{historical}
or \emph{philosophical} truth. It does not, on the other hand, ask about
\emph{the truth of appropriation; it is not infinitely, personally interested,
in passion, with regard to its eternal happiness, in its relation to this
truth}. It is assumed that when \kpage{109}the objectively true has been
procured, appropriation follows of itself; and during the consideration by which
the truth is to be procured, the subject relates itself objectively,
\emph{disinterestedly}.
% TR: evig Salighed = eternal happiness

But it is easily seen, as far as historical truth is concerned, that the
maximum that can be reached with respect to it is \emph{an approximation}, and
that such an approximation is so heterogeneous with an eternal happiness that no
result can come of it, that \emph{the infinite decision} in passion \ed{67.1}
cannot take place through it at all. Whether the inquiry turns to \emph{Holy
Scripture} or to \emph{the Church} in order to secure the objective truth,
approximation remains the highest: attack and defense both lie in an
approximating which in its perpetually continued striving is not dialectical for
an infinite decision on which one builds one's eternal happiness. As soon as
one takes subjectivity away, and from subjectivity passion, and from passion
the infinite interest, there is reached\footnote[37]{The contradiction is
\emph{double}: partly that one wants to relate oneself objectively, i.e.,
disinterestedly, dispassionately, to what is the object of \emph{the subject's}
passionate interest, of its highest passion; --- partly that one wants to let
what at its maximum is only an \emph{approximation} to certainty become that on
which one wants to build an eternal happiness.} no decision whatsoever, neither
in this nor in any other problem. \emph{All decision, all essential decision,
lies in subjectivity}. [``Sensuous certainty, let alone historical certainty, is
uncertainty, is only approximating; the positive and an immediate relation to it
is the negative.'' \ed{\emph{SKS} 7, p. 44}.]
% TR: Afgjørelse = decision

Just as little as the objective consideration of Scripture and Church can lead
to a decision, just as little can \emph{the consideration of world history}
lead to one. Centuries cannot furnish a proof of the truth of Christianity; the
\emph{historical} consequences would always remain incommensurable with the
\emph{eternal} truth, which is more essential than all consequences, and time
becomes incommensurable with the eternal when it is to be a proof of it.
Centuries are no more a proof than one day is.
% "Aarhundreder ere det ikke mere end een Dag": "det" taken as "a proof" (Bevis)

In \emph{the speculative consideration of Christianity} the problem,
\emph{becoming a Christian}, does not come up at all. In its objectivity
\ed{67.2} speculation is precisely \emph{altogether indifferent to the
individual's eternal happiness}. For the one who speculates, the question of his
\emph{personal} eternal happiness cannot come up at all, precisely because his
task consists in getting more and \kpage{110}more away from himself and
\emph{becoming objective}, and so vanishing from himself and becoming \emph{the
beholding power of speculation} \ed{cf. \emph{SKS} 7, p. 60}. Speculation wants
\emph{in time} to be \emph{merely eternal}; in the light of eternity it wants to
contemplate Christianity, see it in its necessity and in its unity with the
eternal truth. ``But suppose now that this whole undertaking were a chimera,
suppose that Christianity is precisely subjectivity, is inward deepening,
suppose, then, that only two classes of human beings can know anything about
it: those who, infinitely interested in passion for their eternal happiness,
believingly build this on their believing relation to it; and those who in the
opposite passion (but in passion) reject it --- the happy and the unhappy
lovers. Suppose, then, that objective indifference can find out nothing
whatever.'' \ed{\emph{SKS} 7, p. 56f.} In relation to all cognition where it
holds that the object of cognition is the very inwardness of subjectivity, it
holds that the one who cognizes must be in this state. Only like is understood
by like. ``In the infinite passionate interest for its eternal happiness,
subjectivity is at its utmost exertion, at the extreme, not where there is no
object (the imperfect and undialectical distinction), but \emph{where God is
negatively in subjectivity}, which in this\footnote[38]{``The infinitely
negative is the only adequate form for the infinite'' (Stages, p. 377
\ed{\emph{SKS} 6, p. 448}.) ``The infinite negative resolution, which is
individuality's\ed{67.3} infinite form for God's being in him.'' (Concluding
Unscientific Postscript p. 21, note* \ed{\emph{SKS} 7, p. 41}.)} \ed{67.3}
interest is \emph{the form of eternal happiness}'' \ed{\emph{SKS} 7, p. 57}.
Here it holds unconditionally that he who is not in this state gains nothing by
all his considering. The fact that the human being is not merely eternal but is
\emph{a synthesis of the temporal and the eternal} means that the happiness of
speculation is an illusion, that\footnote[39]{What is amiss with speculation is
thus: 1) with respect to the relation to Christianity, its objectivity, 2) with
respect to speculation itself, its illusory eternity.} the highest is
\emph{subjectivity's infinite interest in passion for its eternal happiness};
and it is \emph{the highest} because it is the \emph{truer}, which definitely
expresses \emph{the synthesis}.\footnote[40]{(The section from here to 69,6 was
completely reworked in the oral presentation, essentially following the
enclosed leaves).
% Koch: the reworking is printed after the end of lecture 23
\ednote{}}

\kpage{111}But if now the decision of the problem cannot come about by the way
of objectivity, because it leads only into a quantifying, because it leads the
subject away from itself\ins{, because it precisely leaves out that by which the
decision becomes possible,} and because, forgetting the conditions of existence,
it wants to give the subject an illusory eternity, and if the decision is thus
referred to subjectivity, and the form of the decision provisionally indicated
by the determination of the subject as \emph{a synthesis of the temporal and
eternal}, then this concept of \emph{subjectivity, existing} subjectivity, is
now to be considered more closely; and first in general in relation to the
\emph{eternal}, to the truth, without taking closer account of the
\emph{Christian} determinations.

\emph{Objective thinking} was indifferent to the thinking subject and its
existence; \emph{the subjective thinker}, by contrast, is as existing
essentially \emph{interested in his own thinking}, is \emph{existing in} it.
Therefore his thinking has another kind of reflection, namely \emph{that of
inwardness, of possession}, by which it \ed{67.4} belongs to him and to no one
else. While objective thinking puts everything into \emph{result} and helps the
whole of humanity to cheat by copying and reciting the result and the answer,
subjective thinking puts everything into \emph{becoming}
and\footnote[41]{``Certainty is impossible for one who is becoming, and is
precisely the deception.'' \ed{\emph{SKS} 7, p. 75.}} leaves out the result,
partly because this belongs precisely to him since he has the way, partly
\emph{because as existing he is continually in becoming}, which indeed every
human being is who has not let himself be fooled into becoming objective, into
superhumanly becoming speculation. --- \emph{The reflection of inwardness}
% Koch prints "af blive" (for "at blive"); rendered as the sense
% TR: Vorden = becoming
\ed{cf. e.g. \emph{SKS} 7, p. 73f.} is the subjective thinker's \emph{double
reflection}. Thinking, he thinks the universal, but as existing in this
thinking, as acquiring it in his inwardness, he becomes more and more
\emph{subjectively isolated}. He therefore cannot, like the objective thinker,
\emph{communicate himself} directly, but only indirectly, in the form of double
reflection.
% Koch prints "Dobbelt-Reflexionen Form" (for "Dobbelt-Reflexionens"); rendered as the sense

\emph{As existing, the existing subject is in becoming; its thought must then
correspond to the form of existence}. But becoming is the alternation between
being and non-being, the positive and the negative. The \emph{positive} in
relation to thinking can be referred to these determinations: \emph{sensuous
certainty}, \emph{historical knowledge}, \emph{speculative result}. But
\emph{this positive is precisely the untrue}. Sensuous certainty is deceit,
historical knowledge is illusion
\kpage{112}(since it is approximation-knowledge), and the speculative result
is a delusion. For all this positive does not express the cognizing
subject's \ed{67.5} state in \emph{existence}; it therefore concerns a \emph{fictive
objective subject}, and to confuse oneself with such a one is to be,
and to go on being, fooled. Every subject is an \emph{existing subject}, and therefore
this must find essential expression in all its cognition, and express itself
as preventing that cognition from an illusory closure in sense-knowledge, in historical
knowledge, in illusory result. In historical knowledge he moves
constantly in the sphere of approximation-knowledge, only imagines he
has certainty, which is nevertheless had only \emph{in infinity}, which he, as
existing, cannot \emph{be} in, but can only constantly arrive at. \emph{Nothing historical
can become infinitely certain to me except this: that I exist, which is not something historical}.
The \emph{speculative} result is an illusion insofar as the existing
subject wants, thinking, to abstract from the fact that it is existing, and to be
sub specie æterni. --- Precisely because \emph{the negative} is present in existence
and is present \emph{everywhere} (for existence---Tilværelse, Existents---is constantly in becoming),
% Tilværelse and Existents both = existence; Brøchner names both here, so both Danish words are kept
% TR: Vorden = becoming
the only salvation against it is to become constantly
attentive to it. In becoming positively secured, the subject is
precisely fooled. The negativity that is in existence, or rather \emph{the
existing subject's negativity} (which his thinking must essentially
reproduce in an adequate form), \emph{is grounded in the subject's synthesis, that it is an
existing infinite spirit. Infinity and the eternal are the only certain thing, but inasmuch as
it is in the subject, it is in existence, and the first expression of this is its deception, and
this immense contradiction, that the eternal} \ed{67.6} \emph{becomes, that it comes into being}. The existing
subject's thinking must then have a form in which it can reproduce
this, and since the subjective existing thinker, who has infinity
in his soul, always has it, the form must \emph{always} express it,
and it does so by being deceptive, \emph{negative. As existing, the
subjective thinker is constantly just as negative as positive; for his positivity consists
in the continued inward deepening, in which he knows of the negative}.
He is never a teacher, but \emph{a learner}, and when he is constantly just as
negative as positive, he is constantly \emph{striving}. \ins{And this
striving is not a striving that ends when a goal is reached, but
\emph{infinite striving. Existing itself is a striving}, and it is just as pathetic
as it is comic.} In his thinking he constantly imitates his existence,
i.e. he \emph{puts all his thinking into becoming}. --- He has \emph{just as much of the comic
as of pathos}. The relative difference that in immediacy lies
\kpage{113}between\footnote[42]{``In all his thinking the thinker must also think this: that he is existing.''
\ed{\emph{SKS} 7, p. 91.}}
the comic and the tragic vanishes in the double reflection,
where the difference becomes infinite and ideality is thereby
posited. What lies at the ground of the comic and the pathetic is
the disproportion, the contradiction between the infinite and the finite,
the eternal and the becoming. The conception of the disproportion, when it
is seen looking toward the idea, is pathos; the conception of the disproportion when it
is seen with the idea behind it, is comic. If I have not exhausted the whole of the
comic, neither do I have the pathos of infinity; if I have the pathos of infinity,
I also at once have the comic.
% TR: Misforhold = disproportion

[Beyond the negative relation to the eternal, an existing individual can never
be helped by the historical: ``contingent historical truths
can never become proofs of necessary truths of reason''
\ed{\emph{SKS} 7, p. 92}. Here \ed{67.7} a \emph{leap} is required, by which the decision
comes into being.]

\lecture{Lecture 23}

That the existing individual \emph{is constantly in becoming}, i.e. is constantly \emph{striving},
entails as its consequence that for him \emph{no system of existence}
can be given, even though there can be a \emph{logical system. System and conclusiveness
correspond to each other, but existence is precisely the opposite of conclusiveness}. For
\emph{God} existence is a system, but it cannot be one for an existing
spirit, who can gain a conclusiveness outside existence only
fantastically. \emph{In order to think existence, systematic thought must think it
as annulled, hence not as existing}. Existence is the spacing-apart, which
falls asunder: the systematic is conclusiveness, which closes
together. For the existing individual, even if he has reached the highest,
\emph{repetition}, by which he must indeed fill out his existence if
he is not to go backward (or become a fantastical being), will again
be a continued striving, because here again conclusiveness is put at a distance and
deferred. \emph{Continued striving} is the upshot for the existing subject's
ethical view of life. The systematic idea is the \emph{subject-object}, a
unity of thinking and being; existence, by contrast, is precisely separation\kpage{114}.
% TR: Afsluttethed = conclusiveness; det Spatierende = the spacing-apart
From this it by no means follows that existence is thoughtless, but it has
spaced and spaces the subject apart from the object, thinking from being. \emph{Objectively
understood}, thought is \emph{pure thinking}, which just as abstract-objectively
corresponds to its object, which object is therefore again itself, and
\emph{truth} is thought's agreement with itself. This
% Koch prints "obj{jective}" (doubled j); rendered as the sense
objective thinking has no relation to existing subjectivity
and ends by becoming, --- if it \ed{67.8} is, incidentally, possible for
the subject to slip into this objectivity, and if it is not something
that at most it can come to know of through the imagination,
--- the purely abstract co-knowledge in, and knowledge of, this pure relation
between thinking and being, this pure identity, indeed this \emph{tautology},
because by being it is not said that the thinker is, but strictly
only that he is a thinker. The \emph{existing} subject, on the other hand, is existing,
and that indeed every human being is. And \emph{this demand of existence prevents one
% Koch prints "der vi blive deri" (for "vil"); rendered as the sense
from remaining in the metaphysical, strikes the one who will remain in it with the comic}.

If in this way \emph{the existing individual's relation to truth} is now
determined as a \emph{relation of subjectivity}, then the determination now gained,
\emph{subjectivity's} determination, becomes at the same time the determination of
the highest \emph{task} that is possible for the existing individual. ``Objectively
one speaks constantly about the matter; subjectively one speaks about the subject and
subjectivity, and behold, \emph{precisely subjectivity is the matter}. It must constantly
be held fast \emph{that the subjective problem is not something about the matter, but is
subjectivity itself}. For since the problem is the \emph{decision}, and the decision,
as shown above, lies in subjectivity, what matters is \emph{that objectively
no trace of any matter at all remains}; for in the same moment subjectivity will
shirk some of the pain and crisis of decision, i.e.
will make the problem a little objective, and thereby the decision is postponed.''
\ed{\emph{SKS} 7, p. 121.}
% TR: Afgjørelse = decision; Bestemmelse = determination

\emph{Becoming a subject} \ins{, subjectivity's development or transformation,
its infinite concentration in itself over against the representation of
infinity's highest good, an eternal happiness, (subjectivity's
infinite concern about itself.)} \emph{is the highest task set for every
human being, and if it were not, then ethics would have to despair}. Ethics and
the ethical have an irrefutable claim on every existing individual by
being the essential stay of individual \ed{68.1} existence; such an
irrefutable claim that whatever a human being accomplishes in the world,
the most astonishing thing, it is nonetheless dubious if he has not ethically
% Forestilling = representation (Vorstellung), here and p. 120
% TR: evig Salighed = eternal happiness
\kpage{115}been clear with himself when he chose, and ethically clarified his choice. Ethics
must therefore regard \emph{objective world-historical knowledge} with suspicious
eyes, because it easily becomes a snare, a demoralizing aesthetic
distraction for the knowing subject, since the distinction between
what becomes world-historical and what does not is quantitatively
dialectical, which is also why the absolute ethical distinction between\footnote[43]{Only \emph{in himself} can the human being study the ethical --- if he seeks it in history,
he never comes to act. Speculation in ethics wants a now-living individual
to act by virtue of a theory of immanence, i.e. to act by virtue of refraining
from acting.}
\emph{good and evil} is, world-historically, --- aesthetically neutralized in the
aesthetic-metaphysical determination: \emph{the great}, the significant,
to which the bad man has equal access with the good. Ethically seen,
that by which the individual becomes world-historical is \emph{the contingent}, but
ethics regards as\footnote[44]{``The ethical is the only certain thing; to concentrate oneself upon this, the only knowledge that
does not at a last moment possibly change into a hypothesis; to be in it, the
only secure knowledge, where knowledge is secured by something else.'' \ed{\emph{SKS} 7, p. 142.}}
unethical the transition by which one lets go of
the ethical quality in order to try oneself in the quantifying other.
Constant preoccupation with the world-historical makes one unfit to
act, since regard for the \emph{outcome} displaces the will's ethical enthusiasm,
and makes the individual unethical. The true \emph{ethical enthusiasm} lies
in willing to the utmost of one's ability, but at the same time, uplifted in divine
jest, never thinking about whether one thereby accomplishes anything or not.
--- Since a human being cannot by himself, by freedom, by the will to
the good, become a world-historical figure, --- this is namely
dependent on something else, hence contingent --- it is unethical to
be concerned about it. The true ethical individuality would, even if
the governing power so arranged circumstances that he became a world-historical
figure, by a resolution of the will remain ignorant
\ed{68.2} of it, because the external is not in his power, and therefore
means nothing either pro or contra; he \emph{would} not, even in death,
know that his life had had any significance other than that of
having ethically worked through his soul's development. Significance
must always be apprehended as something Governance adds to his ethical striving
in itself, hence not as the fruit of his work; it is a pro
that must be regarded as a spiritual trial fully as much as any contra\kpage{116}.
% TR: Styrelsen = Governance; Begeistring = enthusiasm
The proposition: \emph{world history is the world's judgment}, can, when it is to
express an \emph{ethical} consideration of life, in any case have validity for
God, because in his eternal co-knowledge he possesses the medium that is precisely
the commensurability of the outward and the inward. But the
\emph{human} spirit cannot see the world-historical in this way, even
apart from the difficulties: \emph{a}) that access to becoming world-historical
is quantitatively dialectical, so that what has become world-historical
has passed through this dialectic; \emph{b}) that the world-historical
consideration, as an act of cognition, is an approximation, subject to
the same dialectic as every conflict between idea and empirical reality,
which at every moment will hinder the beginning, and when it has begun,
at every moment threatens a revolt against the beginning; --- \emph{c})
that in world history, being the history of the race, the ethical cannot
be seen, but only the \emph{abstract}, the metaphysical relation of immanence
between cause and effect, ground and consequence. (NB.
besides, the slowness of world history)

\emph{Becoming subjective} is thus the highest task set for a
human being, just as \ed{68.3} the highest reward, \emph{an eternal happiness}, exists only
for the subjective person, or rather \emph{comes into being} for the one who becomes subjective.
And becoming subjective shall give a human being plenty to do
as long as he lives (cf. the examples in Concluding Unscientific
Postscript p. 121ff.) \ed{\emph{SKS} 7, pp. 153--67.}

We are now led by this to the determination of \emph{subjective truth}.

If truth is defined as \emph{agreement between thinking
and being}, and if by being is understood \emph{empirical} being, then
truth itself is transformed into a desideratum and everything is put in becoming, because
the empirical object is not finished, and because the existing
cognizing spirit is itself indeed in becoming, and thus truth is an
approximating, in which the beginning cannot be absolutely posited, precisely
because there is no conclusion, which has retroactive force;
whereas every beginning (if it is not, by being conscious of this,
an arbitrariness), when it is \emph{made}, does not take place by virtue of
immanent thinking, but is \emph{made} by virtue of a \emph{resolution,} essentially
by virtue of faith. --- If, on the other hand, in that definition one takes \emph{being abstractly}
as the abstract rendering of, or the abstract prototype
for, what being in concreto is as empirical being, then nothing
% "abstract bestemmes abstract": the adverb doubled as printed; kept
prevents truth from being abstractly determined abstractly as \emph{something finished}\kpage{117};
for the agreement between thinking and being is, \emph{abstractly
seen}, always finished, since the beginning of \emph{becoming} lies precisely in \emph{concretion},
from which abstraction abstractly looks away. But if being is understood
in this way, the formula is a \emph{tautology}, thinking and being mean one
and the same, and the agreement becomes the abstract identity
with \ed{68.4} itself. All that is expressed is that truth is not
something single, but, quite abstractly, a \emph{duplication}, which is nevertheless annulled in the same
moment. Abstraction can go on paraphrasing this
as much as it likes; it never gets any further. As soon as the being of truth
becomes \emph{empirically concrete}, truth itself is \emph{in becoming}, is indeed again,
\emph{as divined}, agreement between thinking and being, and is
indeed so actually \emph{for God}, but is not so for any \emph{existing
spirit}, because this spirit, existing, is itself in becoming.

For the existing spirit qua existing spirit the question of truth
now comes back. Existence holds the two moments
in truth's abstract reduplication apart from each other, and
reflection shows two relations. For \emph{objective} reflection
truth becomes an objective something, an object, and the point is to look
away from the subject; for \emph{subjective} reflection truth becomes
appropriation, inwardness, subjectivity, and the point is precisely,
existing, to immerse oneself in subjectivity. \emph{Mediation}, which would
% "der vel forbinde" as printed (for "der vil forbinde"?); rendered as the sense
connect these two forms, only leads us back to abstraction; it
does not explain how the eternal truth is to be understood \emph{in the
determination of time} by one \emph{who, by existing, is himself in time}.
% TR: Tilegnelse = appropriation

The way of \emph{objective} reflection makes the subject the \emph{contingent}, and
thereby existence into an \emph{indifferent} vanishing thing. The way to
objective truth goes away from the subject, and while the subject and
subjectivity become indifferent, truth becomes so too, and
precisely this is its objective validity; for interest, like
decision, is subjectivity. The way of objective reflection
now leads to \emph{abstract thinking}, to \emph{mathematics}, to \emph{historical knowledge} of various
kinds; \ed{68.5} it constantly leads \emph{away from the subject}, whose existence
or non-existence becomes, objectively quite rightly, infinitely indifferent.
At its maximum it will lead to the \emph{contradiction} that only objectivity
has come into being and subjectivity has gone out, that is to say, the
\emph{existing subjectivity} that has made an attempt to become what in the
abstract sense is called subjectivity, abstract objectivity's
\kpage{118}abstract form. And yet the objectivity that has come into being is,
subjectively seen, at its maximum either a \emph{hypothesis} or an \emph{approximation},
because all eternal decision lies precisely in subjectivity.

\emph{Subjective} reflection turns inward toward subjectivity,
and wants in this \emph{inward deepening} to be the truth, and in such a way that, just as
in the foregoing, when objectivity had been brought forth, subjectivity
had vanished, here \emph{subjectivity itself becomes the last thing and the objective
the vanishing}. Here it is not forgotten for a single moment that the subject is existing,
and that existing is a becoming, and that therefore that identity of truth,
of thinking and being, is a chimera of abstraction,
and in truth is only the longing of creation, not because truth
is not this, but \emph{because the cognizer is an existing individual, and so
truth cannot be this for him as long as he exists}. Only \emph{momentarily}
can the existing subject be in a unity of finitude and infinity
that goes beyond existing. This moment is \emph{the moment of
passion. Passion is, for an existing individual, precisely the highest of existence}. In passion
the existing subject is infinitized in the eternity of imagination
and yet at the same time precisely most definitely itself.

All \emph{essential cognizing} relates itself to existence\ed{68.6}, i.e. it
relates itself to \emph{the cognizer, who is essentially an existing individual}. Only
\emph{ethical} and \emph{ethical-religious cognizing} is therefore essential cognizing.
But all ethical and all ethical-religious cognizing is essentially relating itself
to the fact that the cognizer exists.

When the question of truth is asked \emph{objectively}, one reflects objectively
on truth as an object to which the cognizer
relates himself. One reflects not on the \emph{relation}, but on its being
\emph{the truth}, the true, that he relates himself to. When what he relates
himself to is only the truth, the true, then the subject is in the truth.
When the question of truth is asked \emph{subjectively}, one reflects
subjectively on the individual's relation; if only the \emph{how} of this relation
is in truth, then the individual is in truth, even if he were thus
relating himself to untruth (NB. here the talk is of essential truth,
or of the truth that relates itself essentially to existence.).
--- \emph{The cognition of God} may serve as an example. \emph{Objectively} one reflects
on this: that it is \emph{the true God}. God is then to be brought forth objectively
through the whole approximating deliberation, \emph{which is never achieved,
because God is subject, and therefore is only for subjectivity in inwardness}. \emph{Subjectively} one reflects
\kpage{119}on this: that the individual relates himself to a something \emph{in such a way} that
his relation is in truth a God-relationship. The existing individual who
chooses the subjective way grasps in the same moment the whole dialectical
difficulty in his having to spend some time, perhaps a long
time, finding God objectively; he grasps this dialectical difficulty
in all its pain, because he must make use of God in the same moment,
because every moment is wasted in which he does not have God. In the same
moment he has God, not by virtue of any objective deliberation,
but by virtue of \emph{inwardness's}\ed{68.7} \emph{infinite passion}. (God thus becomes
a postulate, but this must be understood in such a way that the existing
individual's postulating God is a \emph{necessity}.)

Objectively the accent falls on \emph{what} is said, subjectively on \emph{how} it is said.
At its maximum this how is \emph{the passion of infinity}, and
\emph{the passion of infinity is the truth itself}. But the passion of infinity
is precisely \emph{subjectivity}, and thus \emph{subjectivity is the truth}.
% Koch prints "saaledet" (for "saaledes"); rendered as the sense
Objectively seen there is no infinite decision, and thus it is
objectively correct that the difference between good and evil is annulled
along with the principle of contradiction, and thereby also the infinite
difference between truth and falsehood. Only in subjectivity is there decision,
whereas wanting to become objective is untruth. \emph{The passion of
infinity is the decisive thing, not its content, for its content is precisely
itself. Thus the subjective how and subjectivity are the truth}.
Truth can then be defined thus: \emph{the objective uncertainty, held fast in
the appropriation of the most passionate inwardness, is the truth, the highest truth
there is for an existing individual} \ed{cf. \emph{SKS} 7, p. 186}. But this determination
of truth is a paraphrase of \emph{faith}.

In the proposition that subjectivity is the truth, the \emph{Socratic
wisdom} is contained, which paid heed to the fact that the cognizer is existing.
To this determination corresponds the objective determination
of truth as \emph{paradox}. The paradox is \emph{the objective uncertainty}, which is
the expression for the passion of inwardness, which is precisely the truth.
\emph{The eternal essential truth is the paradox inasmuch as it relates itself to an existing individual}.
The existing individual's \emph{ignorance} is the expression for the objective
uncertainty; but the fact that cognizing is determined as \emph{recollecting} points
to this: that the eternal essential truth \emph{is not in and for itself} \ed{68.8} \emph{paradoxical,
but only by its relation to the existing individual}. (NB. Socrates's relation to
Plato.)

\kpage{120}If one wants to seek a more inward expression than the Socratic for the
proposition: subjectivity is the truth, then one must begin with:
\emph{subjectivity is untruth}, or, in other words, with the Christian
decisive determination of \emph{sin}, which annuls the possibility of recollectively
taking oneself back into eternity. \emph{The individual becomes a sinner by coming
into being}. The way back is barred; he must go forward. The individual must now \emph{in
existence} get hold of the truth, but then the \emph{eternal} truth must
come into being \emph{in time}, and in doing so it \emph{itself} becomes the paradox. The
strongest expression for inwardness now emerges, and it is:
when the retreat out of existence, recollecting, into eternity has been made
impossible, then \emph{with the truth against one as the paradox, in the anxiety of sin and with
this pain, with the immense risk of objectivity, to believe --- by virtue of the absurd}
(The absurd is that the eternal truth has come into being in time, that God
has come into being as this individual human being). \emph{Speculation is here cut off}
(it remains in \emph{immanence} and volatilizes the paradox.) The absurd,
held fast with the passion of infinity, is \emph{faith in the strict sense}. The paradox
is not what is to be explained; what is to be grasped, only more and more
deeply, is \emph{that it is the paradox}. Just as God, as the highest representation,
cannot be explained by anything else, but can only be explained by
immersing oneself in the representation itself, and just as the supreme principles
of all thinking can only be proved indirectly (negatively), so the
paradox, as the ground of an existing individual's relation to an eternal essential
truth, will not be explicable by anything else either, when
the explanation is to be for existing individuals. --- \emph{Faith relates itself self-actively
to the paradox, self-actively in discovering it, and at every moment holding it fast --- in order to
be able to believe}. \ed{69.1} (``Speculation is \emph{objective}, and \emph{objectively there is no
truth for an existing individual, but only approximating, for he is prevented from becoming entirely objective
by existing}. Christianity, by contrast, is subjective;
the inwardness of faith in the believer is subjectivity's eternal decision.
And objectively there is no truth; for the objective knowledge
of the truth of Christianity, or of its truths, is precisely
untruth; to know a confession of faith by heart is paganism,
because Christianity is inwardness''. \ed{\emph{SKS} 7, p. 204})

Now that truth has been determined as subjectivity, \emph{the actuality
of subjectivity} must be considered. \emph{Abstraction} looks away from actuality
and from its difficulty, which is this: that \emph{this determinate something} and \emph{the ideality
of thinking} are to be put together. The \emph{particular}, the \emph{contingent}, is something
% TR: Virkelighed = actuality
\kpage{121}belonging to everything actual and directly opposed to abstraction. \emph{In
the domain of abstraction the principle of contradiction is annulled, but in existence it is
an impossibility to annul it, for then existence itself is annulled}. When I take existence
% Koch prints "selve Existentsens" (genitive); rendered as the sense, "existence itself"
away (abstract), then there is no aut-aut; when I take it
away in existence, then this means that I take existence away, but
then I do not, after all, annul it in existence. For the existing individual, who as
such is \emph{in becoming}, the eternal is not the eternal but \emph{the future}.
When he, acting, relates himself to this --- acting as in immanence
--- he relates himself to it in \emph{infinite passion}, and then an
\emph{aut-aut} is posited. Where everything is in becoming, where there is only so much of
eternity present that it can \emph{hold back} in the passionate decision,
there where eternity relates itself as futurity to the one
becoming, there \emph{the absolute distinction} belongs. \emph{For when I put
eternity and becoming together, I get not rest but futurity}.
% TR: det Tilkommende = the future; Tilkommelse = futurity

\ed{69.2} To exist in truth, thus with consciousness to penetrate
one's existence, at once as it were far out beyond it,
and yet present in it and yet in becoming: that is difficult. To
think existence sub specie æterni and in abstraction is essentially
to annul it. Existence cannot be thought without motion,
and motion cannot be thought sub specie æterni. Insofar as all
thinking is nevertheless eternal, the difficulty is there for the existing individual.
Existence, like motion, is a very difficult thing to deal with.
If I think it, I annul it, and then I do not think it. It
might then seem right to say that there is something that cannot be
thought: existing. But the difficulty is there again, \emph{that existence
puts it together in that the thinker exists}. --- Insofar as existence
is motion, it holds that there is nevertheless something \emph{continuous} that
holds motion together; for otherwise there is no motion.
The \emph{unmoved}\footnote[45]{``The existing individual thinks momentarily, he thinks beforehand and he thinks afterward.
Absolute continuity his thinking cannot attain.'' \ed{\emph{SKS} 7, p. 300.} ``Existence
separates the ideal identity of thinking and being. I must exist in order to
be able to think, and I must be able to think (e.g. the good) in order to exist in it.'' \ed{\emph{SKS} 7,
p. 301.} ``Existence is always the particular; the abstract does not \emph{exist}.'' \ed{\emph{SKS} 7,
p. 301.} ``God does not think, he creates. God does not exist, he is eternal. \emph{The human being thinks
and exists, and existence separates thinking and being}, holds them apart from each other in succession.''
\ed{\emph{SKS} 7, p. 303.}}
belongs with motion as motion's
\kpage{122}goal, both in the sense of τελοσ and of μετρον; otherwise the claim that all is in
motion, if one at the same time wants to take time away and say that all is
always motion, is eo ipso standstill. The difficulty for the existing individual
is now this: \emph{to give existence the continuity without which everything simply
vanishes}. An \emph{abstract continuity} is no continuity, and the fact
that the existing individual exists essentially hinders continuity,
whereas \emph{passion} is the \emph{momentary} continuity, which
at one and the same time holds back and is the impulse of the motion. For an existing individual
the goal of motion is \emph{decision and repetition. The eternal is the continuity
of motion}, but an abstract eternity is outside motion,
and a concrete eternity in the existing individual is the maximum of
passion. \emph{For all idealizing passion is an anticipation} \ed{69.3} \emph{of the
eternal in existence for an existing individual to exist}. But passion's anticipation
of the eternal for an existing individual is still not the abstract
continuity, but \emph{the possibility of approximation to the only true one there
can be for an existing individual} (NB. The eternity of abstraction and of pure thinking,
Concluding Unscientific Postscript p. 236f.
\ed{\emph{SKS} 7, p. 284ff.})

% TR: Bevægelse = motion
% TR: Continuerlighed = continuity
% TR: Virkelighed = actuality; Mulighed = possibility
% TR: ophæve = annul (not Hegel's sense here; sublate where Hegel's aufheben is meant)
\emph{For the existing individual, existing is his highest interest, and interestedness
in existing is actuality}. What \emph{actuality} is cannot be stated in
the language of abstraction. Only by annulling actuality can abstraction
get hold of it, but to annul it is precisely to transform
it into \emph{possibility}. All that is said about actuality in the language of abstraction within
abstraction is said within possibility.
Actuality, existence, is the dialectical moment in a trilogy
whose beginning and end cannot be for an existing individual, who
qua existing is in the dialectical moment. Abstraction joins
the trilogy together, but the one abstracting is, after all, an existing individual, and
as existing therefore in the dialectical moment, which he cannot
mediate or join together, least of all absolutely, as long as
he is existing. When he nevertheless does so, then this must relate
as a possibility to actuality, to existence, where he himself is.
He must explain how he goes about it, i.e., how he
\emph{as existing} goes about it, or whether he ceases to be
existing, and whether an existing individual \emph{has a right to do so}. At the same moment we
begin to ask in this way, we ask \emph{ethically} and assert the ethical's
claim upon the existing individual, which cannot be that he should abstract\kpage{123}
from existence, but that he should exist, which is also the
existing individual's highest interest.

\emph{All knowledge about actuality is a possibility; the only actuality of which an existing individual
is more than knowing is his own actuality, the fact} \ed{69.4} \emph{that he exists, and
this actuality is his absolute interest}. Abstraction's requirement of
him is to become disinterested in order to come to know something; the ethical's
requirement of him is to be infinitely interested in existing. \emph{The
only actuality there is for the ethical individual is his own ethical actuality; all other actuality
he only knows about; but true knowledge is a translation into possibility}. ---
The reliability of \emph{sense perception} is deception. The reliability \emph{that
knowledge of the historical} wants to have is also only deception, insofar as it
wants to be the reliability of actuality, since the knower only
knows about a historical actuality when he has dissolved it into possibility.
\emph{Abstraction} is possibility, the preceding or the subsequent.
\emph{Pure thinking} is a phantom. \emph{Actual subjectivity is}
not the knowing subjectivity, for by knowledge he is in the medium of possibility, but
is \emph{the ethically existing subjectivity}\footnote[46]{NB. The self-relation in relation to knowledge.} (this thinks everything in relation to itself,
infinitely interested in existing.) (NB. The different rank of actuality and possibility
in the domain of the aesthetic and of the ethical, --- \emph{faith's}\footnote[47]{``The object of faith is the actuality of the god in existence, i.e., as a single individual,
i.e., that God has existed as an individual human being.'' \ed{SKS 7, p. 298} --- ``Christianity
is no doctrine, but the fact that the god has existed.'' \ed{SKS 7, p. 298.}``Actuality
is not the external action, but is interiority, in which the individual
annuls possibility and identifies himself with what is thought in order to exist in it. This
is action.'' \ed{\emph{SKS} 7, p. 310.}}
relation to an actuality that is not one's own.
% fn 47: Guden = "the god" (Climacus's usage), distinct from Gud = God
% TR: Indvortesheden = interiority
--- ``\emph{Actuality is ideality}. But aesthetically and intellectually ideality is
possibility. Ethically, ideality is actuality in the individual
himself. \emph{Actuality is interiority infinitely interested in existing, which the
ethical individual is for himself}'' \ed{\emph{SKS} 7, p. 296} --- ``\emph{The individual's own ethical actuality
is the only actuality}.'' \ed{\emph{SKS} 7, p. 522} \ins{(to everything else he has
only a relation of possibility, precisely as a maximum.)} ``\emph{Subjectivity is
truth, subjectivity is actuality}.'' \ed{\emph{SKS} 7, p. 314.})

In the existing actual subjectivity the task is \emph{to give
the individual moments of subjectivity simultaneity}, and this simultaneity is
the opposite of the speculative process. In the direction of existence,
\kpage{124}thinking \ed{69.5} is not at all higher than imagination and feeling, but they
are coordinate. The task is not to raise the one at the other's
expense, but the task is equality, simultaneity, and the
medium in which they are united is in existing.

For a subjective thinker there is required \emph{imagination, feeling, dialectic in existence-inwardness
with passion}. But first and last \emph{passion}, for it is
impossible, existing, to think about existence without coming into passion,
because existing is an enormous contradiction, which the
subjective thinker has not to abstract from, but to remain in.
The subjective thinker is a dialectician in the direction of the existential;
he has passion of thought for holding fast the qualitative disjunction;
he has this absolute disjunction as the final
decision, which prevents everything from spreading out into a quantifying.
He has it at hand, but not in such a way that, by abstractly
recurring to it, he hinders existence. He therefore also has
aesthetic passion and ethical passion; thereby concretion
is gained. He is, however, not a poet, even if he is \emph{also} a poet,
not an ethicist, even if he is also an ethicist, but \emph{also} a dialectician and
essentially \emph{himself existing}, whereas the poet's existence is inessential
in relation to the poem, and likewise the ethicist's in relation to the doctrine
he expounds, and the dialectician's in relation to the thought. He is not
a man of science but an \emph{artist. Existing is an art}. The subjective
thinker is aesthetic enough for his life to acquire aesthetic content,
ethical enough to regulate it, dialectical enough to master it
in thought. His task is: \emph{to understand himself in existence. Every human being
must essentially be assumed to be in possession of what essentially belongs to
being a human being}. The subjective thinker's\ed{69.6} task is to transform
himself into an instrument that clearly and definitely expresses
\emph{the human in existence}. ---

% Koch: two folded sheets (7 written pages) inserted after 67.2
\ednote{}

% TR: Salighed = happiness (evig Salighed = eternal happiness)
\ed{1} \emph{for Lecture 23} (Sheets 67.3--69.6) (Recapitulation of the foregoing:
The problem: to become a Christian. --- The relation to an eternal happiness
(cf. ``Philosophical Fragments.'') --- If the problem is posed objectively, the
decision becomes impossible.) But when by the way of objectivity no
decision can be reached, we are referred to \emph{subjectivity}.

% Tilværelse (below: "outside existence", "system of existence") = existence, as Existents
\kpage{125}First, then, the general relation between subjective and objective
thinking is to be stated. The latter points away from the individual, strives
toward the objective finished result; the former leads back to the individual
so that he can exist in what is thought, and since existence, which always
is in becoming, becomes its domain, it places all thinking in becoming. The
subjective thinker's thought must therefore express the determination of existence:
becoming, --- which, since becoming is the alternation between
the positive and the negative, being and non-being, means: it
must have the moment of negation in it. In other words: the ``positive,''
secured knowledge after which the objective relation strives
is in none of its forms reliable for the existing individual (cf.
67.4f and 69.4); the existing individual, who is a synthesis of the temporal
and the eternal, an existing infinite spirit, --- cannot hold fast
the infinite and eternal, which is the only certain thing --- in the element of infinity,
but since he must hold it fast in the element of finitude\ed{2},
in existence, the highest relation cannot become secured
certainty, but the infinite interest in passion, and the only
certainty about the infinite becomes that which the passion of faith
grasps out of objective uncertainty. The existing individual is thus in his
relation to the truth as much positive as negative; he is constantly
learning, constantly striving, and this striving is an infinite
striving, which lasts as long as existence lasts; for existing itself
is a striving. As existing, the subjective thinker cannot
want to attempt to win that conclusiveness outside existence
from which a system of existence could be formed; for existence,
which has an irrefutable claim upon him, is precisely the
opposite of conclusiveness, and in order to be thought, existence must be thought
as annulled. Nor can he, and nor can he want to, identify
himself with pure thinking's metaphysical identity of subject-object,
the unity of thought and being in the element of abstraction;
he cannot want to fly over existence in order to reach speculation's
illusory eternity. The claim of existence will prevent him from remaining in
the metaphysical, and it strikes whoever wants to remain there with the
comic. (NB. The consequences of this for speculative philosophy
and for ``the science of faith\ed{3}.'' --- From the thought of immortality
% Danish prints "Uødeligheden" (for "Udødeligheden"); rendered as the sense
an example is taken of what has been developed. \emph{Moments}: the immortality
that is asked about is not a fantastic eternity sub specie
\kpage{126}æterni, but the immortality of the mortal (not the immortality of the
human race). --- ``The question of immortality is a question of inwardness,
which the subject, by becoming subjective, must put to
himself. Objectively the question cannot be answered at all, because objectively
one cannot ask about immortality at all, since immortality
is precisely the potentiation and highest development of developed subjectivity.
Only by rightly willing to become subjective can the
question rightly come forth: do \emph{I} remain, or am \emph{I} immortal.''
\ed{\emph{SKS} 7, p. 160} --- Systematically, immortality cannot be proved;
for, seen systematically, the whole question is nonsense. Immortality
is subjectivity's most passionate interest; in the interest
lies precisely the proof. With immortality the point lies in subjectivity
and in subjectivity's subjective development. --- Immortality
in general does not exist at all. --- Through the constant self-relating
to the alternating that is existence, the determinateness of immortality must
become indeterminateness; the highest for the existing individual is to
stake all his passion on the indeterminateness and infinitely passionately
relate himself to the indeterminateness of the determinateness, and this is the
only way in which he as existing can know
about his immortality, because as existing \ed{4} he is so strangely
composite that the determinateness of immortality can be had in determinateness
only by the eternal one, but its determinateness can be had by the existing individual only
in indeterminateness.)

It is thus the assertion of existence, in which the eternal
is only as bound up with the temporal, that posits faith's subjectivity-relation
to the eternal as the highest, since it becomes the only
one possible for an existing individual. But if this subjectivity-relation is
the highest, then the task is accordingly: to \emph{become subjective}, and since what
is posited as the object of a task is what every human being
might seem to be without further ado, becoming a subject acquires
the\footnote[48]{NB: ``to choose oneself \emph{according to one's eternal validity}'' \ed{cf. e.g. \emph{SKS} 3, p. 203.}} definite meaning: to become what the human being is according to his
essence, to let\footnote[49]{cf. Concluding Unscientific Postscript p. 21, note* \ed{\emph{SKS} 7, p. 41.}} subjectivity be developed, be transformed within itself over against
the eternal. Only for such a subjectivity, which is the developed
possibility of subjectivity's highest possibility, does
% "Subjectet, det er henvist til sig selv": the "det" read as relative ("which is referred to itself")
% TR: Betragtning / Betragteren = observation / observer
\kpage{127}the truth come into being. The subject, referred to itself, receives the task:
to be infinitely concerned about itself; but this infinite
concern about itself is the direction toward the eternal in relation to
existence, and it leads toward the ethical and the ethical-religious,
\ins{the essential truth, and} the only essential cognizing,
and this ethical and ethical-religious the subject must seek in
itself; for only there, and not outside itself, and above all not in the objective
observation of the \ed{5} historical, can the individual find it
(NB with regard to the observation of the historical: the transformation
of the categories of ethics. --- The accidentality of its world-historical
significance. --- The outcome. --- \ins{The impossibility for the observer
of coming to act.} --- The task of the ethical individuality.
--- The impossibility for the human spirit of seeing world history
as an ethical world order.)

% Danish prints "Opgaven, Medens" (comma, then capital); rendered as two sentences
% fn 50 reads "i at existere i Virkeligheden" where p. 122 has "i at existere Virkeligheden"; rendered as printed
``To understand \emph{oneself in existence}'' then becomes the expression for the task.
Whereas abstraction wants to look away from actuality and its difficulties,
the subject's\footnote[50]{``For the existing individual, existing is his highest interest, and interestedness
in existing in actuality.'' \ed{\emph{SKS} 7, p. 286.}}
task accordingly becomes to immerse itself in existence,
to put together the ideality of thought and the singularity of existence,
with consciousness to penetrate one's existence, at one and the same
time as it were far above it and yet present in it and yet in
becoming. --- The difficulty for the existing individual becomes to give existence
the continuity without which everything simply vanishes. The
eternal is what gives continuity, and idealizing passion
is an anticipating of the eternal in existence for an existing individual to
exist. That anticipation of the eternal by passion is still not
the absolute continuity, but the possibility of the approximation
to the only true one there can be for an existing individual. (cf. 69.2f.)

Actual subjectivity is not the knowing but the ethically existing
subjectivity. The only actuality there is for the
ethical individual is his own \ed{6} ethical actuality; all other actuality he\footnote[51]{``The individual's own ethical actuality is the only actuality.'' \ed{\emph{SKS} 7, p. 522.}} only
knows about, but true knowledge is a translation into possibility.
--- The determination of actuality as ideality (69.4). --- Subjectivity
is truth, subjectivity is actuality.
\kpage{128}(NB. The simultaneity of the moments of subjectivity in the existing
actual subjectivity. --- The condition for a subjective thinker. ---
``Existing is an art''. \ed{\emph{SKS} 7, p. 321.})

In the foregoing, the subjectivity-relation to the eternal has been posited as
the highest, and the form of subjectivity determined. There now arises again
the question: has not precisely by that determination something still been posited
to which the subject relates itself, thus something objective, a matter,
and by this question we are led to the final decision of the relation
between subjectivity and objectivity. --- (here is inserted
67.8). --- Closer elaboration of the determination that truth is \emph{in
becoming}. (Truth defined as agreement between thinking
and being, either in such a way that being is taken as empirical being,
or abstractly as the abstract reproduction of, or the abstract
prototype for, what being in concreto is as empirical being. In the
first case only an approximating is reached; in the latter case
the formula is a tautology. As soon as being is empirically concrete,
truth itself is in becoming, and can for an existing spirit only \ed{7}
be \emph{divined} as agreement between thinking and being). ---
The necessity of \emph{a choice} between the subjective and the objective
way, the impossibility of a mediation. (68.4). --- The objective direction
ends in a contradiction (68.5 cf. 69.1). --- The subjective direction.
The vanishing of objectivity. Passion as the highest of existence
(68.5). --- The cognition of God from the objective and the subjective
standpoint (68.6) --- The how of the relation is what is essential
in subjectivity. In its maximum this how is the passion of infinity,
and the passion of infinity is the truth itself. But
the passion of infinity is precisely subjectivity, and thus
subjectivity is truth. --- Definition of truth. Faith.
(68.7). --- The Socratic wisdom. Truth as the paradox --- The
Socratic relation to the Christian (68.7f.).

Summing up of what has been developed, a glance toward speculation and
dogmatics. --- The determinations of existence: existence --- becoming ---
synthesis of the eternal and the temporal --- continued striving. --- Subjectivity
--- inwardness --- passion --- the passion of infinity ---
\ins{The single individual} --- faith --- truth as the paradox. NB. the domain of objective
knowledge.

\lecture{\kpage{129}\ed{69.6 continued} Lecture 24}

% TR: Umiddelbarhed = immediacy
% TR: Lykke / Ulykke = good fortune / misfortune
e) \emph{Division of the spheres of existence}. --- Since \emph{existing} has been determined as
the task for the individual, various possible \emph{views of life} appear,
various possible forms of existence, whose relation to one another must
be made clear. Kierkegaard distinguishes \emph{3 spheres of existence}: the aesthetic,
the \emph{ethical}, the \emph{religious}, and as a confinium between the first two he places
\emph{irony}, between the last two \emph{humor}. The \emph{metaphysical} forms no
sphere of \emph{existence}, since no human being exists metaphysically. The metaphysical,
the ontological, \emph{is}, but it \emph{does not exist}; for when it exists,
it is in the aesthetic, in the ethical, in the religious, and when it is,
it is the abstraction of, or a prius for, the aesthetic, the ethical,
the religious.
% Danish prints "Æstehetiske"; rendered as the sense

Aesthetic existence is that of \emph{immediacy}. The aesthetic in a
human being is \emph{that by which he is immediately what he is}; he who lives in and
with and from and for the aesthetic in him lives aesthetically. Within
this sphere the difference can be manifold and great,
but they all possess immediacy, in this: \emph{that spirit is not determined as
spirit, but immediately determined}. Even at the stage where spirituality
shows itself, spirit is not determined as spirit, but as \emph{gift}.
The view of life at all these stages is \emph{enjoyment}, but thereby is posited
a condition \emph{outside} the individual, or a condition \emph{in} the individual such
\emph{that it is not through the individual himself}. Immediacy\ed{69.7} is \emph{good fortune}, and
good fortune its view of life. The opposition comes \emph{from outside}, since
immediacy is \emph{undialectical}; the opposition is then sensed as
\emph{misfortune}, but is not grasped in its \emph{essential} significance.\footnote[52]{NB. The domain of poetry --- The poetic ideality.} Misfortune is \emph{an occurrence}
in relation to immediacy (fate); seen ideally (in the direction
of immediacy's view of life) it is gone, or it must
go.

The \emph{various stages} within the aesthetic can now be nuanced
according as either \emph{1}) the personality is immediately determined,
not spiritually but \emph{physically} (health, beauty posited as
the highest good); --- or \emph{2}) the condition for enjoyment is posited \emph{outside
the individual} (wealth, honor, love etc.); or \emph{3}) the condition
is posited in the individual, but \emph{not through the individual itself} (e.g. the development of talent); ---
\kpage{130}or \emph{4}) the task is posited as this: \emph{to live for one's desire} (here the multiplicity
becomes an object of reflection, but this reflection is \emph{finite},
and the personality remains in its immediacy. --- In desire the
individual is immediate, is \emph{in the moment}). --- ; or \emph{5}) enjoyment is determined
as enjoyment of oneself, \emph{enjoyment of oneself in enjoyment} (Epicureanism,
Cynicism), --- or \emph{6}) when \emph{grieving} becomes the meaning of life, when
the individual enjoys himself in sorrow in egoistic idealizing; --- or \emph{7})
when \emph{despair} itself is made a view of life, but as \emph{despair of thought}.
% TR: Tungsind = melancholy; Veemod = wistfulness
Also in this simple\footnote[53]{NB. The standpoint of irony.} form of aesthetic view of life, which has to
a certain degree taken up into itself the consciousness of the nothingness of such a view of life,
the personality remains in its immediacy. Here too
life is a passionate \emph{seizing of the moment}, \ed{69.8} in which \emph{continuity
is lacking}, an alternation of energy and indolence, an arbitrary
holding fast of the particular and the accidental. The \emph{mood} in the individual is
between wistfulness and transient sorrow. Common to all the aesthetic
stages is thus\footnote[54]{The one who lives aesthetically develops \emph{with necessity}, not with freedom. ---
\emph{He is in the moment}, does not acquire continuity, \emph{not transparency}.}
this: \emph{that spirit has not taken itself out of its immediacy,
has not become conscious of itself according to its eternal validity}. Where this movement has
begun but been broken off, \emph{melancholy} appears. But common to all
aesthetic stages is also \emph{despair}, conscious or unconscious. The
immediate, on which the aesthetic view of life builds, is the
accidental, the \emph{transitory}. Its annihilation, the non-appearance of
the external condition, of good fortune, the coming of misfortune, brings the aesthetic
individual to despair; but the fact that he \emph{comes to} despair
shows only that he \emph{was} in despair, for the transitory has \emph{become
manifest to him as what it was}. The difference, then, is that before he \emph{was not
conscious} of being in despair. \emph{But every aesthetic view of life}\footnote[55]{The aesthetic individual is \emph{undialectical} in himself, and has his dialectic outside himself. Existing
is one thing, the contradiction something else, which comes from outside.}
\emph{is despair,
and everyone who lives aesthetically is in despair, whether he knows it or not}.
% Danish prints "æstetisk" (without h); rendered as the sense

But \emph{despair is the means of healing}. When I despair \emph{over the finite},
the transitory,\footnote[56]{``To \emph{will one's despair} in the infinite sense is identical with \emph{absolute devotion}.''
\ed{cf. \emph{SKS} 3, p. 212.} --- ``Finite despair is hardening, absolute despair
is infinitization.'' \ed{\emph{SKS} 3, p. 212.}}
and in despair separate myself from it, I find\kpage{131}
\emph{the infinite}, the \emph{eternal human being}. But \emph{despair is a choice}. I
can doubt without choosing it, but I cannot despair without
choosing it, and when I despair, I choose again, \emph{and choose myself, not
in my immediacy, not as this accidental individual, but in my eternal validity.
Despair is the doubt of the personality}. Its condition is \emph{the will}, \ed{70.1} and
when one has in truth willed to despair, then one is in truth \emph{beyond
it}; one has found one's \emph{self}, the personality is set at rest, the absolute
won, in myself, \emph{in my freedom}. And with this the transition is
formed to another view of life, another sphere of life: the \emph{ethical}.

By \emph{the choice}, then, \emph{the ethical} is posited. The choice in the stricter sense is
therefore the \emph{ethical} choice, not the aesthetic, which is no choice.
The only \emph{absolute either/or} there is, is the choice between good and
evil, but that is also absolutely ethical. In the first instance, however, the ethical
choice is not the choice \emph{between} good and evil, \emph{but the choice by which the distinction:
good and evil, is chosen or excluded}.

From the standpoint of ethics, which is the standpoint of \emph{self-assertion}, the
choice is seen as \emph{the self's own deed}; it endures despair and finds
itself immanently\footnote[57]{NB. the infinite resignation (cf. Fear and Trembling.) \ed{cf. \emph{SKS} 4, p. 132 ff.}} in despair. The choice is then determined more closely
thus: when I \emph{choose absolutely}, I choose despair, and in despair
% fn 58 "det Ethiskes Forhold til Immanentsens": elliptical, read as "to that [the relation] of immanence"
I choose\footnote[58]{(This is developed further orally, and reference is made to the relation of the ethical to
that of immanence.)} \emph{the absolute}; for \emph{I myself am the absolute, I posit the absolute
and am myself the absolute}; but as entirely identical with this I must
say: \emph{I choose the absolute that chooses me, I posit the absolute that posits
me}; for if I do not remember\footnote[59]{to choose oneself = to receive oneself.} that this second expression is just as absolute,
then my category: to choose, is untrue, for it is precisely the identity of
both. What I choose I do not posit, for if it were not
posited, I could not choose it; and yet, if I did not posit it
by choosing it, then I would not choose it. It \emph{is}; for if
it were not, I could not choose it; it \emph{is not}, for it first comes to be
by my choosing it, and otherwise my choice would be an illusion. --- When
I choose the absolute, I choose \emph{myself in my eternal validity}. Something
\emph{other than myself} I can \ed{70.2} never choose as the absolute: for
if I choose something else, then I choose it \emph{as a finitude}, and choose
\kpage{132}it therefore \emph{not absolutely}. But this \emph{self} of mine is at once the
most abstract and the most concrete of all: it is \emph{freedom}. The self that
chooses has a multiplicity within it, insofar as it has a \emph{history}, in
which I acknowledge the identity with myself. In this history
I stand in relation to \emph{the whole race}, and this history contains
something painful, and yet I am who I am only through this history.
None of it can I give up when I would choose myself, but then
the expression for \emph{this acquiring of myself is repentance. I repent myself back into
myself, back into the family, back into the race, until I find myself in God}. Only
when I choose myself as guilty do I choose myself, if I am
to choose myself at all in such a way that \emph{choosing} oneself does not
become identical with \emph{creating} oneself. Since the self has come forth
through a choice, it is the consciousness of \emph{this particular free being}, which is \emph{itself}
and no other. It contains within it a rich concretion, a
multiplicity of determinacies, of qualities; in short, it is \emph{the whole}
aesthetic self that has been chosen ethically, that has been \emph{taken over by itself}. When I
\emph{think}, I also infinitize myself, but not absolutely; for \emph{I
vanish in the absolute}; only when I absolutely \emph{choose} myself do I infinitize
myself \emph{absolutely}; for only myself can I choose absolutely, and
this absolute choice of myself is \emph{my freedom}, and only when I have absolutely
chosen myself have I posited an \emph{absolute difference}: that between \emph{good
and evil. The good is through my willing it, and otherwise it is not at all; it is the in-and-for-itself,
posited by the in-and-for-itself, and} \ed{70.3} \emph{this is freedom}. As
for \emph{evil}, repentance is the expression for its belonging to me
essentially, and at the same time for its not belonging to me essentially. If
it did not belong to me essentially, I could not choose it, and if there were
something in me that I could not choose absolutely, then I would not be choosing
myself absolutely at all, then I would not myself be the absolute,
but only product.
% TR: uendeliggjøre = infinitize
% TR: det Iogforsigværende = the in-and-for-itself (an und für sich)

% Koch prints a stray underscore in "det,_hvorved"; not reproduced
Through the ethical choice there thus takes place \emph{the movement of infinity in} the individual;
the ethical is \emph{that whereby the human being becomes what he becomes}; what
despair brings forth is the universal \ins{and \emph{continuity}
becomes the form of the ethical life (NB. \emph{disclosure} in the universal).}
% TR: det Almene = the universal
% TR: Aabenbarelse = disclosure (aabenbar = disclosed)

He who chooses himself \emph{according to his freedom} has not chosen himself \emph{abstractly} in a
contentless eternity; however concrete his self may be, he has chosen
himself according to his possibility; in repentance he has ransomed himself in order to
remain in his freedom, but he remains in it only \emph{by constantly realizing}\kpage{133}
\emph{it}. He is therefore eo ipso \emph{acting}. One can choose oneself
according to one's freedom only when one chooses oneself ethically, but ethically one can
choose oneself only by \emph{repenting} oneself, and only by repenting oneself
does one become concrete, and only as a \emph{concrete} individual is one a \emph{free} individual.
When one chooses oneself abstractly, one does not choose oneself ethically. Only
when in the choice one has \emph{taken over oneself}, put oneself on, penetrated
oneself totally, so that every movement is accompanied by the
consciousness of a \emph{responsibility for oneself}, only then has one \emph{chosen oneself
ethically}, only then has one \emph{repented oneself}, only then is one concrete, only then
% Koch prints "abs{solut}" (double s); rendered "absolute"
is one, in one's total \emph{isolation}, in absolute \emph{continuity with the actuality
to which one belongs}. When I choose myself repentantly, I gather
myself together in my whole finite \ed{70.4} concretion, and in having thus
chosen myself out of finitude, I am in the most absolute continuity
with it.
% TR: Virkelighed = actuality (virkelig = actual)

The ethically chosen actual concretion is the individual's \emph{task. Task
is the category of ethics}, it leads immediately to action, it expresses
the individual's independence of the external, his self-assertion.
The condition of the task is thoroughgoing consciousness, \emph{the individual's
transparency to himself}.
% TR: Selvhævdelse = self-assertion

The task for the ethical individual is, more precisely expressed, this: since his
self, as immediate, is accidentally determined, \emph{to work the accidental together
with the universal}. He who lives ethically, he expresses in his life the
universal: he makes himself \emph{the universal human being}, not by taking
off his concretion, for then he becomes nothing at all, but by
putting it on and penetrating it with the universal. But
\emph{this universal the individual must himself be}, otherwise the ethical falls as an \emph{abstraction}
outside him and cannot be realized. In \emph{conscience} lies
this secret, \emph{that the individual life is at once an individual
life and also the universal}, if not immediately as such, then still
according to its possibility.

The \emph{self} that the ethical individual knows through an act of
self-scrutiny is at once the \emph{actual self} and the \emph{ideal self},
which the individual has outside itself as the image in whose likeness it is to
form itself, and which it nevertheless, on the other hand, has within itself, since it is that
self. If one does not hold fast to the individual's having the ideal self \emph{in
himself}, then his thoughts and aspirations become \emph{abstract}. The task becomes: \emph{to
transform oneself into the universal individual}, but this is possible only when I
\kpage{134}κατα δυναμιν have this \emph{in me}. For the universal can very well subsist
with and in the distinctive without consuming \ed{70.5} it; it is like
the fire that burned without consuming the thornbush. In the
act of despair the universal human being came forth, and is now behind the concretion
and breaks forth through it. \emph{The accidental in the individual is ennobled by
the choice}.

The possibility of \emph{fulfilling duty}, the universal, is present only when
the individual is at once conceived\footnote[60]{
Ethically regarded, the single individual becomes \emph{at once moment and the whole}.}
as the universal and the particular.
If I posit duty as something \emph{external}, then the difference between
good and evil is annulled, for when I am not myself the universal, I can
enter only into an \emph{abstract} relation to it, but the difference between good
and evil is incommensurable for an abstract relation. When the personality
has felt the intensity of duty with all his energy, then he is
ethically matured, and duty will then break forth in himself. As soon as
the personality has found himself in despair, has absolutely chosen
himself, has repented himself, he has himself as his task under
an eternal responsibility, and thus duty is posited in its absoluteness. Since,
however, he has not \emph{created} himself but \emph{chosen} himself, duty is
the expression of his absolute \emph{dependence} and absolute \emph{freedom} in
identity with each other. The \emph{particular duty} he will learn \emph{by himself}, and in vain
seek information about it from anyone else, and yet here again he will
be an \emph{autodidact} just as he is a \emph{theodidact}, and conversely. Duty
then in no case becomes something abstract for him, partly because it is not
something external to him; for when that is the case, it is always abstract,
partly because he himself is concrete, for when he chose himself ethically, he chose
himself \emph{in his whole concretion} and renounced the abstractness of arbitrariness.

Since the \emph{universal} is \emph{the norm of ethics}, the consequence is \emph{that nothing universal can be given up
except by virtue of a resolution}, for which the individual must at every moment
bear the responsibility; that: \emph{not to realize the universal}, is, from the standpoint
of ethics \ed{70.6}, \emph{the imperfection that is to be repented}. It follows further
from this that \emph{it is every human being's duty to have a calling}; for this is
the expression of there being a rational order of things, in which every
human being, if he will, fills his place in such a way that he
at once expresses the universally human and the individual.
\kpage{135}\emph{In the universal the human being becomes disclosed: disclosure} is therefore the
commandment of ethics.

In relation to the aesthetic view of life and sphere of life, the
distinctive character of the ethical is, then: \emph{freedom, choice of personality according to
its eternal validity, acting, infinity, continuity, disclosure}.
The sphere of action is \emph{the inner: the individual's transformation according to the}\footnote[61]{
the ``\emph{accomplishing of one's work}'' \ed{cf. \emph{SKS} 3, p. 281} is the ethical individual's highest, but
this work is not identical with the external result, the outcome.}
\emph{ethical
ideal}. The external, \emph{the outcome}, is the ethically indifferent: \emph{the resolution, the willing
of the good}, the ethical reality. \emph{The requirement of infinity} is present at every
moment in its absoluteness, but before this requirement the individual
collapses, and \emph{the ethical shows itself to be a transitional sphere}. Faced with the requirement
% Koch prints "Indidvidet" (for Individet); rendered as the sense
the individual comes to know itself as \emph{guilty}; in \emph{repentance} it now seeks
to regain itself; but \emph{this it cannot do by itself}. When the individual\footnote[62]{
The ethical individual is dialectical inwardly in itself, in \emph{self-assertion}, in such a way \emph{that
therefore the ultimate} ground does not become dialectical in itself, since the underlying self is used
to overcome and assert itself. (the determination of the religious is self-annihilation
before God.)}
despairs, it brings itself to despair, and therefore it can indeed
by itself despair of everything, but \emph{not}, when it does this\emph{, come back
by itself}. In the moment of decision it needs \emph{a higher assistance}, and
here the \emph{religious} breaks forth from the ethical. ---\footnote[63]{
``The ethical is only a \emph{transitional sphere}, and therefore its highest expression is repentance
as a \emph{negative action}.''\ed{\emph{SKS} 6, p. 439.}}
% Koch prints a stray underscore in "høiere_Bistand"; not reproduced
% TR: Gjennemgangssphære = transitional sphere
% TR: Selvtilintetgjørelse = self-annihilation

The acting of the ethical was \emph{inward} acting, but the acting of inwardness
must show itself as \emph{suffering}; for the individual is not able to\footnote[64]{
``While aesthetic existence is essentially \emph{enjoyment}, ethical existence essentially
\emph{struggle and victory}, religious existence is \emph{suffering}, and not as a transitional moment, but
as constantly accompanying''. \ed{\emph{SKS} 7, p. 263.}}
create itself anew, therefore \emph{to suffer} is the highest acting in the inward.
The human being's highest inward acting is to \emph{repent}. To repent is
not a \emph{positive movement outward} or toward; but a \emph{negative movement
inward},\footnote[65]{
cf. The Concept of Anxiety p. \emph{130}. \ed{\emph{SKS} 4, p. 419f.} Fear and Trembling p. \emph{106}. 107f. \ed{\emph{SKS} 4,
p. 187f., p. 188f.}}
\ed{70.7} \emph{not a doing, but a letting oneself by oneself undergo}\footnote[66]{
``The \emph{ethical} already posits a kind of relation of opposition between the outward
and the inward, insofar as it posits the outward in indifference, --- --- the \emph{religious}
posits the opposition definitely between the outward and the inward, definitely as
opposition, therein lies precisely suffering as the existence-category for the religious,
but therein lies precisely also the inner infinity of interiority, inwardly.'' \ed{\emph{SKS} 7,
p. 270f.}}
\emph{something. Repentance}\kpage{136},
when it is held out, but in any case by means of \emph{the sympathetic-dialectical}
in it, leads on to \emph{the religious}, to \emph{despair over the possibility of
self-assertion, to self-annihilation before the Deity}.
% TR: Indvorteshed = interiority (kept apart from Inderlighed = inwardness)
% "føre hen til": the Danish has the infinitive, a finite verb or modal wanting; rendered "leads on to"

% Koch: describes the leaf on which Brøchner summarized lecture 24
\ednote{}

\ed{1} \emph{Lecture 24. Recapitulation of the previous lecture}. (The
assertion of existence against speculation. --- The relation of faith as the highest.
--- The task: to become subjective. --- The ethical and the ethical-religious
as the essential cognizing and the essential truth. ---
The actual subjectivity. --- Actuality. --- The choice between the
subjective and the objective way. --- Definition of truth. ---
The Socratic and the Christian. --- The categories of existence. --- The significance
of the Kierkegaardian polemic.)

\emph{The spheres of existence}. 1) \emph{the aesthetic}. --- The determination: immediacy.
--- Enjoyment as view of life. --- Good fortune. (the undialectical in immediacy.
The conception of misfortune.). --- The various nuances of the
aesthetic. --- What is common to them. --- Despair in the aesthetic
view of life.

% Koch prints a stray underscore in "evige_Gyldighed"; not reproduced
2) \emph{the ethical}. --- Despair as a means of healing: in despair
I find \emph{the infinite}, the \emph{eternal human being}. --- Despair is a \emph{choice}
(NB. the significance of the will): I choose myself in \emph{my eternal validity}. ---
Through the choice the \emph{ethical} comes forth. The choice is in the stricter sense
always ethical (the choice by which the distinction: good and evil, is chosen
or excluded.) . --- In the choice I choose \emph{myself as the absolute}
(NB. the dialectic of the choice). The choice brings forth \emph{the eternal human being}, the
\emph{universal, freedom}. --- The determination of the \emph{self}. --- \emph{Repentance} (NB. the standpoint
of ethics as the standpoint of self-assertion.) --- The return of the aesthetic
self. --- \emph{Good and evil} (the relation of repentance to
evil). Summing-up of the determinations of the choice in relation to the
\kpage{137}aesthetic. \ed{2} The \emph{concretion} of the self as won by the choice and repentance.
The self's \emph{taking over of itself}. --- The self as \emph{task for itself}. Working together
of \emph{the accidental and the universal}. The ethical individual expresses in his
life the \emph{universal human being}. --- The being of the universal in the individual. --- The
\emph{actual} self and the \emph{ideal} self. --- The fulfillment of \emph{duty}. --- \emph{Realization
of the universal}. The exceptions to this. --- Having a \emph{calling}. --- \emph{Disclosure}.

Comparison of the aesthetic and the \emph{ethical}. The transition of the \emph{ethical}
to the \emph{religious}. --- Comparing of the basic determinations of the ethical and the
religious.

\sepline

At the end of the lecture the work: \emph{The Point of View for My Work as an
Author} is discussed. Attention is drawn to a) its relation to
the work: On My Work as an Author, --- b) the relation of the communications
given in it to my presentation of Kierkegaard's personal
development; --- c) the apparent discrepancy between
my division of his writings and his own, --- a discrepancy
which is nevertheless removed on closer consideration of the matter.

\lecture{\ed{70.7 continued} Lecture 25}

The first expression of \emph{religious} existence, which recalls the ethical's
starting point, \ins{infinite resignation,} is \emph{the absolute}\footnote[67]{
NB. The relation to the ethical.}
\emph{direction}
\emph{(respect) toward the absolute τελοσ,}\footnote[68]{
infinite resignation. \ed{cf. \emph{SKS} 4, pp. 123--49.}}
\emph{expressed in action in the transformation of existence}.
The \emph{pathos} that is hereby posited is not the aesthetic, which
expresses itself in the word, and in its truth can signify that the individual
abandons himself in order to lose himself in the idea, whereas \emph{existential pathos}
arises through the idea's relating itself transformingly to the individual's
existence. Poetic pathos is imagination with the ideality of
possibility, existential pathos relates itself to the ideality of actuality,
and in relation to actuality the \emph{deed} is the highest pathos.

% Koch prints "Existende" (for Existerende); rendered as the sense
Existence is composed of infinity and finitude, the existing individual
is infinite and finite. If now an \emph{eternal happiness} is his highest\kpage{138}
good, this means that the moments of finitude are once
and for all, in action, reduced to what must be given up in relation
to the eternal happiness. If it does not transform his existence absolutely for him,
then he is not relating himself to an eternal happiness; if there is something
he will not give up for its sake, then he is not relating himself to an
eternal happiness. ``\emph{To will absolutely}\footnote[69]{
(the finite cannot be willed absolutely.) cf. Concluding Unscientific Postscript
p. 323. \ed{\emph{SKS} 7, p. 384f.}}
is to will the infinite, and to will
an eternal happiness is to will absolutely, because it must be capable of being willed at every
moment. And therefore it is so \emph{abstract}, and, aesthetically seen, the
poorest idea, because it is \emph{the absolute τελοσ for one willing who
absolutely wills to strive} --- and not thoughtlessly imagine himself finished, and
not foolishly engage in haggling, whereby he only loses the
absolute telos; and therefore in a finite \ed{70.8} sense it is foolishness,
precisely because in an infinite sense it is the absolute τελοσ. \emph{And
therefore the one willing does not even want to know anything about this τελοσ,}\footnote[70]{
``For a finite being the \emph{negative infinity} is the highest, and the positive a
dubious reassurance.'' (Stages. p. 345.) \ed{\emph{SKS} 6, p. 411.}}
\emph{other than that
it exists}; for as soon as he gets to know something about it, he already begins
to be slowed down in his course.'' \ed{\emph{SKS} 7, p. 359.} --- The absolute τελοσ is
only when the individual relates himself absolutely to it, and as an eternal
happiness relating itself to an\footnote[71]{
``In temporality the \emph{expectation} of an eternal happiness is the highest reward'' \ed{\emph{SKS} 7, p.
366.} --- ``The first true expression for relating oneself to the absolute τελοσ is \emph{to renounce
everything}; but if in the same moment the retrogression is not to begin, then the true
understanding of it is \emph{that this renouncing of everything is nevertheless nothing, if it were to merit the highest good}.''
\ed{\emph{SKS} 7, p. 368.}}
existing individual, they cannot possibly get
each other or rest in belonging to each other in existence, i.e. in temporality;
the \emph{whole} time, existence, is the time of being in love. Mediation gets
no place here; for the existing individual is prevented from gaining such a
foothold outside existence that from it he could mediate what,
moreover, by being in becoming, withdraws itself from completion. In
relation to the absolute τελοσ, that there is mediating means precisely
that the absolute τελοσ is reduced to a relative τελοσ. \emph{Nor is it
true that the absolute τελοσ becomes concrete in the relative ones}; for the absolute
distinction of resignation will at every moment secure the absolute
τελοσ against all fraternizing. It is true that the individual with the direction\kpage{139}
toward the absolute τελοσ \emph{is in the relative ones}, but he is not in these
in such a way that the former exhausts itself in them. It is true that for God and
for the absolute τελοσ we are all equal; but it is not true that
for me or for a single individual \emph{God or the absolute τελοσ} is equal
to everything else. The relation to the absolute τελοσ cannot be said
to exhaust itself in the relative ends, since the absolute relation
can require the renunciation \emph{of}\footnote[72]{
NB.}
\emph{them all}. On the other hand, he who relates
himself to the absolute τελοσ can very well \emph{be in the relative ones} precisely in order,
by renouncing, to practice the absolute relation. The individual is not to
go out of the world, --- for \emph{externality is the untrue expression} --- the task
is to be that the individual strives, while practicing the absolute
relation to the absolute τελοσ, to reach \ed{71.1} the maximum:
\emph{to relate himself at one and the same time to the absolute τελοσ and to the relative ones} --- not by
mediating them, but by relating himself \emph{absolutely} to his absolute
τελοσ and \emph{relatively} to the relative ones. To relate oneself now and then to
one's absolute τελοσ is to relate oneself to it relatively, but to relate oneself
to it relatively is to relate oneself to a relative τελοσ, for \emph{the relation}
is what is decisive. The task, then, is \emph{to practice one's relation to the absolute
τελοσ so that one constantly has it with one, while one remains in the relative
ends of existence}. In the absolute respect for the absolute τελοσ
a yawning chasm is fixed between it and the relative ones. \emph{The individual does indeed
% Koch prints "sit liv Liv" (a doubled word); rendered as the sense
live in finitude, but no longer has his life in it}. As in the great
moment of resignation there was no mediating, but \emph{choosing}, so the task is
to acquire proficiency in repeating the choice of passion, and,
existing, to express it, in that he no longer exists \emph{immediately}
in finitude. In existence the word is always: forward, and
as long as the word is: forward! the point is to practice \emph{the absolute distinction}.
% "svælgende Dyb befæstet": echoes Luke 16:26 ("a great gulf fixed")
% TR: indøve = practice
% TR: Færdighed = completion (p. 138) / proficiency (p. 139): the Danish plays on one word

\emph{Since there is an absolute difference between the human being and God, the human being's
absolute relation to God must express this}. The human being expresses himself most perfectly
when he absolutely expresses the difference. \emph{Worship is the
maximum of a human being's God-relationship}, and thereby of his likeness to
God, since the qualities are absolutely different. But worship means
precisely \emph{that God is everything to him}, and the one who worships is in turn the
one who absolutely distinguishes. \emph{The one who absolutely distinguishes relates himself to his}
\kpage{140}\emph{absolute τελοσ, but eo ipso also to God}. The absolute distinction
is not between finitude and infinity, but between \emph{existing
finitely and infinitely}. For infinity and finitude are put together in
existing and in the existing individual, who is not to trouble himself
about creating existence or about imitating existence in thought,
but all the more about existing. Existing, then, he is
not to form \ed{71.2} existence out of infinity and finitude, but,
composed of finitude and infinity, he is to \emph{become} one of the two;
and one does not \emph{become} both at once, as one \emph{is} both
by \emph{being} an existing individual. Ethics has not the medium of being but of \emph{becoming},
and, regarding every existing individual as its serf,
it will constantly forbid him to begin on an abstracting from
existence, whereby becoming is exchanged for being.

\emph{Only through the venture of relating himself absolutely to the absolute τελοσ does
the individual become infinitized; through the venture he becomes another}. The existence of the \emph{absolute
ethical good} can be expressed only by the individual himself,
existing, expressing that it exists. The existing individual's relation to the
eternal is \emph{the relation of ignorance, of expectation, to what is to come}. Precisely
this gives the highest pathos. The eternal happiness, as the
absolute good, can be defined solely and only \emph{by the mode of acquisition},
nothing else can be said of it than that it is \emph{the good that is attained
by absolutely venturing everything}.
% Tilværelse ("the good's Tilværelse ... that it er til"): existence as being-there, not Existents
% TR: evig Salighed = eternal happiness

Since existential pathos was determined as action, as the transformation of
existence \ins{(by this transformation it becomes more
and more concrete.)}, and the task is posited as this: at one and the same time to
relate oneself absolutely to one's absolute τελοσ and relatively to the relative ones,
then, since this task is to be realized in existence, the beginning
becomes this: \emph{to practice the relation to the absolute τελοσ by renunciation}, and thereby
to \emph{take the power away from immediacy}. For the actual individual is in
immediacy, and to that extent properly \emph{absolutely in the relative ends}.
That renunciation is the condition for its being possible to begin on
the \emph{ideal} task. But even when the individual has overcome immediacy,
it is nevertheless, with the victory, once again in \emph{existence}, and thereby once again
prevented from \emph{absolutely} expressing the absolute relation to an absolute
τελοσ.

The acting that belongs to existential pathos \ed{71.3} now receives its
essential religious expression in being determined as \emph{suffering}. The
\kpage{141}essential existential pathos relates itself to this: essentially
existing; and essentially existing is \emph{inwardness}, and \emph{the acting of inwardness
is suffering}; for the individual is not able to \emph{create} itself anew, and therefore
to suffer is the highest acting in the inward. Inwardness
(the ethical and ethical-religious individual) grasps suffering \emph{as the essential},
the religious individual constantly has suffering with him, demands suffering in the same
sense as the immediate individual demands good fortune, and demands
suffering and has suffering even if misfortune is not there in the outward.
(The religious address has essentially \emph{to uplift through suffering}.) \emph{Religiousness
is inwardness, and inwardness is the individual's relation to itself before
God, its reflection into itself, and it is precisely from this that suffering comes}. --- The
\emph{religious suffering} is \emph{actual} suffering, and by the actuality of suffering is understood
its \emph{continuance} as essential for the pathetic relation to an
eternal happiness, so that suffering is not deceitfully revoked or the individual
goes beyond it. (Aesthetically, suffering relates itself as accidental
to existence.). Nor is the actuality of suffering revoked
by its acquiring significance for the eternal happiness, so that
suffering could thus seem to turn into joy; --- for the \emph{existing individual}
cannot make the dialectical conversion by which suffering is converted
into joy; for his knowledge of the eternal happiness is to be held fast
\emph{in the medium of existence}, not in eternity, and the perfection of joy must
then fail, since it is to be had \emph{in an imperfect
form}. The individual cannot, at the same time as it suffers religiously,
be beyond suffering in joy over this suffering's significance as the relation;
for suffering concerns precisely this, \emph{that he is separated from joy},
but it also signifies\ed{71.4} the relation, so that being without suffering
signifies that one is not religious. --- As for an existing individual
the highest principles of thinking can be proved only negatively, and
wanting to prove them positively at once betrays that the existing individual, insofar as
he is after all surely an existing individual, is in the process of becoming fantastical:
so also for an existing individual the existence-relation to
the absolute good can be determined only by \emph{the negative}: the relation to
an eternal happiness by \emph{suffering}, just as the certainty of faith, which relates
itself to an eternal happiness, is determined by uncertainty (NB. the significance of
\emph{spiritual trial} for religious suffering. ``Spiritual trial expresses
the reaction of the boundary against the finite individual'' \ed{\emph{SKS} 7, p. 417}, it
is ``precisely the reaction against the absolute relation's absolute expression''
\kpage{142}\ed{\emph{SKS} 7, p. 417}, ``the\footnote[73]{``The suffering in the inward, \emph{the suffering of the God-relationship}.'' \ed{\emph{SKS} 7, p. 418.}}
absolute's own resistance.'' \ed{\emph{SKS} 7, p. 417.})

% TR: Afdøen fra Umiddelbarheden = dying away from immediacy
% TR: Tilintetgjørelse / Selvtilintetgjørelse = annihilation / self-annihilation
The religious suffering, which continually persists, is the individual's \emph{dying away
from immediacy}. The meaning of suffering is: to express, in existing:
that \emph{the individual is capable of nothing at all, but is nothing}\footnote[74]{``The human being is capable of nothing at all before God \emph{except becoming conscious of this}.'' \ed{\emph{SKS} 7, p. 419.}}
\emph{before God}; for here again
the God-relationship is recognizable by the negative, and \emph{self-annihilation}
is the essential form of the God-relationship.
% TR: Formaaen / Ikke-Formaaen = being-able / not-being-able
% TR: Religiøsitet = religiousness (Religiøsiteten A / B = Religiousness A / B)
And this not-being-able
the individual is to have continually with him, for its disappearance is the disappearance
of religiousness. (This does not mean that the human being is to will
to undertake nothing at all because everything is after all vanity and
trumpery; but to work with one's utmost ability and then nevertheless to understand that all our
being-able and not-being-able is a jest.) The religious person lies bound
% Forestilling here, and Gudsforestilling(en) throughout this span = idea (the idea of God), not Hegel's Vorstellung
in finitude with the absolute idea of God present to him in
the lowliness of the human being. The suffering of annihilation is there for him when
in his nothingness he has the absolute idea, but no reciprocity.
He wills absolutely to hold fast to the idea, and precisely this annihilates
him; he wills to do everything, and while he is doing it, the powerlessness
begins; for \ed{71.5} for a finite being there is, after all, a meanwhile; he wills
to do everything, he wills to express this relation absolutely, but he cannot
make finitude commensurable with it. --- And yet the individual is to
exist in finitude, and cannot there uninterruptedly lead the life of
eternity. When accidental finitude is put together with the idea of God,
the difficulty becomes, after that first difficulty
has been overcome: \emph{to be able to do it with God}. (The \emph{ethical} steps in
as regulative.) (The \emph{equality} of love expresses itself in that between God and
% Forskjel / Forskjellighed rendered difference / unlikeness; the construal of "ved" here is doubtful
human being there is \emph{the absolute difference} through the \emph{absolute unlikeness}, and
its form is the \emph{humility} that fully acknowledges its human
lowliness with humble confidence before God.) ---\emph{The actuality of suffering}
% TR: Virkelighed = actuality
% TR: Bestemmelse = determination
belongs to inwardness, and must not express itself in the outward; the
% forklarende: "explanatory" (possibly "transfiguring"); kept literal
religious is \emph{the inwardness whose explanatory determination is to be hidden}
(the determination of the ethical was disclosure). The martyrdom of hidden
inwardness lies in the suffering of annihilation that is dying away
from immediacy, in the very resistance of the divine
\kpage{143}against an existing individual who is prevented from relating himself absolutely, and finally
in life in the world with this inwardness without having an expression
for it. ---

Suffering was determined as the specifically religious expression of existential pathos;
but it finds a still more decisive expression
in \emph{guilt}. In existence the individual is a concretion, time concrete,
and even while the individual deliberates, it is ethically responsible for the use
of time. Existence is not abstract haste, but \emph{striving} and \emph{a
continuing meanwhile}; even in the moment the task is posited, something
has already been wasted, because existing has gone on meanwhile and the beginning has not
been made at once.

% TR: evig Salighed = eternal happiness (salig = blessed)
% "skal gjøres til Afdøen": construed as the beginning to be made on (toward) dying away
\ed{71.6} Just as the beginning is to be made on dying away from immediacy,
it is discovered that, since time has meanwhile passed,
\emph{a wrong beginning} has been made, and that the beginning must be made by becoming
guilty; and \emph{from that moment the total guilt, which is decisive, accrues usurious interest in new
guilt. Guilt is the most concrete expression of existence; and precisely for that reason it is, for the existing individual,
the expression of his relation to the eternal happiness}. The more abstract the individual
is, the less it relates itself to an eternal happiness, and the
more it also distances itself from guilt; for abstraction puts existence
into indifference, but guilt is the expression of existence's strongest
self-assertion, and it is, after all, an existing individual who is to relate himself to an
% Koch prints "Existentsens selv" (for "Existentsen selv"?); rendered as the sense
eternal happiness. (\emph{Yet the individual cannot shift the guilt onto the one who has placed
him in existence, or onto existence itself}, cf. Concluding Unscientific
Postscript p. 405f. \ed{\emph{SKS} 7, 477f.}). \emph{The totality of guilt} comes into being
for the individual by putting his guilt, even if it were only a single one,
even if it were the most insignificant of all, together with \emph{the relation to an
eternal happiness}. The putting-together yields the determination of quality. Since
the idea of God is included, it transforms the determination of guilt into a
determination of quality. Put together with the comparative as
standard, guilt becomes a quantifying; over against the absolute
quality, guilt becomes dialectical as quality.

% TR: Misforhold = disproportion (relation to eternal happiness; not the self-relation of The Sickness unto Death)
The \emph{essential consciousness of guilt} is the expression for the relation to an
eternal happiness by expressing \emph{the disproportion}. Yet, however decisive the consciousness
may be, \emph{it is nevertheless always the relation that bears the disproportion}, only
the existing individual cannot get hold of the relation, because the disproportion
continually places itself in between as the expression for the relation. But
% TR: Brud(det) = (the) break
on the other \ed{71.7} hand they do not, after all, repel each other so
\kpage{144}(the eternal happiness and the existing individual) that \emph{the break} constitutes itself
as such; on the contrary, only by being held together does the disproportion repeat
itself as the decisive consciousness of guilt essentially,
% TR: Immanents = immanence
not of this or that guilt. \emph{The consciousness of guilt still lies within
immanence}. It is \emph{the same} subject that, by holding guilt together
with the relation to an eternal happiness, becomes essentially guilty,
but the subject's ideality is nevertheless such that guilt does not make the subject
into \emph{another}, which is the expression for the break. But the break,
in which lies the \emph{paradoxical} accentuation of existence, cannot enter
into the relation between an existing individual and the eternal, \emph{because the eternal everywhere
embraces the existing individual}, and the disproportion therefore remains \emph{within immanence}.
If the break is to constitute itself, the eternal must \emph{itself} determine
itself as something \emph{temporal}, as in time, as \emph{historical}, whereby the existing individual
and the eternal in time get eternity between them. This is
\emph{the paradox}.

% TR: Erindren = recollecting; Vedvaren = persistence
The decisive expression of the consciousness of guilt is this consciousness's
\emph{essential persistence}, or \emph{the eternal recollecting of guilt} because it is continually
put together with the relation to an eternal happiness. The consciousness is
that the subject \emph{is decisively changed}, while the subject's \emph{identity} nevertheless lies
in this, \emph{that it is he himself who becomes conscious of it} by putting the guilt together
with the relation to an eternal happiness. Guilt is the expression of \emph{freedom}.

The religiousness treated in the foregoing, \emph{the religiousness of hidden
inwardness}, is \emph{not} the specifically \emph{Christian} one. The kind of religiousness
described (\emph{Religiousness A.}) is not undialectical, but it is \emph{not
paradoxically dialectical}. It is the dialectic of inward deepening, it is the relation
to an eternal happiness not conditioned\ed{71.8} by a something, but is
the dialectical inward deepening of the relation, thus conditioned only by the inward deepening,
which is dialectical. \emph{This religiousness, which can also be present in
paganism, must be in the individual before there can be any question of becoming attentive
to the dialectical determinations of the Christian. It has only human nature in
general as its presupposition}.

The upbuilding in this sphere is that of immanence, is \emph{the annihilation
in which the individual puts himself aside in order to find God, since it
is namely the individual himself who is the hindrance}. The upbuilding is thus here quite
rightly recognizable by \emph{the negative}, by the self-annihilation that finds
the God-relationship \emph{within itself}, that, suffering through, sinks into the God-relationship,
% "Gud er i Grunden": also "God is, at bottom"; the play on "grundes" is lost
\kpage{145}is grounded in it, because God is in the ground, if only everything that is in
the way has been cleared away, every finitude and first and foremost the individual
himself in his finitude, in his insistence on being in the right against God.

% TR: Confinium = border territory
The border territory of this religiousness is \emph{humor}. The humorist grasps the meaning
of suffering in relation to existing, \ins{he feels the pain
and is sympathetic}, but then makes the deceitful turn and revokes
the suffering in the form of jest. He continually puts the idea of God
together with something else, and brings out the contradiction,
but he does not himself relate himself in religious passion to God; he
transforms himself into a jesting and yet profound point of transit
for all these conversions, but he does not himself relate himself
to God. Humor puts the eternal recollecting of guilt together with
everything, but does not itself, in this recollecting, relate itself to an eternal happiness.
Humor rests in immanence; it is, when it makes use of the
Christian determinations, not Christianity but a pagan speculation
that has come to \emph{know} everything Christian. It is \emph{the equilibrium of
the comic and the tragic}.

\ed{72.1} In the same relation as humor stands to religiousness,
\emph{irony} stands to the ethical. Irony arises by continually
putting the particulars of finitude together with \emph{the ethical infinite
requirement} and letting the contradiction come into being. But the ironist does
not in himself relate himself to the absolute requirement of the ethical. Irony is an \emph{art of existence},
he who essentially has it has it \emph{always}; for irony is
the infinity in him. --- ``The ironist levels on the basis of \emph{abstract humanity},
the humorist on the basis of \emph{the abstract God-relationship}.'' \ed{\emph{SKS} 7, p. 408.}
--- ``In irony there is \emph{no sympathy}, there is self-assertion, and its sympathy
is therefore altogether indirectly sympathizing, not with any
human being, but with the idea of self-assertion as every human being's
possibility.'' \ed{\emph{SKS} 7, p. 502.}

Since humor and irony become the border territories of the religious and the
ethical, they can become the \emph{incognito} of the religious person and the ethicist.

% Koch: an inserted one-page summary leaf, apparently belonging to lecture 25
\ednote{}

\emph{The religious}. 1) the absolute orientation toward the absolute τελοσ,
expressed in action in the transformation of existence. (aesthetic --- existential
\kpage{146}pathos?). --- Relinquishing the moments of finitude over against
the eternal happiness. --- The abstraction of the eternal happiness. ---
\emph{The expectation} of an eternal happiness. --- The absolute τελοσ in relation to
the relative ones (NB. God and the absolute τελοσ). --- At one and the same time to
relate oneself absolutely to the absolute τελοσ and relatively to the relative ones.
--- The absolute distinction. --- Worship as the God-relationship.
--- To \emph{become} infinite. --- Inward infinitizing. --- Determination
of the essence of the eternal happiness.

2) The practice of the relation to the absolute τελοσ. Suffering --- the acting of
inwardness is suffering. --- The relation of suffering to the God-relationship.
--- The actuality of suffering is its \emph{persistence}. (NB. Spiritual trial)
--- Dying away from immediacy. That the individual is capable of nothing at all,
but is nothing before God (self-annihilation). --- The norm's determination
of suffering. --- Existence in finitude. --- Religiousness
as hidden inwardness. --- The martyrdom of hidden inwardness.

3) Guilt as the specific expression of the religious. --- Guilt as
the most concrete expression of existence. (NB. pushing the guilt away from oneself).
--- The totality of guilt. --- The consciousness of guilt; its relation to immanence.
--- The eternal recollecting of guilt.

The relation between Religiousness A. and the Christian. Humor ---
irony.

\lecture{\ed{72.1 continued} Lecture 26}

% TR: Tilegnelse = appropriation
f) \emph{Christianity}. --- From the preceding religious sphere the
Christian, \emph{paradoxical religiousness} (Religiousness B), is distinguished by this,
that it has \emph{conditions} that are not the dialectical deepening of inward deepening,
but a \emph{specific something that more precisely determines the eternal happiness},
not in that it more precisely determines the individual's appropriation of it, but
more precisely determines the eternal happiness, yet not as a task for
thinking, but precisely as \emph{paradoxically repelling toward new pathos, the pathos of
% "lærende" (learning, or teaching) as printed; not emended to "troende"
the absurd}. The learning Christian both has and uses his understanding, respects
the universally human, does not explain it by lack
of understanding if someone does not become a Christian, but in relation to
Christianity \emph{he believes against the understanding}, and here too uses the understanding
--- in order to take \ed{72.2} care that he believes against the understanding. \emph{Nonsense
\kpage{147}he therefore cannot believe against the understanding}, for\footnote[75]{NB. \emph{Religiousness A must come in between as terminus a quo for paradoxical religiousness}, only
thereby is it prevented that the Christian determinations run together with the aesthetic ones.}
the understanding will precisely see through
that it is nonsense, and prevent him from believing it, but he
uses the understanding so much that through it he becomes \emph{attentive to
the incomprehensible}, and now he relates himself to this, believing against the understanding.
\emph{Faith belongs essentially in the sphere of the paradoxical-religious}; all other
faith is only an analogy that is none. In faith there is a break with the understanding
% Koch prints "Imannentsens" (for "Immanentsens"); rendered as the sense
and with immanence; immanence's last foothold, \emph{eternity
behind}, is lost, and the subject exists hard pressed at the extremity
of existence, \emph{by virtue of the absurd}.

That in which the Christian can be concentrated is: \emph{that an eternal happiness
is decided in time by the relation to something historical}. All the other Christian
determinations lie in this one, and can be consistently derived from
it. It forms the sharpest contrast to paganism; in
it lies \emph{the new dialectical element that sharpens pathos}.

The dialectical contradiction is then \emph{first} this: that \emph{an eternal happiness is expected
in time through a relation to something else in time}. Existence is here paradoxically
accentuated, the distinction \emph{here and hereafter} absolutely determined \emph{in that the eternal
itself has come into being in a moment of time}. Since the eternal has come into being in a moment of time,
% TR: Guden = the god (as against Gud = God)
the existing individual does not come in time to relate himself
to the eternal or to bethink himself of his relation to it (this is
A.), but \emph{in time to relate himself to the eternal in time}, so that the relation
is thus within time, which relation is equally in conflict with all thinking,
whether one reflects on the individual or on the god.

% TR: Modsætning = opposition (Modsigelse = contradiction)
% "forledt af" (led astray by) as printed; "forladt af" (abandoned by) would give easier sense
\emph{The paradoxical-religious posits the opposition absolutely between existence and the
eternal}, for precisely the fact that the eternal is at a specific \ed{72.3} moment of time is
the expression for this, \emph{that existence is led astray by the hidden immanence of the eternal}. In the
religious A. the eternal is ubique et nusquam, but hidden by the actuality
of existence; in the paradoxical-religious the eternal is at a specific
place, and this is precisely \emph{the break with immanence. The individual is paradoxically dialectical},
every remnant of original immanence annihilated, and all connection
cut off, the individual hard pressed at the extremity of existence. Existing
becomes the absolute contradiction of immanence. There is
no underlying immanent kinship between the temporal
and the eternal, because the eternal itself has come into time and wills \emph{there} to constitute
\kpage{148}the kinship. --- \emph{The teacher} becomes \emph{the god in time}, the disciple a \emph{new
creation}, the individual, who was not eternal, \emph{becomes so}, thus does not bethink itself
of what it is, but becomes what it was not, and, be it noted,
becomes something that has the dialectic that, as soon as it is, it \emph{must have
been}; for this is the dialectic of the eternal. What is inaccessible to all thinking
is \emph{that one can become eternal although one was not}.

% TR: Kundskab = information (Viden = knowledge)
The dialectical contradiction is \emph{next} this: \emph{that an eternal happiness is based
on the relation to something historical}. For all historical knowledge and information
it holds that it is only an \emph{approximation} even at its maximum.
The contradiction is: to build one's eternal happiness on an approximation,
which can be done only \emph{when one has no eternal determination in oneself}
(which again cannot be thought, any more than how one
then hits upon it; \emph{therefore the god must provide the condition as well}), for which reason this
again hangs together with the paradoxical accentuation of existence.
The contradiction arises in that the subject, at the extremity of its subjective passion
\ed{72.4} (in concern for an eternal happiness), is to base
this on a piece of historical information whose maximum becomes an
approximation. This presupposes that the existing individual must have
lost continuity with himself, must have become \emph{another} (not different
from himself \emph{within himself}), and now, by receiving the condition from the god,
become a new creation. The contradiction is \emph{that becoming a Christian
begins with the miracle of creation}, and that this befalls one who is
created, and that Christianity is nevertheless proclaimed to all human beings,
% "tilværende" = existing (there, in being), not "existerende"
who by it must be regarded \emph{as not existing}, since the miracle by which they
come into being must come in between, as actual, or as the expression for
the break with immanence and the resistance that absolutely makes the passion of faith
paradoxical so long as one exists in faith, i.e., for\footnote[76]{``Christianity makes existing paradoxical, for which reason it becomes the paradox so
long as one exists, and only eternity has the explanation.'' \ed{\emph{SKS} 7, p. 512.}}
the whole of life; for he always has, after all, his eternal happiness based
on something historical. Here too is a break with all thinking.

The dialectical contradiction is \emph{finally} this: \emph{that the historical in
question is not something simply historical, but formed from what can become historical only against its
essence, thus by virtue of the absurd}. The historical is that
the god, the eternal one, has come into being at a specific moment of time as a
single human being. That the one who according to his essence is eternal comes into being in time,
% Koch prints "efter sit i Væsen Evige" (garbled); rendered as the sense "efter sit Væsen Evige"
\kpage{149}is born, grows, dies, is a break with all thinking. In all \emph{speculative} talk
% TR: Forstaaelse = comprehension (Forstand = understanding)
of an eternal becoming of God the scene is in the medium of \emph{possibility}\footnote[77]{``Faith is the objective uncertainty with the repulsion of the absurd, held fast in the passion
of inwardness, which is precisely the relation of inwardness raised to its highest power.''
\ed{\emph{SKS} 7, p. 554f.} ``The appropriation [the believer's] is the paradoxical inwardness.'' \ed{\emph{SKS} 7,
p. 554.}}, but when
it is moved over into that of \emph{actuality}, \emph{the absolute paradox} emerges. Here
possibility is not higher than actuality, but actuality
is the highest and is the absolute paradox. Actuality, i.e., that
this and that has actually happened, is the object of faith, and yet surely not
any human being's or humanity's own thoughts; for thought is
then at most possibility, but possibility as comprehension is
precisely the \ed{72.5} comprehension by which the retrograde step is taken
that faith ceases. In the fantasy-medium of possibility God can very
well, in imagination, fuse together with the human being, but
in actuality, with the single human being, lies precisely the paradox.

In Religiousness A, which rests in immanence, \emph{every} human being
is\footnote[78]{``Religiousness A makes existing as strenuous as possible (outside the
sphere of the paradoxical-religious), but does not base the relation to an eternal happiness on
its existing, but lets the relation to an eternal happiness be the basis of the transformation of existence.''
\ed{\emph{SKS} 7, p. 522.}}
essentially regarded as a partaker in happiness and becomes so in the end. \emph{Religiousness
B is separating, polemical}: only \emph{on this condition} do I become blessed,
and as I absolutely bind myself by it, so I exclude
everyone else. Having one's eternal happiness based on something historical
makes the Christian's good fortune recognizable by \emph{suffering}; the chosen one
is not directly understood as the fortunate one, but understood with such
difficulty that for anyone other than the chosen one it would have to be something
to despair over.

% TR: Syndsbevidsthed / Skyldbevidsthed = consciousness of sin / of guilt
Through the dialectical contradictions in Religiousness B \emph{pathos is sharpened}, and
this sharpened pathos is a) \emph{the consciousness of sin}. This is the expression
for \emph{the paradoxical transformation of existence}. Sin is the new medium of existence;
to exist now means that \emph{the individual, by having come into being, has become a
sinner}. The consciousness of sin is the break with immanence; by
coming into being the individual becomes another, or: in the now in which he is to come into being,
he becomes, by coming into being, another. From eternity the individual is not
a sinner, but when the being eternally disposed, which comes into being in birth,
\kpage{150}becomes a sinner in birth or is born as a sinner, then it is existence
% Koch prints "ethver" (for "ethvert"); rendered as the sense
% TR: Tilblivelse = coming into existence
that so closes in around it that every communication of immanence
by way of recollection, through being taken back into the
eternal, is broken off, and the predicate ``sinner,'' which first, but also at once,
appears with the coming into existence, acquires such paradoxically overwhelming
\ed{72.6} power that the coming into existence makes him another. This is the consequence
of the god's appearance in time, which prevents the
individual from relating himself backward to the eternal, since he now, going forward,
comes in time to become eternal through the relation to the god in time.
\emph{The consciousness of sin, therefore, the individual cannot, like the consciousness of guilt, come to know
by himself}; for the consciousness of sin is the very change of the subject,
which shows that there must be \emph{outside} the individual the power that informs
him that he, by having come into being, has become other than
he was, has become a sinner. This power is \emph{the god in time}. --- In
the consciousness of sin the individual becomes conscious of himself \emph{in his unlikeness to
the universally human}, which only by itself becomes conscious of what
it is to exist as a human being. For since the relation to that historical something
conditions the consciousness of sin, \emph{this cannot have existed during all the time
in which this historical something did not exist}. The believer extends the consciousness of sin
% TR: Arvesynd = hereditary sin
to \emph{the whole race},\footnote[79]{NB. \emph{Hereditary sin}. The Sickness unto Death. p. \emph{90. 97}. \ed{\emph{SKS} 11, p. 202f., p. 208.}}
but does not at the same time know the whole race saved,
inasmuch as the single individual's salvation will, after all, depend on his being brought into
relation to that historical something, which precisely because it is historical cannot
be everywhere at once, but takes time to become known to human beings,
so that countless generations, without any guilt of their own, may remain
ignorant that it has existed. To have an existence in this
determination is sharpened pathos, both \emph{because it cannot be thought} and
\emph{because it is isolating}. For sin is no teaching or doctrine for
thought, in that case the whole thing becomes\footnote[80]{\emph{The forgiveness of sins} is the paradoxical satisfaction by virtue of the absurd.}
nothing; it is a \emph{determination of existence}
and precisely cannot be thought.

b) \emph{The possibility of offense} or the \emph{antipathetic collision}, which is immediately
given with the paradox, which requires \emph{faith against the understanding}. Christianity
is the only power that can in truth arouse offense.
For the believer offense is at the beginning, and its
possibility is the \ed{72.7} constant fear and trembling in his existence.

\kpage{151}c) \emph{The pain of sympathy} on account of the believer's not, like the
religious person A., latently sympathizing, and not being able to sympathize, with \emph{every
human being qua human being}, but essentially \emph{only with the Christians}. For
the believer it holds that outside that condition there is no
happiness, and for him it holds, or it may come to hold for
him, that he must hate father and mother. For is it not like
hating them, to have one's happiness bound to a condition that he knows they
do not accept? He can to the very last be willing to do everything for them, fulfill
all a faithful son's duties with the greatest enthusiasm; in that way
Christianity does not bid him hate, and yet when this condition separates
them, separates them forever: is it not as if he hated
them? And is this not a terrible sharpening of pathos in relation
to an eternal happiness?

The specific determination of the Christian is thus\footnote[81]{(\emph{with regard to the section 72.7--73.4, cf. the enclosed leaf.})
% Koch: cross-reference to the leaf printed after the end of lecture 26
\ednote{}}
\emph{sin} and
its opposite \emph{faith}, and the way of faith is through \ins{\emph{despair} and
through} \emph{the possibility of offense}; its object is the \emph{absolute paradox},
which \emph{is believed by virtue of the absurd}.

% TR: Spændkraft = tensile force
With \emph{the doctrine of sin} Christianity begins, and \emph{the consciousness of sin
becomes the tensile force}\footnote[82]{cf. Practice in Christianity \emph{p. 76}. 166. \ed{\emph{SKS} 12, p. 79f., p. 157f.}}
\emph{in the individual that bears faith in the paradox and Christianity's
ideal requirement}. The category of sin is the category of singularity;
% TR: Enkelthed = singularity (Enkeltheder = particulars)
it cannot be thought speculatively, just as one cannot
think the single human being, since singularity lies beneath the concept,
but only the concept human being. But before God, in religion, the
human being is always the \emph{single} human being, alone over against God, alone in
this enormous exertion and this enormous responsibility.

\emph{The definition of sin} is: after having been enlightened by a revelation from God
about what sin is, before God --- or with the idea of God
--- in despair not to will to be oneself, \ed{72.8} or in despair to will to
be oneself. Sin is\footnote[83]{Sin is a position, not a negation cf. The Sickness unto Death. p. \emph{97}. \ed{\emph{SKS} 11, p. 208.}}
\emph{despair}, but the opposite of sin is
not that despair is not and has not been, but that the possibility
of being able to be in despair\footnote[84]{NB. Sin and grace.}
is annihilated and in every moment is being annihilated.
\kpage{152}The definition of the state in which despair no longer
is, is this: in relating itself to itself, and in willing to be itself,
the self is grounded transparently in the power that posited it. But this
formula is the definition of\footnote[85]{This definition of faith is compared with the one given earlier, and it is shown
how they lead to the same point.}
\emph{faith}. In this formula, that is, \emph{faith
in the forgiveness of sins} is contained; for it is the condition for the individual's
being able to relate himself to himself and to will to be himself before God. But in
the forgiveness of sins there lies in turn the faith that God, also for the sake of this
self, let himself be born, became man, suffered and died, the faith
\emph{that Christ was God}. But in faith in the forgiveness of sins and in the faith
that this individual human being is God lies \emph{the possibility of offense}; for \emph{between
human being and God the yawning abyss of quality is fixed by sin}; therefore
% Koch prints "Foragelsen" (for "Forargelsen"); rendered as the sense
offense is there when a human being would forgive sins, and when
the qualities of God and of the human being are put together in this human being,
in \emph{unrecognizability}, and this lowly human being says that he is God,
then the possibility of offense is there again. The words: \emph{blessed is he who takes no offense
at me}, belong to the proclamation about Christ. But he who
does not take offense, he \emph{worships}, believing.

% TR: Samtidighed = contemporaneity (samtidig = contemporary)
% TR: Fornedrelse = abasement (den Fornedrede = the abased one)
% TR: Virkelighed = actuality
% TR: Salighed = happiness (evig Salighed = eternal happiness; the adjective salig = blessed)
The condition for believing in Christ is to\footnote[86]{``In relation to \emph{the absolute} there is only one time: the \emph{present}. For him who is not
contemporary with the absolute, it does not exist at all. And since Christ is the absolute,
it is easily seen that in relation to him there is only one situation, that of contemporaneity'' \ed{\emph{SKS} 12,
p. 75.} --- ``Only the contemporary is actuality \emph{for me}. And thus every human being
can only become contemporary: with the time in which he lives, --- and then with one thing more, with
Christ's life on earth, for Christ's life on earth, the sacred history, stands here by
itself, outside history.'' \ed{\emph{SKS} 12, p. 76.}}
become \emph{contemporary} with him,
to see him in his true form and in the surroundings of actuality,
as he walked on earth, in \emph{the form of abasement}, for in glory no one has
yet seen him, and no one can see him before his return in
glory. The believer \emph{believes} that \ed{73.1} Christ is now in glory,
but it was under the conditions of abasement\footnote[87]{``\emph{Christ's life on earth is nothing past}, whence in each generation true Christians
% doubtful construal: the Danish ("hvad der i hver Generation af sande Christne, ere samtidige") does not construe as printed; rendered by the evident sense
are contemporary with Christ, and have nothing to do with the Christians of the previous
generation, but only with the contemporary Christ. His life on earth follows the race,
and follows each generation in particular as \emph{the eternal history}; his life on earth has eternal
contemporaneity.'' \ed{\emph{SKS} 12, p. 76.}}
that he spoke of his return
in glory; therefore this can only be \emph{believed}, but one cannot
\kpage{153}have become a believer without having come to him in his state of
abasement; otherwise he does not exist, for \emph{only thus has he existed}.
That he will come in glory is \emph{expected}, but it can only be expected
and believed by the one who has held and holds to him as he has
existed. What Christ has said and taught, every word he has spoken,
becomes eo ipso untrue when we give it the appearance of its being Christ
in glory who says it. No, \emph{he} is silent; the \emph{abased one} speaks.

From history one can come to \emph{know} nothing about Christ; for about
him one can \emph{know} nothing at all; he is \emph{the paradox}, the object
of faith, \emph{there only for faith}. One cannot prove his divinity from the \emph{consequences} of his life;
this, which is incommensurable with all historical
consequences, is reached only by faith.

Christ is thus the object only of faith, and the Christ who
is the object of faith is \emph{Christ in the state of abasement}, and from him \emph{the possibility
of offense} is inseparable. Offense relates essentially
to \emph{the composition of God and human being}, or to \emph{the God-man}. The God-man
is the unity of God and \emph{an individual} human being. The relation to
the God-man is impossible except as beginning with the situation of
contemporaneity.

% forraade: "betray" in both senses, give away and reveal, as in English
Offense has 2 forms: in the direction of \emph{loftiness}; one takes offense
at an individual human being's saying that he is God, acting and speaking
in a way that betrays God; --- or in the direction of \emph{lowliness}: that
he who is God is this lowly human being, suffering as a lowly
human being. The God-man is the paradox, absolutely the paradox;
therefore it is quite certain that the understanding must come to \ed{73.2}
a standstill at it. --- (To the latter kind of offense, offense in
the direction of lowliness, there then corresponds a possibility of offense: \emph{that the
Christian must come to suffer in the world}, that truly to be a Christian,
since it is to be \emph{Christ's follower}, becomes \emph{likeness to him in suffering}. ---
% doubtful construal: "Frivillighedens og Forargelsens Mulighed" may also be read "the possibility of voluntariness and of offense"
The decisive thing in \emph{Christian suffering} is the \emph{voluntariness} and the \emph{possibility of
offense} for the sufferer. For the natural human being the \emph{infinite}
passion with which Christianity apprehends eternal\footnote[88]{NB. Christian compassion --- the Christian idea of the misery of the
% Forestilling here = idea
human being (cf. Practice in Christianity I.)}
happiness is an offense. The \emph{absolute} is an offense to the understanding.)

% TR: Ophævelse = annulment
\kpage{154}But the possibility of offense is \emph{the inescapable}. The true God
\emph{cannot} \ins{as \emph{spirit}} become \emph{directly recognizable}; therefore he can only become
the object of \emph{faith}: only\footnote[89]{Practice in Christianity. p. 145. 147. \ed{\emph{SKS} 12, p. 139f., 141.}}
as the object of faith is he \emph{true} God, and
only as true God can he \emph{save the human being}. By his love
he must therefore plunge the human being into the decision of offense. The annulment
of offense was, now, \emph{faith};\footnote[90]{NB. The self's choice of Christ (Practice in Christianity. p. 170ff. \ed{SKS 12, p.
161ff.})}
but faith is a \emph{choice}, by no means
the direct reception, --- and it is the receiver who becomes
manifest, whether he will believe or take offense. If one abolishes the possibility
of offense or faith, one abolishes the God-man as well; but
if one abolishes the God-man, one abolishes Christianity as well.

% TR: Efterfølgelse = imitation (Efterfølger = follower)
% TR: Forbilledet = the prototype
The Christian's true relation to Christ was the relation of
contemporaneity; his task is \emph{to become Christ's}\footnote[91]{Cf. Practice in Christianity. p. 116, cf. p. 126. \ed{SKS 12, p. 114f., p. 123.}}
\emph{follower}, but Christ's follower
\emph{in the state of abasement};\footnote[92]{Christ as ``the prototype'' \ed{cf. e.g. SKS 8, p. 332 and SKS 10, p. 133.}}
for only in this can he become
the object of imitation. But Christ came into the world \emph{in order to suffer};
the imitation of him then becomes, since it is the transformation of our life into
likeness with him, the \emph{voluntary}\footnote[93]{cf. The Moment VI, 18f. VIII, 12. \ed{SKS 13, p. 268f., p. 351f.}}
taking on of suffering. \emph{To will truth
absolutely and to give the truth actuality in one's life} is the Christian's task,
but it belongs to truth just as essentially to suffer in \emph{this}
world as to triumph in another, in the world of \emph{truth}. \ed{73.3} Likeness
to Christ \emph{in love}, which as \emph{love of the neighbor} embraces \emph{all},
it being the love of \emph{self-denial} in contrast to the egoism of preferential
love; likeness to him in \emph{unconditionally willing the true}; likeness
% TR: Afdøethed = deadness (to the world); Afdøen = dying to (the world)
in \emph{deadness} to the world, in the giving up of the earthly, must bring
\emph{Christian suffering}: not only the \emph{inward} suffering, which is given in the very annihilation
of the selfish, of the ego's\footnote[94]{cf. The Moment V, 3f. 11. ``\emph{to love God} in hatred of humankind, \emph{in hatred} of oneself and thereby
of other human beings.'' \ed{\emph{SKS} 13, p. 227f., p. 233f.}--- cf. The Moment VI, 14 \ed{\emph{SKS} 13, p.
266f.}, The Moment VII. 31f. 36. 38. \ed{\emph{SKS} 13, p. 303f., p. 307f., p. 309f.}}
immediacy, in the possibility of despair
and of offense, which must be overcome every moment, in
the uncertainty that in temporality constantly follows the acquired
\kpage{155}certainty, \ins{in the pain of the consciousness of sin, in spiritual trial,} and in the
suffering of sympathy that suffers through the untruth and
sin of others, in compassion for them, but also in suffering \emph{from without}, in hatred,
persecution and mockery, \ins{in suffering for the sake of self-contempt.}\footnote[95]{cf. The Moment VIII, 9. 12f. \ed{\emph{SKS} 13, p. 349f., p. 351ff.} ``What Christ Judges
of Official Christianity'' p. 11f. \emph{13} \ed{\emph{SKS} 13, p. 179f., p. 180f.} ``the Christianity of the New
Testament is: in the fear of God to suffer for the doctrine at the hands of men.'' cf. The Moment IX,
16. \ed{\emph{SKS} 13, p. 383.}}

The last-named follows necessarily from \emph{the relation of Christ and of the Christian to
the world}, and from this relation there follow in turn the consequences that contain
the judgment on \emph{established Christendom}, --- the judgment that was pronounced
on it \emph{through} Kierkegaard's last writings.\footnote[96]{cf. the enclosed leaf
% Koch: the leaf follows the end of lecture 26
\ednote{}}

% TR: Ueensartethed / Uensartethed = heterogeneity (Eensartethed = homogeneity)
% TR: Bestemmelse = determination
Decisive for Christ's relation to the world is the God-man's
\emph{essential heterogeneity} with every other human being and with the whole
race. \ins{(NB. Sinfulness)} Christianity relates itself essentially
to \emph{eternity};\footnote[97]{cf. The Moment \emph{VI, 23, IX, 16}. \ed{\emph{SKS} 13, p. 273f., p. 383.}}
life here on earth is, for each individual in particular
of the countless human beings who have lived, \emph{the time of testing}. It is indeed
Christianity's intention that it is to be proclaimed \emph{to all}, but from this it does
not follow that it should ever happen that all should\footnote[98]{cf. The Moment \emph{V, 16ff.} \ed{\emph{SKS} 13, p. 237.} (NB the defectiveness in the dialectic). ---
Practice in Christianity p. 234f. \ed{\emph{SKS} 12, p. 216 ff.}}
accept
it and become true Christians. Christianity never becomes a
development \emph{within} the determination: the human race; it relates
itself constantly in the same way to what Christianity calls
\emph{the world}: paganism and the natural human being in Christendom.
\emph{The world always remains essentially the same; like the sea, like the air, in short like an
element; for it is, and is meant to be, the element} \ed{73.4} \emph{that can furnish the test for
that which lies in having to be a Christian, who in this world is always a member of the Church
militant}. --- ``The concept Christian is a polemical concept; one can
be a Christian only in opposition or oppositionally!'' \ed{\emph{SKS} 13, p.
187} --- ``To be a Christian is to live, to love God, in a relation of opposition.''
\ed{\emph{SKS} 13, p. 216} --- Christ's saying ``my kingdom is not
of this world'' \ed{John 18:36} is an eternally valid saying, always at every time
equally valid, about the relation between Christ's kingdom and this
world; as soon as Christ's kingdom makes peace with this world and
\kpage{156}becomes a kingdom of this world, Christianity is abolished. When,
on the other hand, Christianity is in truth, it is indeed a kingdom \emph{in} this
world, but not \emph{of} this world, that is: it is \emph{militant}. When the
Christian Church claims to have triumphed in this world, then it is not the Church that
has triumphed; then the world has triumphed. Then the heterogeneity between
Christianity and the world has vanished, the world has won, Christianity
lost. Then Christ is no longer \emph{the God-man}, but an \emph{eminent
human being}, whose life is homogeneous with the development of the race.
Then \emph{eternity is abolished} and the scene of the consummation of all things
placed \emph{within temporality}.

% Koch: describes the loose inserted leaf, an outline of 72.7-73.4
\ednote{}

\ed{1} (Kierkegaard)

\emph{The Christian}. I.) The relation to Religiousness A. --- The paradox --- The understanding's
relation to it. The fundamental determination of the Christian: \emph{that an eternal
happiness is decided in time by the relation to something historical}. --- The \emph{dialectical}
contradictions in this. --- The intensification of \emph{pathos} lying in this (the consciousness of sin.
--- The possibility of offense. --- The pain of sympathy.)

II.) The coherence of the fundamental Christian determinations. --- Christianity's
relation to the world:

a) \emph{The consciousness of sin} (the \emph{single individual}). --- The definition of sin (the relation to
despair). --- The opposite of sin is \emph{faith}; definition of this.
(Faith's relation to God and human being with respect to
its origin.).

b) \emph{Faith} as faith in the forgiveness of sins, which, since the forgiveness of sins is only
possible through God, becomes faith in \emph{the God-man}.

c) The composition: \emph{the God-man} --- Closer determination of this.

d) The inescapability of \emph{offense}. --- Its forms.

e) The relation of \emph{contemporaneity} to Christ (only contemporaneity gives the relation.).
--- The Christ to whom the believer relates himself is
Christ \emph{in abasement} ( --- further development).

\ed{2} compare 72.8, 73.1.).

\kpage{157}f) \emph{The task for the believing Christian}: the \emph{imitation} of Christ (Christ is \emph{the prototype.)}
i.e. imitation of the abased Christ --- of Christ's
suffering.

g*) The relation of the Christian to \emph{the world}. --- The world seen under the determination:
sin. The necessary relation of opposition. --- The heterogeneity
between Christ and the world. The impossibility of Christianity's
fusing with the worldly. --- \emph{Christian suffering}.

h) The consequence of this for the judgment on \emph{the established order}.

*) Since the world is seen under the determination \emph{sin} and the single individual as
sinner (\emph{hereditary sinner}), the relation of opposition becomes present;
each individual who in faith relates himself to the Christian separates
himself from the mass of sinfulness (massa perditionis), but this mass
continues to form the opposition. The possibility that \emph{all individuals}, cf.
The Moment \emph{5, 16ff.} \ed{\emph{SKS} 13, p. 237ff.}, should become Christians is contradicted
by Christianity's conception of time as the time of testing
% Koch prints "SKS 12" for The Moment no. 3, which is in SKS 13 (cf. note 109: The Moment III,8 = SKS 13, p. 191); kept as printed
(NB. Christian self-denial, The Moment 3, p. 8. \ed{\emph{SKS} 12,
p. 191f.}) --- Christianity's \emph{naturalization} is therefore impossible (NB.
becoming a Christian as a child). This follows also from \emph{the essential heterogeneity
of Christ and of Christ's history with the race and the history of the race}. It remains
the eternally contemporary absolute, which is never assimilated.

\lecture{\ed{73.4 continued} Lecture 27.}

Herein, then, lies \emph{the judgment on the established order}, on \emph{Christendom}. It can
be summed up thus: \emph{``Christendom'' is the putrefaction of Christianity;
% TR: Bladartikler = S. Kierkegaards Bladartikler (1857, ed. R. Nielsen), the collected
% newspaper articles, which Brøchner cites by page: left in Danish as the title of that volume
a ``Christian world'' is: the falling away from}\footnote[99]{cf. Bladartikler p. 109 \ed{\emph{SKS} 14, p. 193, col. 2}.}
\emph{Christianity} \ed{\emph{SKS} 14, p. 173, col. 1.}.
The \emph{established order} is the homogeneity between Christianity and the world;
it is the direct commensurability in which being a Christian
gets its corresponding expression in the external, through external recognition,
through esteem, honor, wealth, power, \ins{it is \emph{naturalized} Christianity.}
\emph{The relation of opposition} is thus taken away, \ins{and since it is
taken away, this has happened in such a way \emph{that the Christian has gone back to worldliness's}
world. From this, in the first place, the consequence becomes}
% doubtful construal: base text and \ins{} join awkwardly in the Danish ("Følgen} af den christelige Lidelse er den udenfra kommende er borttaget"); rendered by the sense
that of \emph{Christian suffering} the one coming \emph{from without} is taken away, the one that is given \ed{73.5}
\kpage{158}by the relation of opposition, and thereby the \emph{inner} one too, the one that is
the reflex of the outer in sympathy. \ins{But next, then, the consequence becomes that
Christianity in this fusion has lost \emph{its absoluteness}, and
that the \emph{suffering} that stands in relation to this must also disappear. Becoming
a Christian becomes something one simply is, by being born
into Christianity; the earnestness of decision disappears; \emph{all are Christians}.}
As the absoluteness of Christianity disappears, not only does
the \emph{ideal requirement} disappear: the proclamation of Christianity is toned down,
the specifically Christian fades and is veiled, and this toned-down
Christianity is then set up as the highest; the ideal is declared to be
extravagance and fantasy; but since Christianity only
is \emph{as the absolute}, the fact that this is lost is \emph{the cessation of Christianity}.\footnote[100]{NB. ``The Condition of Christendom.''}

% TR: Sandsebedrag = illusion
It is against this that \emph{Kierkegaard's last writings} direct the decisive
objection. The Christian had been the task of his life and of his authorship,
not change\footnote[101]{cf. The Moment VIII.7 \ed{\emph{SKS} 13, p. 348}.}
in the external, but \emph{the inward deepening
of Christianity}. In all his writings his task had been: \emph{to introduce
the ideals}, to fight \emph{against the illusions by means of the ideals}. Against the illusion
of Christendom, against the concealment of the decisive, the
\emph{ideal} form of Christianity therefore had to be brought forth. This happened
in ``The Sickness unto Death'', still more definitely in ``Practice in Christianity''.
Here the requirement is forced up to the highest of ideality
so that it could be heard, so that one, with the consciousness of it,
might learn ``not only to flee to grace; but to flee to it
\emph{also in relation}\footnote[102]{cf. Bladartikler p. 123. p. 59 note. \ed{\emph{SKS} 14, p. 213, p. 124, the footnote.}}
\emph{to the use of grace}'' \ed{\emph{SKS} 12, p. 15.}, that one ``before
God might sincerely humble oneself under the requirements of ideality.''
\ed{\emph{SKS} 12, p. 79.} The book was \emph{a question} put indirectly to
% construal: "den" (the Christianity "som den bekjender") taken as the book, not the established order
the established order: whether it would acknowledge the Christianity that
the book professes, and in this acknowledgment make the admission concerning its own
Christianity. It could thus at one and the same time become
the only possible defense of the established order, and the judgment on the
established order. \ed{73.6} It became the latter through the established order's own
% Skyld here in its everyday sense, fault; Koch prints "uundgaaeelig", rendered as the sense
fault, and the polemic against the established order now became inescapable.
It is superfluous to discuss the \emph{external history} of this polemic; it is still\kpage{159}
so fresh in memory. I will only point out that it
does not begin at the moment \emph{when the article on Mynster} was published in
``Fædrelandet'' \ed{cf. \emph{SKS} 14, p. 151}, but that, just as it had long
been foreseen by Kierkegaard and hinted at in conversation, it is also hinted at
in his writings. This holds not only of such \emph{more general
hints} as this one: ``There are cases where an established order can
be of such a character that the Christian ought not to put up with it,
ought not to say that Christianity is precisely this indifference to the
external'' (Bladartikler. p. 53 \ed{\emph{SKS} 14, p. 114, col. 2}), --- but of the
\emph{most definite statements}, such as this one: ``I have, after all, never in Christendom heard
any discourse or sermon about which I, if it were before God
that this question was put to me, would unconditionally dare to say that it was
Christian'' \ins{(Practice in Christianity. p. 241 \ed{\emph{SKS} 12, p. 222})
and of the whole portrayal in ``For Self-Examination'' \emph{p. 13ff.} \ed{\emph{SKS} 13, p. 46ff.},
where Mynster is characterized without being named (cf. Bladartikler
p. 69 \ed{\emph{SKS} 14, p. 132, col. 2 -- p. 133, col. 2}) and Kierkegaard's
own position designated.} --- And for nearly every one of the articles in
``The Moment'' one will find the catchwords in Practice in Christianity, especially in its last
part.

% TR: Sandhedsvidne = witness to the truth
In the conflict itself only \emph{the course of development}, leaving aside the
personal, will occupy us. In what form the \emph{first protest} was lodged
had significance at the moment it was lodged; that it already
no longer has now. The passions that were set in motion at the time by
misunderstanding and distortion of Kierkegaard's action have
subsided again; it is now not difficult to gain the calm to see that
the protest was not directed at the individual personality, that it was not directed at
the usage of the word \emph{witness to the truth}, but that it was directed at the matter itself:
% \ed{sic}: Koch marks the double definite "den ideale Christendommens Sandhed", which English cannot show
over against the ideal truth of Christianity \ed{sic}, which
% Koch prints "Indvidet" (for "Individet"); rendered as the sense
is truth only \emph{as life}, as realized by the individual \ins{in the imitation of Christ
in suffering,} to set up something else as truth, something \emph{lower}, and this
only \emph{as}\footnote[103]{``Christianity is no doctrine but expresses an existence-contradiction and is
an existence-communication'' (i.e. a doctrine that is to be realized in existence). Concluding
Unscientific Postscript. p. 289. \ed{\emph{SKS} 7, p. 345f.}}
\emph{an object of contemplation}, and then let the proclamation of this
entitle one to the name witness to the truth. With that conception of truth,
the concept of the witness to the truth \ed{73.7} had to become inseparable from
\kpage{160}\emph{suffering}; for since Christianity is heterogeneous with the world, the
witness to the truth must always be recognizable by heterogeneity with
this world, by offense, by suffering. (cf. what was developed earlier.)
In that protest against the use of the name witness to the truth for the most
eminent representative of established Christendom in Denmark,
the judgment was thus pronounced on this whole Christendom itself,
and the judgment ran, as has already been said, that ``Christendom'' is
the distortion of Christianity. It was pronounced \emph{that official Christianity
was a blasphemy}, that\footnote[104]{cf. Bladartikler. p. 103 \ed{\emph{SKS} 14, p. 180, col. 2 -- p. 181}. The Moment. V, 32 \ed{\emph{SKS}
13, p. 251.}}
% lone slash after "impropriety" as Koch prints it after "Mislighed"; no matching slash
the \emph{pastor's} position was an impropriety /
(cf. Bladartikler p. 111 \ed{\emph{SKS} 14, p. 197, col. 2}, The Moment II,7.
III,7f. V, 30, VII, 12. 40) \ed{\emph{SKS} 13, p. 149f., p. 191f., p. 249f., p. 289, p.
310f.}, that \emph{Christianity had vanished} in the illusion of a\footnote[105]{cf. Bladartikler. p. 85. 106. \emph{130}. \ed{\emph{SKS} 14, p. 155, col. 2, p. 189 col. 2 -- 190, col. 1, p.
220.} The Moment II,9. III,3. IV,16. VI,11 \ed{\emph{SKS} 13, p. 151, p. 187, p. 216, p. 263f.}.}
\emph{Christian world}, in which we all by birth are\footnote[106]{cf. The Moment. 7, 19f. \ed{\emph{SKS} 13, p. 297.}}
Christians, that \emph{indifferentism}
had stepped into\footnote[107]{The Moment. 6, p. 9. \ed{\emph{SKS} 13, p. 262f.}}
the place of religion, that humanity was not even
``\emph{in a state of religion}'' \ed{\emph{SKS} 13, p. 315}, and that the sorrowful question
was raised: is \emph{faith to be found on}\footnote[108]{``Christianity --- does not exist at all'' (Bladartikler. p. 89 \ed{\emph{SKS} 14, p. 163, col. 1}).
``The Christianity of the New Testament does not exist at all'' (ibid. p. 93 \ed{\emph{SKS} 14, p. 169 col. 1}).
cf. The Moment. VI,16 \ed{\emph{SKS} 13, p. 267.}}
\emph{earth}? it was raised by
one who himself did not claim to be a Christian, but only \emph{to know
what Christianity} was, to strive after it, and to will \emph{honesty} (cf.
Bladartikler. p. \emph{100-103}. \ed{\emph{SKS} 14, p. 177-81.})

\ins{In ``What Christ Judges of Official Christianity'' Kierkegaard himself
describes the course of development in the conflict in the following way:}
The beginning was made in \emph{Practice in Christianity} with \emph{the pseudonymous
poetic presentation} of ideal Christianity. The aim was taken at
official Christianity, which abolishes the imitation of Christ, so that
one relates oneself to the prototype \emph{only through the imagination}, oneself
lives in quite other determinations, which is to relate oneself only \emph{poetically}
to Christianity, and finally throws the relation away altogether,
and lets what one is, mediocrity, pass more or less for the ideal.
\kpage{161}Faced with this \emph{poetic} presentation, those concerned were secure,
they noticed nothing; only the article on Mynster startled them,
and now that presentation was made out to be exaggeration etc. Then
the poet transformed \ed{73.8} himself; he brought out the \emph{New Testament}, reminded them of
the pastor's oath on it, and now there was silence. He now spoke, for the time being
\emph{in his own name}, but ever more and more decisively; he began
by putting the matter as favorably as possible to the opposing party, by
pointing to the gentlest reshaping of the established order, but
as this was spurned, the polemic sharpened; the contradiction
% Koch prints "officiele" (for "officielle"); rendered as the sense
in official Christianity, and in the pastor's position, the distortion
of Christianity, was brought forward ever more sharply, and finally
he took it upon himself to say in his own name: \emph{that it is guilt, a great guilt, to
take part in the public worship of God}. Then at last \emph{Christ's own judgment
on official Christianity was brought forward}, and from
this moment the character of the polemic took on the utmost severity; Christianity's opposition
to the world, its requirement of dying to the world and of suffering, was presented
in its whole unconditionality. \ins{Suffering \emph{from without} was stressed ever more
sharply. God's love expressed through suffering. The Moment
VI,14f. VIII,12. \ed{\emph{SKS} 13, p. 265f., p. 351f.}} With the deepest pathos
there mingled annihilating \emph{satire}: ``Laughter, employed in fear and trembling,
became the scourge'' \ed{\emph{SKS} 13, p. 301}, which struck the established order, in which
the relation is now so reversed that Christianity, which came
into the world as the truth one \emph{dies for}, has now become the truth
% "hvoraf man boer" (lit. "of which one dwells"): rendered by the sense of the antithesis dies for / lives off
\emph{off which one lives}, comfortably set up, living by depicting that Christ
was sacrificed, --- struck the self-contradiction that the pastor, who at
ordination relates himself to a kingdom that is not of this world,
at the same time takes an oath to the State, which surely is of this world,
thus the opposite of that kingdom;\footnote[109]{cf. The Moment III,8 \ed{\emph{SKS} 13, p. 191.}}
struck the secularization that came
with this relation to the State and made the proclamation of strict Christianity
by such pastors impossible. The judgment was pronounced:
\emph{the clerical estate is, Christianly,}\footnote[110]{cf. The Moment V,30 \ed{\emph{SKS} 13, p. 249f.}. The Moment VII,40 \ed{\emph{SKS} 13, p. 310f.}.
(compare earlier \emph{milder} utterances).}
\emph{in and for itself of evil}, is a corruption, a
human selfishness that turns Christianity exactly the opposite way round
from the way Christ turned it; it was said: naturally, there is in the clergy
\kpage{162}not one single honest man; \ed{74.1} pastor became synonymous with
perjurer. The Moment VII \ed{\emph{SKS} 13, p. 302}. In the whole deception of
Christendom an immensely comprehensive \emph{criminal history} was sought,
the fundamental confusion, the source of the distortion, is placed \emph{in the concept Church} or \emph{Christendom}.\footnote[111]{cf. Bladartikler, p. 95 \ed{\emph{SKS} 14, p. 173, col. 1}.}
% "søgtes" (was sought) kept literally: Kierkegaard is said to have searched out in Christendom's deception a vast criminal history
Therefore everything was condemned by which this deception was supported, everything
by which the decision of Christianity became insignificant; protest was made
against \emph{becoming a Christian as a child}, against the taking of the vow at
a childish age, against the untrue ``\emph{hidden inwardness}'' \ed{cf. \emph{SKS} 12, p.
214}, which is only the hiding place of indifference (cf. remarks on Protestantism,
Bladartikler p. 95f. \ed{\emph{SKS} 14, p. 173, col. 1.}) And as / the race's
% the single slash is Koch's, as printed there; no closing slash follows on this page
\emph{perdition in sin} was sharply accentuated, marriage was condemned (cf.
The Moment VII, 21, 28f. \ed{\emph{SKS} 13, p. 296, p. 301ff.}), and the ideality
of the task called the individual away from every direction toward the external
back to himself, to the \emph{infinite task in himself}, to the \emph{God-relationship}.
Christianity was seen so ideally that not only did our age not satisfy
its requirement, but it was even declared: ``\emph{Christianity
has actually not come}\footnote[112]{cf. The Moment V,7, VI,16 \ed{\emph{SKS} 13, p. 203f., p. 267}.} \emph{into the world}; it stopped with the prototype
and at most the apostles, but even they already proclaimed it so
strongly in the direction of propagation that here the irregularity already began''
(The Moment V,7 \ed{\emph{SKS}. 13, p. 231}). [Therefore
% TR: Forbillede = prototype
% Koch prints "{SKS}. 13" with a stop after SKS; kept as printed
Kierkegaard did not give himself out to be a \emph{true Christian} either, but laid
claim only to stating truly \emph{what}\footnote[113]{cf. 73.7.} \emph{Christianity is}. (cf. Bladartikler.
p. 102f., 109 \ed{\emph{SKS} 14, p. 180f., p. 193, col. 2}. The Moment
VI,15. IX.5 \ed{\emph{SKS} 13, p. 266f., p. 374f.}. Practice in Christianity p.
241 \ed{\emph{SKS} 12, p. 222f.}. For Self-Examination. p. 13ff. \ed{\emph{SKS} 13, p. 46ff.})]

Thus the struggle continued ``in agonies such as a human being has
rarely experienced them, in an exertion of spirit that would rob another of his understanding''
\ed{\emph{SKS} 13, p. 124}, in such exertions and agonies ``that
often enough his wish was death, his longing the grave, his desire
that this wish, this longing might soon be fulfilled.'' \ed{\emph{SKS}
13, p. 134} And yet he could say: ``Yet your love \ed{74.2} O
God moves, the thought of daring to love you inspires me --- under
the possibility of being omnipotently compelled --- gladly and gratefully to
\kpage{163}will to be what is the consequence of being loved by you and of loving
you: one sacrificed, sacrificed on a race for whom the ideals are a fool's prank,
a nothing, the earthly and temporal the earnest thing, a race which
worldly shrewdness in the guise of teachers of Christianity has shamefully,
in the Christian sense, demoralized.'' \ed{\emph{SKS} 13, p. 124.} His wish
and longing were fulfilled. Only for a few months did the struggle last;
then he died, but his thoughts did not die with him, nor did the memory
of his personality.

\sepline

By the account I have given in the foregoing of Kierkegaard's
conception of the Christian, I hope it will have become clear
both \emph{how strict a coherence} there is in it --- how unjustified, therefore,
the scornful rejection by speculative dogmatics of a ``fragmentary,
unsystematic thinking'' was --- unjustified to an even higher
degree because it came from a quarter where there appeared a complete
lack of clarity and sharpness of thought, an almost comic confusion;
--- and also \emph{how many problems have been brought forward}, how many
new \emph{categories} have been drawn into the light, how many \emph{concepts}
have been more sharply and more deeply determined, how perspicuously the \emph{spheres of existence}
are ordered according to their inner relation, with how sharp a \emph{dialectic}
guard has been kept over \emph{the claim of existence}, over the ethical's and \ed{74.3} the religious'
inescapable claim on the individual in existence.

With regard to this, what is needed will have been made clear by my account,
but with more definite regard to the problems that were
set up when I concluded the critique of \emph{Feuerbach}, I will here point out
the contributions Kierkegaard has made to their solution, so as, when
that is done, to bring out the \emph{difficulties} his theory seems
to me to leave behind, whereupon the \emph{final result} of all these lectures
can then be drawn.

We then set up \emph{first} the problem: \emph{what becomes religion's relation
to reason}? does thought become religion's enemy, or can it
enter into an inner and serving relation to it? With regard to this
problem Kierkegaard has maintained the essence of religion as an \emph{existence-relation},
as \emph{life}, in opposition to every conception of it
as a merely \emph{thought}-relation, as a stage of development in \emph{cognition}; he
has sought to determine \emph{the domain of objective thinking}, to set the boundary\kpage{164},
% TR: Existentsforhold = existence-relation
% TR: Gebet = domain
% TR: Opfattelse = conception
where the \emph{paradox} emerges and where thought, in discovering
it, is used to guard it; --- he has sought, in the finite consciousness
--- in the \emph{consciousness of sin} --- what can bring the human being to
hold fast \emph{what conflicts with} \ed{74,4} thought --- and in the consciousness of sin
% TR: Syndsbevidsthed = consciousness of sin
% Koch prints "{74,4}" with a comma for the stop; kept as printed
at the same time the possibility of \emph{a relation to the divine as something qualitatively
higher than the human being}. Dialectic is thus seen as serving the
religious relation; the spheres of \emph{thought} and of \emph{faith} are separated.

We set up \emph{secondly} the problem: \emph{what becomes religion's relation
to the moral}? can it assert an independent place beside
this, or must it be absorbed by it, or must it stand in opposition
to it? --- With regard to this problem Kierkegaard
has determined the relation of the \emph{ethical sphere of life} to the religious one, he has demonstrated
the \emph{dialectic in the ethical} which prevents one from holding it fast in isolation,
the necessity that the religious must come forth, and he has
sought to maintain the concept of \emph{an absolute relation to the absolute}, an \emph{absolute
duty toward the absolute}, which is not exhausted in the ethical duties.
% TR: det Sædelige = the moral

We asked \emph{thirdly: how, when religion's independence
is maintained, are its essence and distinctiveness to be determined}? --- This Kierkegaard
has described with the analyses given in the foregoing
of the religious and the specifically Christian. By these analyses
the position of the religious among the spheres of existence, its relation
to thought\ed{74.5}, to irony and humor, and to \emph{poetic ideality}
has at the same time been determined. \emph{Religion's reconciliation} is, as the reconciliation of actuality, distinguished
from poetry's illusory reconciliation in the sphere of possibility, of the impersonally
ideal.
% TR: Virkelighed = actuality

We asked \emph{fourthly: what is the boundary between what truly issues
from religion's principle and what only through a ready relapse from the
religious to the lower stages of spirit is brought into connection with the religious and confused
with it}? --- Here the norm for the separation is given in the analysis of
religion's\footnote[114]{NB. The development in Works of Love. --- The indirect polemic against
Feuerbach's development of the relation between faith and love. (NB. The inward-turnedness
in faith; the cessation of comparison --- the conception of predestination.)}
essence, in the categorical distinction of the aesthetic
and the religious, and this distinction keeps out of
the religious not only everything that has not gone through the \emph{dialectic of infinity} and
\emph{absolute resignation}, hence all \emph{immediate existence}, everything \emph{egoistic}
\kpage{165}that does not itself pretend to be the religious, but also, of what
appears with the claim to be the religious, even the religious in a
higher sense, everything that has a remnant of the aesthetic's
immediacy in it, hence everything that seeks rest in an \emph{unshakable,
undialectical objectivity}, or relates itself \emph{poetically} to the infinite, so
that \emph{the personality itself} has not made the movement of infinity, or
that relates \ed{74.6} itself to the absolute without the \emph{possibility of offense}
becoming at every moment the annihilated passageway of the relation and
without the fear and trembling of \emph{uncertainty} becoming the ground on which the certainty
of faith rests. That separation thus gives the norm for judging
\ins{many of the ``religious'' phenomena Feuerbach has
attacked, and} such expressions as: Christian world, Christian\footnote[115]{NB. religious awakening. (The critique of Grundtvigianism, with and without
mention of names, in Concluding Unscientific Postscript).}
state, Christian art, etc.

We asked \emph{finally}: \emph{what form of religion will be able to assert itself, either beside
the moral and rational or in unity with it}? To this
Kierkegaard answers decisively: \emph{the Christian religion is the absolute}, which eternally
asserts itself as the highest form of existence, which, as heterogeneous with
history, is never absorbed by it in such a way that a new and higher form of religion
could develop; which never at any time comes
to stand in such a relation to the race's development that it needs
accommodation, reasons, defense, since it is the unconditioned,
which contains the judgment on the race, and before which
the race has to defend itself. Another, higher form of religion
is the impossible; every attempt at it is rebellion against Christianity,
a \emph{sin against the Holy Spirit}, is the intensified form of \emph{despair}.

\ed{74.7} Thus the problems set up have been partly solved, partly
at least brought considerably nearer their solution by the penetrating
analyses and the sharp determinations of concepts. I shall now
attempt to point out the \emph{difficulties} that still remain in
Kierkegaard's theory, difficulties which at least prevent
me from following him to the last point. I present these
difficulties as \emph{doubts}, for there is in Kierkegaard's theory
one point that bears the contradictions and the difficulties, at least
the essential ones, but to reach this point there is required
\kpage{166}a movement I cannot make. As such doubts, --- and as
doubts that are not to tear down but only to clarify, I ask that what
I develop in the following be regarded; I put them forward with
caution before a thinker whom I admire, with modesty
before a man who has been dear to me as few have.

% Koch: description of an inserted leaf carrying a summary of lecture 27
\ednote{}

\emph{Kierkegaard} (Lecture 27.)

The judgment on Christendom, as it is contained in what has been
developed. --- The consequences of homogeneity.
% TR: Ensartethed = homogeneity (uensartet = heterogeneous)

The character of ``Christendom.''

Kierkegaard's last writings. --- The ideal Christianity. --- ``Practice
in Christianity.'' --- Intimations of the polemic earlier than
the article against Mynster.

The course of development in the polemic: ``witness to the truth'' --- the judgment on
the established order: official Christianity --- the pastor's position ---
the Christian world, --- indifferentism. --- ``Will faith be found on
earth?'' \ed{\emph{SKS} 12, p. 17.}
% TR: Sandhedsvidne = witness to the truth

Kierkegaard's account of the course of the struggle in ``What Christ
Judges etc.\ of Official Christianity'' Sharpening of the problem.
--- The satire (``Laughter in the service of fear and trembling.'' \ed{\emph{SKS} 13, p. 301}!)
--- Emphasis on suffering from without. --- Unconditional condemnation of
the clerical estate. --- Childhood Christian faith and confirmation. --- Condemnation
of marriage. --- The exclusivity of the God-relationship. --- Doubt
whether Christianity has ever existed.
% Brøchner writes "Hvad Christus dømmer o.s.v." and Koch supplies the rest of the title after the "etc."; both kept

\sepline

\emph{Retrospect}: The coherence in Kierkegaard's theory. --- New problems
and categories. --- Sharper determinations of concepts --- The ordering of the
spheres of existence. --- The How of existence.

Problems that were posed in the critique of Feuerbach.

I. Religion as existence-relation, as life. --- The domain of objective
thinking. --- The paradox. --- The consciousness of sin as bearing
the paradox, and as maintaining God's qualitative difference from
the human being. --- Dialectic --- The sphere of faith and of thought.

\kpage{167}II. The ethical and the religious.

III. The essence of religion. --- its place in the spheres of existence. --- Humor
--- Poetic ideality.

IV. The aesthetic and the religious. --- Separating out the immediate.

V. The eternity of Christianity.

\lecture{\ed{74.7 continued}. Lecture 28}

In order not to let the critique lose itself in scattered detail, I will keep
to \emph{two} main points at which the essential difficulties
come forth; namely the determination of subjectivity's relation to
objectivity and the determination of the religious.

a) \emph{Subjectivity and objectivity}. --- As the existence-relation was maintained
against speculation, and as, in the determination of the human being as a \ed{74.8}
synthesis of the temporal and the eternal, the starting point was taken for a
polemic against speculation's illusory eternity, the result was reached
that \emph{the existing individual} had to \emph{choose} between the objective and
the subjective relation to the truth, that they \emph{could not be mediated}, and that,
since the objective relation ended in a contradiction, the \emph{subjective} became
the only possible one. And in this subjective relation \emph{the truth} \ins{
(i.e., the essential truth)} was for the subject only in the form of \emph{faith};
its definition became to be \emph{the objective uncertainty, held fast in the appropriation of the most passionate}\footnote[116]{NB. \emph{God} as \emph{postulate} (Concluding Unscientific Postscript p. 147 \ed{\emph{SKS} 7,
p. 182f.}).}
% fn 116: Koch leaves the round bracket opened before the title unclosed; closed here
\emph{inwardness}; the \emph{infinite passion} became its own
content; the \emph{negative infinity} became the one true expression in existence
for the infinite,\footnote[117]{``In the infinite passionate interest in his eternal happiness the individual is in
his utmost exertion ------ where God is \emph{negatively} in the subjectivity, which in this interest
is the form of eternal happiness''. (Concluding Unscientific Postscript p.
36. \ed{\emph{SKS} 7, p. 57.})}
the infinite negative resolution
became the individuality's form for God's being in it. And subjectivity's
highest form became \emph{faith in the paradox}, in which passion,
precisely through\footnote[118]{NB. the paradoxical (the relatively paradoxical) prior to the sphere of the Christian.
(the Socratic ignorance.)} the repulsion of the absurd, is potentiated to the highest. ---
% TR: Beslutning = resolution
\kpage{168}\emph{Objective thinking} got its own limited domain, a domain that
did not comprise the personality's deepest interest, and whose truth,
since the decisive conclusion was not possible, was always only an
\emph{approximation}, hence really a \emph{hypothesis}. But in this domain
thought gets its validity, and its warrant extends right to the
point where the sphere of faith \ed{75.1} begins; that far the use of
it is warranted, and \emph{within} the sphere of faith it gets the significance:
\emph{to guard the paradox}. But here essential \emph{difficulties} arise. It
must be asked: must not the subjective\footnote[119]{NB. the insufficient developments in Nielsen's Propædeutik pp. 74--80. By his
separation of subjectivity and objectivity, subjectivity becomes the thoughtless, the irrational,
the life of consciousness irreconcilably split. --- (NB. The subject is surely also
a member in the objective order of things.) (NB. the ethical resolution. Thought's share
in it.)}
relation, \emph{as a relation of consciousness},
translate itself into an expression in \emph{thought}, and must not then, when this has
happened, the contradiction in the thought-expression for \emph{the absolute truth} react
back upon the rest of consciousness's thought-content and \emph{annihilate its
truth}? In other words: \emph{can the paradox's contradiction in thought be isolated}? and,
if it cannot, does not all distinction between
truth and meaninglessness vanish, does not \emph{precisely the use of the understanding as a
guard for the paradox become impossible}? When the paradox states the \emph{unthinkable}: the union
of God and the single human being in the person of the God-man;
the coming into existence and becoming of the eternal \ins{NB. The divergence in the development of
this point in ``Philosophical Fragments'' and in ``The Sickness unto Death''.
--- (The becoming of the \emph{necessary})} within\footnote[120]{The coming into existence of the eternal (NB. Religiousness A and the Christian).} finitude, when
% TR: Tilbliven = coming into existence; Vorden = becoming
% "Philosophical Fragments": the Danish opens the quotation mark and does not close it; the English supplies it
by the determination sin it states the transformation of human nature,
the individual's paradoxical transformation, the loss of eternity, \ins{the
eternal which the human being nevertheless cannot get rid of (cf. The Sickness
unto Death, p. 11 \ed{\emph{SKS} 11, p. 132f.})}, which in the forgiveness of sins, in faith,
is won again, and won in such a way that it in the same moment eternally presupposes
itself, do not all these determinations then contain what,
translated into an expression in thought, comes into conflict with all the \emph{laws and
forms of thinking}, and, in expressing \emph{the truth}, makes what conflicts with it
the untrue. \emph{Or} \ed{75.2} is logical thinking not to be touched
by the contradiction because it belongs in the sphere of \emph{abstraction},
and not in that of \emph{actuality}? But whether Kierkegaard defines
% "der ... atter vinder": the subject of "vinder" (wins) is not expressed; construed as eternity won back, rendered passive
\kpage{169}abstraction as the abstract \emph{prototype} for or the abstract
\emph{reproduction} of concrete actuality, there nevertheless remains such a relation
between it and actuality that the transformation of the latter changes
the former. \emph{Or} is actuality to be capable of being split in such a way that one domain
becomes thought's, another that of the unthinkable but absolutely true?
Are we not thus led to the unthinkable representation of a \emph{divided
whole of existence}, two defective universes, of which the untrue \emph{grounds} the
true and makes it \emph{finite}. And if we let this representation stand as
belonging \emph{under the paradox itself}, the consequence nevertheless remains
that in the disrupted human nature, which has become another
with sin, thought, which surely is \emph{the expression of the essence}, has remained the
same and remains so when nature is again transformed. There remains the
difficulty: \emph{how the reception of revelation, which presupposes a re-creation
of the individual, can let the individual's ideality in its self-consciousness persist}.
% TR: Forestilling = representation (Hegel's Vorstellung)
% "Tilværelseshele": Tilværelse = existence, here "whole of existence"
The paradox is \emph{the miracle in the world of thought}, but just as the miracle in the world
of nature annulled the concept of nature, and of natural law\ed{75.3}, so
the miracle in the world of thought must overturn the laws of thought, bring
\emph{skepsis} in, and make doubtful the thinking that is to distinguish the paradox
from nonsense and guard the former in order to hold it fast as paradox,
as the absolutely unthinkable. Life seems, under this presupposition,
to dissolve into \emph{confusion}, \ins{the life of consciousness to be effaced, the
highest forms of individual existence to fall into unconditional discord
}, and the task to become \emph{also in the sphere of thought to cast the world away in order to
be only in the paradoxical God-relationship}. --- Whether the \emph{consciousness of sin} has the tensile force
to bear these difficulties, to keep the spheres apart, I do
not decide; I only point out the difficulties.
% TR: hæve / Ophævelse = annul / annulment

In the account of Kierkegaard's development of the subjectivity-relation
I have conducted myself only historically, as a reporter. Here \emph{is}
the place to call attention to what applies
both to the difficulties I have already pointed out and to those
I will point out in what follows: that the \emph{philosophical presupposition} from
which his polemical development proceeds is \emph{the conception of speculative
philosophy as philosophy's highest form}, and of the \emph{Hegelian} as
speculation's most consistent execution. Against speculation's
volatilization of existence the polemic is directed which maintains existence;
but when thought's relation to existence is determined, the
starting point becomes \emph{speculation's conception of the eternal}, the very determination\kpage{170}
% TR: Forflygtigelse = volatilization
of the human being \ed{75.4} as a synthesis of the temporal and the eternal
shows through there, and the conception of \emph{actuality} becomes the \emph{idealist} one;
it is presupposed everywhere that \emph{the sensuous is the deceptive}, the
sensuous world is illusion, \ins{sensuous actuality is seen
everywhere as something secondary, an \emph{inexplicable} addition to the eternal.} From
this presupposition it follows, when the determination of finitude is accentuated
in opposition to speculation \ins{(and for this speculation
itself gives the warrant)}, that eternity --- which must after all always
be present in order to give continuity --- determines itself for the existing individual
as \emph{something to come}, / which is always grasped only in \emph{uncertainty} and in
% the single slash is Koch's, as printed there; the next pair of slashes encloses "(which again ... inner)"
the form of \emph{abstraction}, that it therefore only in \emph{passion} gets momentary concretion,
and that \emph{the individual's own ethical actuality} \ins{(which again is conceived as
something \emph{inner})} remains as the only actuality for the individual
besides the paradoxical actuality of the \emph{object of faith}. In these consequences
lie the difficulties: that the eternal\footnote[121]{cf. the enclosed leaf.
% Koch: description of the enclosed leaf (a summary and reflections on the eternal)
\ednote{}}
gets momentary existence
and becoming, \emph{without}\footnote[122]{NB: to \emph{become} infinite.} \emph{yet being determined for thought by this}\footnote[123]{(we get the unthinkable before the paradox proper.)} \emph{becoming}.
This eternity becomes \emph{the abstract eternity of speculative philosophy}, only\footnote[124]{We get no starting point for such a concrete knowledge, on which the ethical
resolution must rest.}
in such a way that it first gets the reality of being \emph{when becoming is annulled}. But
the fact that becoming \emph{is to be annulled} posits becoming as something transient, \emph{something that
is not to be}, that eternity as what is properly true --- for the beings existing in eternity's
one true existence.
% TR: Tilkommende = something to come

If --- and to this the development of philosophy points --- \emph{another
philosophical starting point} is taken, the matter seems to present itself differently, and
that determination of the subjectivity-relation and of thought's relation
to existence to receive a \ed{75.5} different conception. If the
\emph{one, all-embracing} absolute moving itself in an \emph{eternal} process, \emph{as
it was determined by the speculative philosophy that volatilizes actuality}, does not become
the definitive truth, then the \emph{single}, the individual, must get essential
original significance, existence and thought must get a different position toward
each other. \ins{--- Thought gets an \emph{essential} hold, a real fullness in the
\kpage{171}individual consciousness, whose foundation is not something merely vanishing;
it} \emph{becomes immediately concrete through consciousness}, \ins{(thus it becomes
an activity that is not merely the activity of ``the eternal'')} \emph{and the abstract universal}
(thought's \emph{negative}) becomes \emph{what is identical in the concretion}, the \emph{a priori},
\emph{negative} foundation; \ins{The a priori --- (in its abstraction and negativity)
--- becomes in the single existence, precisely when this is seen as
a member of a comprehensive total organism, the abstract but only possible
thought-expression for the all-embracing ideal-real. (NB. its real
existence in feeling, in passion, in love.)} \emph{The single existence}
is seen here as \emph{the concretion of ``the eternal''}, the single \emph{sensuous-spiritual}
existence as something true that has content precisely through its individuality
and its unity with the other individualities in the a priori, which
is present \emph{whole} in every moment, precisely because it \emph{does not become}. That\footnote[125]{NB. to become \emph{consciousness}.} \emph{which is to come}
is then no more than \emph{that which is present} something eternal: \emph{the eternal is the individual
that in the a priori has the form of infinity and of isolation's}\footnote[126]{The individual that becomes conscious of itself as concrete and in its concretion becomes
conscious of the a priori gets in this consciousness a guarantee of an objective truth.
What cannot be left unthought without the individual's consciousness being disturbed,
without its falling into discord with itself, is for the individual \emph{the objectively true}, but the
true \emph{by virtue of the subjective consciousness}.}
% fn 126: Koch prints "Indvidet" (for "Individet"); rendered as the sense
% TR: det Aprioriske = the a priori
\emph{annulment}
(NB. the always identical a priori's relation to reality.
--- The problems of logic and of ethics meet here.). (In \emph{passion} the
individual is in its highest \emph{self-affirmation} while at the same time it \emph{affirms the universal}:
only in passion is it therefore\footnote[127]{Consciousness gets for its object an actuality outside the ethical actuality.} productive, \emph{creative} --- both spiritually
and bodily.) --- [The concept \emph{guilt}, \emph{reconciliation} and the like is not annulled
by that starting point, but reconciliation is sought \emph{within
the human} in the \emph{supplementing} individuals who in love \emph{communicate} their
essence.]

\ed{75.6} b) \emph{The determination of the religious}. --- In the foregoing the dialectic has been demonstrated
by which Kierkegaard lets the \emph{ethical} standpoint
point beyond itself to the \emph{religious}, the necessity of a \emph{higher
assistance for the individual} in order to regain himself, and the determination of the relation
to this higher as \emph{an absolute relation to the absolute}. In \emph{the God-relationship
there is} therefore a \emph{particular} relation given, which \emph{does not exhaust itself or become concrete}
\kpage{172}in the particular ethical relations, and, as the \emph{absolute duty}, can require
the renunciation of them \emph{all}. With this distinction essential difficulties
arise once again. When the requirement is set that the human being
\emph{shall at one and the same time relate himself absolutely to the absolute τελοσ and relatively to the relative ones,}
then it must be \emph{asked}: \ins{since to relate oneself to an absolute
τελοσ is \emph{always} to relate oneself to it, for to relate oneself to it \emph{now and then} is
to relate oneself to it \emph{relatively}, but to relate oneself to it relatively is to
relate oneself to \emph{a relative τελοσ,}} with what power does the human being
then relate himself to the relative ones? If it is a power \emph{different} from the one
with which he relates himself to the absolute, then the latter surely becomes
limited, \emph{relative}; but if it is a part of the latter power, then he surely relates
himself, \emph{in} the relation to the relative ends, to the absolute.
--- Can the human being, it must \emph{further} be asked, renounce \emph{all} the relative
ends? He can do so in one moment, in which, in the resolution, he
gathers up existence \emph{in an abstraction} over against \emph{the abstraction of the eternal};
but \emph{in actuality}, in which he exists in this determinate
concretion, he renounces only the one ``finitude'' \ins{(the one relation)}
\emph{with the power by which he exists in the other.} And even the
power by which the \emph{resolution} is taken, with which one abstracted
\ed{75.7} from the determinateness of existence, is a determinate,\footnote[128]{cf. Concluding Unscientific Postscript p. 361 \ed{\emph{SKS} 7, p. 428.}.} \emph{limited}, finite
power, \emph{in the direction of the abstract.} --- But when it thus seems,\footnote[129]{NB. The expression: the ethical-religious.}
that the religious does not come to be distinguished from the ethical, but that
it becomes \emph{the soul in the ethical,} which latter cannot be thought through
without the former, then this also finds its confirmation when it is held fast
\emph{that the ethical comes forth only as the movement of infinity is made}. The
\emph{eternal}, to which the relation is in the religious, thus does not come forth \emph{first}
on the domain of the religious; it first comes forth there as \emph{the
all-embracing eternal} (in that it is not sought merely in the individual, which \emph{asserts
itself by itself),} but precisely \emph{through the abstraction}, \emph{in which it remains} (as \emph{the
infinite negative), it is not to be distinguished from the eternal of ethics}, and it seems
one-sided to hold fast the single ethical duty as \emph{merely single}, finite,
without the eternal and the relation to it coming forth in each single
one. Should the eternal, on the domain of religion, have the closer
\kpage{173}determination of \emph{personality}, then\footnote[130]{That the individual, which in the consciousness of guilt collapses before the ideal requirement of the ethical, must seek a higher assistance in order to regain itself, Kierkegaard at once determines more closely thus: that the higher, to which the individual comes into relation, becomes a higher \emph{outside} individuals, so that the relation to the higher thus becomes a \emph{particular} relation, an absolute relation to the absolute, which constitutes an absolute duty toward the absolute.} determinations would be required here
that are \emph{not} given (--- it is only said: God is \emph{subject}, and as such exists
only for subjectivity in inwardness), and his compatibility
with the absolute cannot be demonstrated for thought. --- And \emph{finally} it must
be asked: does not precisely the dialectic by which
Kierkegaard arrives at giving the ethical subject \emph{content} again presuppose a
different\ed{75.8} fundamental philosophical view from the one from which
Kierkegaard takes his point of departure, and \emph{precisely the one that was
intimated above} and that led to other possibilities both in the solution of the
problems touched on earlier and in this one? When difficulties thus arise
in the determination of \emph{the relation of the ethical to the religious},
no fewer arise in the determination of \emph{the relation of the Christian
to the preceding spheres}. What Kierkegaard calls \emph{Religiousness
A}, and which, without corresponding to a determinate \emph{historically manifested} form of
religion, is to comprise the actual religiousness that is not yet
the Christian (thus presumably also the \emph{Jewish}, although e.g. in ``Fear
and Trembling'' \emph{Abraham} is depicted as the father of ``faith''), and which presumably most nearly
recalls \emph{the conception of God in Greek thinking} --- contains a \emph{relation
to the eternal}, and it is conceded that it at least \emph{could} take
place in paganism, since it had only \emph{human nature as such} as
its presupposition. But when the determination of Christianity comes
forth: the determination \emph{hereditary sin}, the loss of eternal being, the ignorance
of this loss, then the question becomes: \emph{in what sense, then, can
the human being, from this point of view, have related himself to the eternal}? and if he has
related himself to the eternal, \emph{what difference then remains between the eternal to which
pagans related themselves, and the eternal that becomes the determination of the Christian God}!?
On this point there seems\ed{76.1} also to be an \emph{uncertainty} in
Kierkegaard, for while on the one hand the concession touched on above
is made, \emph{there is missing} on the other hand \emph{that the pagan
had representations of God} (The Sickness unto Death, p. 119 \ed{\emph{SKS} 11. p.
228}), and when the consequences of the determinations of the Christian are
% doubtful construal: "saa mangler der ... at Hedningen havde Guds-Forestillinger" kept literal; the sense seems to be that K. withholds from the pagan any representation of God
% Forestilling = representation (Hegel's Vorstellung) here and below
% Koch prints "SKS 11. s. 228" (full stop for comma); kept as printed
\kpage{174}rigorously carried through, we arrive not only at the result \emph{that Christianity
has never come into the world}, but also \emph{that Religiousness A becomes only
a possibility} in uncorrupted human nature, and so, under the presupposition of
hereditary sin, has never become actual either;
and finally, \emph{that neither has the ethical ever been able to exist}, \ins{
not merely understood in the sense that it has not been able to exist \emph{justifiably
as} ethical,} \emph{since} the ethical is only a sphere of transition, which cannot
be thought through without reaching the\footnote[131]{therein lies the consequence that the one possible existence in actuality becomes the one that recognizes the ground of concrete individuality as something essential.} religious, and which, held fast
as the ethical in its isolation, would become sin, despair;
--- but in the sense that it has not been able to exist at all,
since it presupposes a relation to the eternal. This leads us on to
a \emph{new difficulty}. When we proceed from the Christian conception of
existence, from the \emph{relation of opposition} between Christianity and the world,
from the requirement\footnote[132]{Must not the content of existence vanish in the \emph{absolute opposition} to the eternal truth?} to \emph{die away}, to hate one's natural life, \emph{does it not
then become impossible to get}\footnote[133]{cf. the abstract character of the historically manifested ``Christian ethics''.} \emph{an ethical content to subsist together}\footnote[134]{(NB. the abstraction of eternal happiness. Concluding Unscientific Postscript pp. 300--06 \ed{\emph{SKS} 7, pp. 356--65.}).} \emph{with the Christian}?
Is not the incompatibility conceded in essentials by the polemic
against marriage? \ins{cf. The Moment VII,32 \ed{\emph{SKS} 13, p.
304f.}.\footnote[135]{cf. the consequences of Kierkegaard's polemic against the established order demonstrated in previous lectures.} Practice in Christianity p. 267. \ed{\emph{SKS} 12, p. 244f.}}
Can what is touched on in Concluding Unscientific Postscript p. 374 \ed{\emph{SKS} 7, p.
442f.} lead to an opposite result? Does it become possible to let
the ethical keep a place, just as the understanding was to keep
its place in spite of the paradox? \ed{76.2} But if it does not become possible, \emph{will not
the ethical then constantly remain standing as a justified protest against the Christian}?
will not the one who sees something divine \emph{in} the ethical precisely
thereby find an insurmountable obstacle to making the leap into
the\footnote[136]{NB. Comparison of this critique's results with the earlier set-forth determination of the Christian view of life as \emph{abstract-spiritualistic}.} domain of the Christian?
% Tilværelse (existence at large) and Existents (the existence of the one existing) are both "existence" in this span
% TR: Virkelighed = actuality
% TR: Bestemmelse = determination
% TR: Religiøsiteten A = Religiousness A
% TR: Arvesynd = hereditary sin
% TR: Salighed = happiness (evig Salighed = eternal happiness)

\kpage{175}These difficulties in the determination of the relation of the ethical to
the religious, and of the relation of the Christian to the preceding spheres of life,
seem to me to be the essential ones alongside those developed under
a). Yet alongside them stand \emph{others}, which are not
of slight weight; \emph{thus} when God is determined as the qualitatively
% Koch prints "Bestemnelse" (for "Bestemmelse"); rendered as the sense
higher than the human being and yet the determination of the divine
is kept in the form of ignorance and of negation; --- \emph{when} the relation of the human
and the divine in faith is to be determined,
where on the one hand faith becomes a \emph{requirement}, a ``You shall,'' and yet
faith can be given only by God, so that what is left over for
the human being becomes only this: by himself to become nothing, self-annihilation
(\emph{NB.} its possibility), his only capability this: to become
conscious that he is capable of nothing (NB. Freedom and grace --- predestination.
The remark about Socrates in Kierkegaard's last writing
\ed{cf. \emph{SKS} 13, p. 405, The Moment 10}. Faith in the absurd in the direction
of redemption); --- \emph{when} the Christian dogmatic determinations
are rendered in a form in which they seem to acquire another content,
e.g. when an \emph{eternal consciousness} is posited as \ed{76.3} a synonym for \emph{eternal happiness;}
and \emph{finally} when the dialectic employed, by which the determination
of guilt is posited in its totality,\footnote[137]{NB. The determination of the essence of the human being, the self's relation to God. (The Sickness unto Death --- Either/Or II.)} seems to throw the guilt from the individual
over onto existence or the author of existence, and\footnote[138]{NB. the ``losing the condition for the truth'' (Philosophical Fragments \ed{SKS 4, p. 223f.}) The concept: \emph{eternal truth.} --- the extent in which it is used by Kierkegaard.} to make
guilt a \emph{metaphysical} concept, which becomes = \emph{finitude} (cf. Concluding
Unscientific Postscript p. 405f. \ed{\emph{SKS} 7, p. 479ff.}) (NB.
Kierkegaard's relation to \emph{Protestantism}. --- Faith\footnote[139]{Bladartikler p. 95 \ed{SKS 14, p. 173.}} --- imitation,
suffering.)
% TR: Efterfølgelse = imitation
% TR: Selvtilintetgjørelse = self-annihilation

The difficulties developed seem to me to remain
in Kierkegaard, and they prevent me from following him to the
final goal where he found rest. I do not, however, present them as
\emph{definitive} objections to his theory, but have on the contrary pointed
to the point in the theory itself in which the theory's tensile force
lies, and it is perhaps --- I do not decide that --- for the one who can
\kpage{176}make the movement by which he reaches this point, great
enough to bear all these difficulties. For him the demonstration
of them will then only give the relation to the paradoxical a new intensity; for
others they must help to make the \emph{choice} with clearer consciousness.\footnote[140]{But that point: the consciousness of sin presupposes faith: the theory thus closes itself in a circle into which only a leap leads.}

\sepline

From these developments, which I do not want regarded as a critique
properly carried through, least of all as a rejecting one, I can now pass
on to concluding these lectures, in that for this purpose I first
briefly recapitulate their course, \emph{and the}\ed{76.4} \emph{partial results gained},
in order \emph{then} to \emph{gather the final result} from these. ---

\sepline

% Koch: the inserted leaf Brøchner referred to at 75.4 follows
\ednote{}

\ed{1} \emph{Kierkegaard}: (Lecture 28.)

1) \emph{Subjectivity and objectivity}. --- can the contradiction for thought in the paradox be
isolated? --- NB. Abstraction and actuality --- the splitting of actuality
--- the relation of thought to essence ---receptivity to
revelation. The identity of the individual. --- The paradox as ``the miracle in
the world of thought'' --- Only the God-relationship seems to have to remain.

Kierkegaard's philosophical presupposition. --- The conception of
the eternal. --- Actuality. The sensuous. --- Consequences. ---
The difficulties: the becoming of eternity, its abstraction. --- The untruth
of becoming.

Intimation of another philosophical point of departure (see opposite
page).

2) \emph{Determination of the religious}. Kierkegaard's determination of the relation
between the ethical and the religious. --- Difficulties: at one and
the same time to relate oneself absolutely to the absolute τελοσ and relatively to the
relative ones. --- Can the human being renounce \emph{all} the relative ends? --- The
eternal \emph{in} the ethical in relation to the eternal in the religious. --- the
single duty. --- God's personality --- The dialectic by which, in
Kierkegaard, the ethical gets content.

\kpage{177}\emph{The relation of the Christian to the preceding spheres}. --- Religiousness A. --- has
it existed? (NB. The relation in it to the eternal). --- Uncertainty in
Kierkegaard --- The consistent result.

Christianity's relation to the ethical.

% Koch prints "M{enne{skets}" (his second brace turned the wrong way); rendered as the sense
\emph{Other difficulties}. 1) God as something qualitatively higher than the human being.
2) The relation of God and the human being to faith. 3) The rendering
of the Christian dogmatic determinations. 4) The dialectic of guilt.

\sepline

Transition to what follows.

\ed{2} Intimations of another point of departure (see opposite page) What
is to be \emph{avoided} by them is: the contradiction that the eternal acquires \emph{momentary}
existence, and a \emph{becoming that does not determine it}. What is to be \emph{gained}
is: 1) that eternity becomes just as fully something \emph{present} as something to come,
--- 2) that thought in individual existence gets a ground
that gives it \emph{content}, and that this content, since the ground is something \emph{essential},
acquires \emph{efficacy}, so that the certainty of the abstract eternal does not become the
only certainty, the ideal ethical actuality not the only \emph{actuality}.
3) that in this knowledge of an essential ground, in which a
knowledge of the a priori is contained, there is given a guarantee of an \emph{objective}
knowledge, a \emph{real knowledge} by virtue of the subjective consciousness.
% Koch prints "at de i denne Viden" (for "der"); rendered as the sense
% TR: Enkeltexistents = individual existence
% TR: Grundlag = ground

\sepline

Determination of ``\emph{the eternal}'' (cf. 75.5)

The relation of subjectivity. --- The certainty that is necessary in the ethical
domain in order to be able to begin to act. --- The certainty in the God-relationship.
---

\sepline

\emph{Life} and thinking.

\lecture{\ed{76.4 continued} Lecture 29.}

1. \emph{Recapitulation of the content of the lectures}.

(With regard to the part of the lectures that falls in the \emph{previous
semester}), the draft for Lecture 1 of this semester is used. ---
The \emph{relation of divergence,} which already then showed itself between the principle of Christianity\kpage{178}
and that of modern philosophy and of the modern age in general,
now finds its closer determination through what has been developed in \emph{this semester},
especially through \emph{Feuerbach's} and \emph{Kierkegaard's} theories. --- Both
agree in conceiving the relation as a \emph{relation of opposition}, in
wanting to keep the two spheres apart from each other; although the interest is
different in them, and the individual analyses essentially divergent,
they nevertheless agree in --- what also became our \emph{provisional} result --- \emph{that Christianity
is not the ruling spiritual power of the modern age}. By the intimations given in the critique
of Kierkegaard's theory it has also been confirmed
that what, --- in opposition to Christianity --- governs
the spiritual life of the modern age is a \emph{positive} principle of life, which is able to
unfold itself into a coherent view of existence which
lacks neither the distinctive character of the \emph{ethical} nor that of the \emph{religious}.)
\ed{76.5} With regard to what has been developed in \emph{this semester}, the presentation
and critique of the \emph{attempts at restoration} are first recalled. It is
first pointed out in general how this critique on the whole led to
the same result as the critique of the similar attempts discussed in the previous semester,
namely: \emph{that} they were unsuccessful, and \emph{that} what prevented
them from succeeding was precisely the influence of the tendency heterogeneous
with Christianity and characteristic of the modern age.
% TR: Retning = tendency
% TR: Restaurationsforsøgene = the attempts at restoration

In particulars, then, with regard to \emph{Fr. Baader}, the result gained
by the critique of him is recalled: that all \emph{theosophy is fantasy}, and
\emph{all immediate theosophical cognition of nature fantastical,} --- \emph{the
deviation from the Christian} in him is recalled. With regard to \emph{Schelling's}
philosophy of revelation, the critique of the concepts is recalled:
``\emph{the individual being that is all}'', and ``\emph{free cognition}'' --- the determination
of the philosophy of revelation as \emph{fantastical theosophy}; the demonstration
of the influence of the representation of \emph{God's personality}, and of
the reason for the \emph{deviations from the Christian}. With regard to \emph{speculative
dogmatics}, the critique of the determination of \emph{the essence of
religion}, of \emph{the science of faith,} of the \emph{general} \emph{conception of existence} is recalled,
and the critique of the development of the \emph{individual dogmas}.

With regard to the \emph{negative tendency,} which \ed{76.6} was discussed as
the opposite of speculative dogmatics, the \emph{problems} brought forward
with it are recalled: the relation of reason as an inwardly moving
power in the religious, and: the relation of philosophy and religion to
the individual.

\kpage{179}Passing then to the last section of the lectures:
\emph{the separation of religion and philosophy}, Feuerbach's and Kierkegaard's
views are first provisionally juxtaposed. Then, with regard
to \emph{Feuerbach}, the most characteristic \emph{results} at which he
arrives are brought out: ``The absolute essence is the human being's\footnote[141]{``It is impossible for human thought to think something qualitatively higher than the essence of the human being.''---} own essence'';---
``The infinity of the object is the expression of the organ's own infinity'';
--- ``Ignorance of the divine is the negation
of the divine'', --- ``Religion is the human being's first
and indirect self-knowledge.'' --- Religion's systole and diastole;
--- the subjective and practical character of religion; --- the contradiction
in the religious representations; their necessary dissolution into the
anthropological; --- the opposition between faith and love,
the consequences\footnote[142]{The necessity of ``atheism''.} of this; --- the development in Feuerbach's later
writings (Wesen der Religion --- Theogonie.) --- (Reference to the critique
of Feuerbach.).
% Koch prints "F{eurbach}'s"; rendered as the sense

As regards \emph{Kierkegaard}, reference is made only to what was developed in the previous lectures,
and \ed{76.7} since this returns to what was set forth
at the beginning of the hour, it proves to be the task still
remaining to set the two opposites, which show themselves
historically to stand over against each other, against each other, so that the
final choice can then be made. This becomes the task for the concluding lecture.

\sepline

% Koch: the last lecture sketched very briefly on an inserted leaf
\ednote{}

\ed{1} \emph{Lecture 30}.

Connection with the preceding lecture.

Of what character will the ethical-religious view of life that is different from Christianity
be? The answer negative --- positive.

Results of the critique of the realistic series of development --- of the
idealistic.

The task of philosophy. --- Feuerbach's attempt at a solution of it.
\kpage{180}--- The difficulty at which he stopped (compare
Herbart.)

Outline of a conception of existence: individual existence --- the whole
of existence. --- Unity of nature and spirit --- The ideal-real.
--- The natural ground of the ethical. --- How the ethical comes
forth. --- The definition of ethics. --- --- The a priori in the ethical.
--- The relation to time --- Avoidance of the egoistic. --- Resignation.
--- Guilt, repentance --- reconciliation, higher assistance.

Relation to the modern age's conception of nature --- relation to Feuerbach.

The religious moment in this ethical view of life. --- The relation of faith
--- The absolute. Its determination. (Ignorance). --- The
pantheistic. --- God's personality. --- In what the personal relation
to the absolute consists. --- The persistence of personality after
death.

(verte)

\ed{2} Elements for this view of life in the modern age. --- The one-sidedness
and superficiality in the various theories. Their ground.
--- The possibility of avoiding them.

Relation to the view of life of Greek antiquity. --- Influence of
the Christian --- The conception of the ethical --- of the natural.

The Christian and the humane. --- The determination of both --- The mixed forms.
--- The essential element in them. --- The necessity
of a distinction and a choice.

\sepline

Conclusion.

\lecture{\ed{76.7 continued} Lecture 30}

2) \emph{Final result}. --- Since the \emph{provisional} result gained earlier with
regard to the relation of Christianity and the modern age's view of life
has been confirmed by the theories last presented; since
in particular the Kierkegaardian analysis of the Christian,\footnote[143]{(Through Kierkegaard's analysis of the Christian a double result has been gained: 1) the disagreement between Christianity and philosophy, 2) that our age does not live in the Christian determinations. But since, as earlier, the tendency of the modern age has been shown to agree with the principles of modern philosophy, that double result is gathered into one.)} which with
\kpage{181}strict consistency attaches itself to its most concentrated expression,
and only develops what is given in it, has led to seeing the divergence
in its sharpest form; and since at the same time the critique of the
Kierkegaardian determination of the relation between the ethical and the
religious has intimated the possibility of a \emph{unity} of these that is \emph{different
from the Christian}, then, in conclusion, this view of life is accordingly to be
worked out more definitely in its relation to the Christian, so that the choice
between them can then be made.

% Koch prints "af hvad Beskaffende" (for "Beskaffenhed"); rendered as the sense
The question is then asked: \emph{of what character will the ethical view of life} \ins{\emph{be,
for which the elements are given in the tendency of the modern age and which, as it combines
these elements into a thoroughgoing whole, comes to}} \emph{stand over against the
Christian}? \emph{Negatively} \ed{76.8} the answer to this question is given with
the historical presentation and critique of the two main tendencies
of modern philosophy; for a \emph{positive} answer the elements lie \ins{in the earlier
analyses of the principles of modern philosophy, and} in the intimations
that are given in the critique of Kierkegaard's view.

The \emph{negative} result that was gained by that presentation was
this: \emph{that neither one-sided realism nor one-sided idealism was able to
form the foundation for an ethical outlook on life that could be carried through}. It became
evident in the critique of the \emph{realistic} series of systems that it
negatively had to end in \emph{skepticism}, positively in a \emph{materialism}, in whose fundamental concepts
the dualism of medieval thinking asserted itself,
so that it could not be carried through without falling into the
contradiction with itself of having to take up concepts which it could
not deduce but could not do without either, and, by taking up what is
inexplicable by its principle, itself admitted its one-sidedness, which was given
by the fact that this tendency's fundamental concept, \emph{matter}, which was posited as
\emph{all}, was determined as it had to be determined \emph{within the
dualistic opposition to spirit}, and therefore constantly had to be supplemented
unconsciously by the opposite. --- At this tendency's terminal point
every ethical foundation was lacking, \emph{everything became physics}, and the first concepts of physics\ed{77.1}
stood unexplained and inexplicable.
% TR: Livsbetragtning = outlook on life

In the critique of the terminal point of the \emph{idealistic} series, the \emph{Hegelian}
system, it likewise became evident that neither could this tendency, in its
% join with k172-181: Danish "at ikke heller denne Retning i sin | Gjennemførelse kunde constituere"; the English here assumes p. 181 ends "that neither could this tendency, in its"
% TR: Virkelighed = actuality (Hegel's Wirklichkeit); virkelig = actual / really, by context
\kpage{182}full execution, constitute an ethical view of life. It became
apparent that a one-sidedness similar to the one that had made materialism
end in a self-examination made the idealist tendency
end in a \emph{volatilization of actuality}, whereas precisely the requirement was made
that it be taken up into thinking. The unity of being and thought became
merely the unity of monistic thinking with itself, actuality
became merely a vanishing phenomenon in the eternal process of the absolute,
nature came to be only for human thought, while it was eternally
sublated in the absolute idea's closing together with itself; the
foundation of the concrete ethical disappeared, as the eternity of thought --- and
the eternity of \emph{abstract} thought --- became the one thing that is, in which nature
and existence in time became merely something vanishing, which did not make
the idea richer in content. Ethics was here transformed into \emph{metaphysics},
and all the problems of \emph{existence} remained unanswered.

The one-sidedness of the two tendencies pointed toward another
form of philosophy, whose \emph{task} could be expressed negatively and quite abstractly
thus: \emph{to recognize the distinctiveness of spirit and of nature
without dualism}. An attempt at this task's \ed{77.2} solution was made by
\emph{Feuerbach}, in that he asserted the significance of the \emph{sensuous}, the \emph{natural}, the \emph{material},
the essentiality of \emph{individual existence} as the ground of the abstract
universal, and sought in nature not merely dead matter but
\emph{life and spirit}. But however far his \emph{sensualism} departed from the
materialism of the preceding century, reawakened in the natural sciences of our time,
the critique of his system nonetheless showed how
he, in developing the consequences, ran aground on difficulties
he could not remove, how his very premises pointed
beyond \emph{that} at which he stopped; the dissolution of existence into atomistic
individuals --- a dissolution that is also characteristic of
another, richer system, which makes the concept of individual being
its fundamental concept: the \emph{Herbartian} system.

% Tilværelse = existence (the whole of what is), as distinct from Existents = existence (the individual's)
From the point at which Feuerbach stopped, we are thus
directed, \ins{when the results of the historical critique are recalled,} toward a
view of existence in which the \emph{essentiality of individual existence}
is asserted, but \emph{not atomistically}; for the independence and eternity of the atomistic
contain a contradiction for thought; but in such a way that
the thought of a \emph{unity and wholeness of existence}, an \emph{organic order of existence},
without which thought loses its hold in itself and dissolves into skepticism\kpage{183},
is joined with it. And as the individual existence is seen as an \emph{inseparable
unity of nature and spirit}, as the being of spirit is \emph{indissolubly} bound to
the being of nature, \ed{77.3} \emph{the ideal that unites existence becomes at the same time
something real}, something actual, something \emph{natural}. Human existence then acquires
a \emph{natural foundation}, determined not as that which is to be given up, but as
\emph{that which is to be taken up into consciousness, to be permeated by it}. The \emph{ethical} emerges
as the single individual in his\footnote[144]{NB. Kierkegaard's dialectic with regard to the ethical.}
individual distinctiveness
knows himself \emph{as a member of that comprehensive whole}, has himself \emph{as a task as such},
and thereby makes \emph{the movement of infinity}. \ins{Ethics becomes \emph{the doctrine of the
realization of the human essence through free human action}, but the
human essence is realized as its natural foundation is permeated
by the life of consciousness and consciously ordered into a higher whole.}
This ethical then acquires a \emph{content}, precisely through \emph{the individual's natural determinateness},
and it acquires \emph{its regulative in the a priori}, which in the domain of ethics
acquires the same validity as in that of logic. It has \emph{time} and
\emph{actuality} as its stage, not the abstract eternity of speculation;
it does not lose itself, as materialism does, in the egoistic, \emph{for
the individual knows himself only as a member of a higher whole}. It does not attempt to
hold itself fast in an isolation, it rejects \ins{not \emph{resignation},} not
the determination of \emph{guilt} and \emph{repentance}, not the need for a \emph{reconciliation}\footnote[145]{NB. can philosophy recognize a disturbance in the whole of existence? If
it cannot, what then becomes the consequence with regard to the concepts freedom,
guilt, etc.? (cf. The Sickness unto Death).}
and a
\emph{higher assistance}. But resignation remains when it is given an ethical \emph{content
by virtue of the ethical itself}; as each single form of the ethical is seen
as an expression of a relation to the absolute, there is no absolute
relation \emph{outside} the ethical, but it becomes \emph{omnipresent in the
ethical}. The whole compass of the ethical cannot then be given up for an ``absolute\ed{77.4}
duty'' that does not express itself in it, but the ethical \emph{from a
lower sphere} can with justification be given up for the ethical \emph{in a higher,
more comprehensive one}. The ethical individual who recognizes the ethical's
\emph{infinite requirement} bows under the feeling of \emph{guilt}, recognizes the justification of \emph{repentance},
and the need for help \emph{from outside}; but as \emph{the whole
of existence} is seen as borne by the all-embracing, \emph{its harmony as
a whole therefore cannot be thought of as disturbed} \ins{and therefore the individual existence's\kpage{184}
\emph{essence} cannot be thought of as transformed either,} then the help is sought not
outside existence, but the reconciling powers are sought \emph{in} existence,
in the \emph{other, supplementing individual beings} with whom the individual joins
\emph{in love} \ins{in \emph{mutual communication of essence}.} In love for
the human being as human being \emph{reconciliation} is then sought, and in it is seen
not merely a relation to the single individual but \emph{a relation to that which binds together,
an expression of the unity}.

This view thus appropriates \emph{the modern age's conception of nature, divergent
from Christianity}. It recognizes the \emph{essentiality of the natural},
its \emph{positive} significance for spirit, its significance as the \emph{foundation}
of a \emph{concrete ethical life}. Nature becomes not merely that which \emph{gives
thought fullness}, that without which thought loses itself in the abstract or the
fantastic, but it becomes \emph{that which bears the whole of human life} and
gives it fullness and content. \emph{To master nature} through the cognition
of nature's laws can become an \emph{ethical} task; \emph{in all ethical relations
nature becomes the foundation}.

Likewise this view appropriates \ed{77.5} what is true in \emph{Feuerbach's}
theory. \emph{Ethical relations are seen as religious}, precisely because the relation
to the absolute is expressed \emph{only} through them; \emph{in the relation of human being to
human being} the human being realizes \emph{the highest}; this relation has \emph{absolute}
worth in and for itself.

When the \emph{religious} moment in this ethical is to be determined more precisely,
it becomes that by which the human being sees every relation as
resting in the absolute, sees the all-embracing eternal at every
point of existence as bearing it. There comes to be room for a relation of \emph{faith};
for this absolute, eternal, which holds existence together,
is not \emph{immediately} visible to the single consciousness; it becomes \emph{a
postulate} of the individual consciousness, which without it would become confused in
itself, would lose its hold in itself. Faith then rests \emph{on the human
essence's very consciousness of itself}; it does not come into conflict with
this, but \emph{confirms} it. It becomes a faith in that which is \emph{higher than
the individual}, in that it embraces \emph{all} individuals and \emph{the whole of} existence. But
thought has no determinations for it higher than those that express
its own essence as \emph{human} thought; it cannot negate itself
by thinking something that is \emph{qualitatively}\footnote[146]{NB. Feuerbach's developments.}
\emph{higher} than the human\kpage{185};
it therefore has only \emph{abstract} determinations for the all-embracing absolute;
% TR: Forestilling = representation (Hegel's Vorstellung); here "religiøse Forestillinger" = religious representations
% Danish word order: "forsaavidt bliver[147] Uvidenheden Formen"; the anchor stays on "becomes"
\emph{to that extent} \emph{ignorance} becomes\footnote[147]{(In this way contradictions in the common religious representations are avoided.)}
the form of the relation
to the absolute, and the abstraction of ``ignorance'' first acquires \ed{77.6}
\emph{determinations} through \emph{the cognition of the real universe}, which is seen
as revelation of the absolute, and \emph{through the ethical}\footnote[148]{This religiousness becomes \emph{a religiousness without dogmas}, --- (different from the dogma-less
religiousness of the philosophy of faith in that it takes seriously the unity of the divine
with the human and the natural; --- different from the Schleiermacherian
in that the absolute is in the actuality of the world itself); --- its \emph{cultus} becomes \emph{action},
ethical action; its ideal is that \emph{the whole of life} becomes cultus.}
\emph{relations} that
include a relation to the absolute.

If one objects against this religiousness that it is \emph{pantheistic},
I shall not attempt to repel this objection. \emph{All consistent philosophy
must become pantheistic}; for as an attempt to bring about a unity and wholeness
of cognition, it can posit no break in
the world of thought; but the highest thought must be linked
with the whole series of thoughts, and becomes that which bears the others within itself;
but \emph{this highest thought becomes philosophy's expression for the divine}. The determination
\emph{personality} philosophy will therefore not be able to apply
to it; every attempt to prove God's personality by way of thought
must, as has already been developed at an earlier point of these lectures,
fail; for in these attempts determinations must be
left out without which we could not conceive of personality
or see it given in actuality, \ins{and precisely the distinguishing of the
divine as something personal leads to all the contradictions in the
religious representations that \emph{Feuerbach} has demonstrated, and that we
earlier found already brought out by \emph{Spinoza}.}

The \emph{personal relation to the absolute} then becomes \emph{the personal relation, permeated by faith in
the all-embracing absolute, to other personalities, the communication
of the essence of one's personality to them in love. The appropriation of their personality
in reconciling love}; and this communication and this \ed{77.7}
appropriation give the single personality life and duration in \emph{the form of
personality} even \emph{after the dissolution of personality by death}.

In the \emph{fundamental view} on which modern philosophy and
the whole life of modern times rest, and to which we have traced back the various
nuances of its forms of appearance, we find the \emph{traces}
\kpage{186}of such an ethical-religious view of life, different from the Christian.
They emerge in \emph{the faith in the human}, in the undisturbable
good in the human essence, in \emph{the recognition of nature's} essential
significance for human thought and human life, in
% TR: Fornufttro = rational faith (Vernunftglaube)
% TR: Humanitet = humanitarianism; det Humane = the humane
\emph{rational faith}, \ins{in \emph{moralism},} in \emph{humanitarianism}, in \emph{the philanthropic}, in the
fundamental thought from which the more significant \emph{socialist} systems have proceeded.
They have been present everywhere, more recognizably or less recognizably,
and everywhere in connection with one-sidedness or superficiality.
Their \emph{one-sidedness} has been grounded in the opposite one-sidednesses of realism and
idealism; their superficiality in the fact that the ethical problems
were not grasped more deeply, that the faith in the human, that
humanitarianism and philanthropy did not have \ins{\emph{resignation} and}
\emph{the consciousness of guilt} as their background, that therefore ethical earnestness was lacking
and the problems of will and freedom did not stand out sharply, the need for reconciliation
and the religious conception\ed{77.8} of ethical relations
were lacking. But, as has been shown above: \emph{these defects and one-sidednesses
can be supplemented and have in part been supplemented without the ground of the view of life
being given up}, without the significance of nature and of the human, of human
thinking and of human ethical relations being distorted
or disappearing.

In this view of life the humane and natural view of life of \emph{Greek antiquity}
is preserved, but it has \emph{through the Christian} not merely
acquired the \emph{universal} character that antiquity lacked, but antiquity's
aesthetic-metaphysical conception of the ethical has been replaced by a
\emph{really ethical} conception, in which all the determinations that antiquity
hurried past, the determinations of \emph{will, freedom} and \emph{guilt}, are more sharply
specified, and \emph{the end point of the ethical} is no longer the \emph{political
relation to the state} but the \emph{ethical-religious relation to humanity}\footnote[149]{NB. Works of Love.}
\emph{and to the absolute
revealing itself in it}. The conception of the \emph{universe} has become a different one,
as antiquity's \emph{metaphysical} conception of the physical has been replaced
by a \emph{physical} one, for which the natural \emph{as such} becomes something \emph{real}, and
as something real acquires the possibility of becoming \emph{a content-rich material for thought},
a \emph{foundation for an ethical relation}, both of which acquire with it the fullness of \emph{actuality}.
% "der med det faae": plural verb, construed as thought and the ethical relation (cf. p. 184: nature gives thought fullness)

\emph{Two powers} thus stand opposed to each other in the modern \ed{78.1} age,
\emph{two different views of life} assert themselves in the new age. We can
\kpage{187}name the opposites: the \emph{Christian} and the \emph{humane}. What the
\emph{Christian} is, Kierkegaard in particular has stated sharply and clearly; through
his analysis it has become plain what lies in Christianity's
fundamental determination; it has been shown that the Christian so determined
decisively distinguishes itself from all the phenomena that I have earlier
brought forward, \ins{which is what its relation to \emph{thought} and to
the whole of \emph{natural} existence becomes, --- (NB. abstract spiritualism --- dualism.)
---} and that, if the Christian is \emph{not} what he states it
to be, then it cannot be distinguished categorically from what historically
precedes it, that it would thus have no specific
distinctiveness.
% doubtful construal: "hvilket dets Forhold bliver til Tanken" (right-column note) read as "which [distinguishing] is what its relation to thought ... becomes"

What it is that I call the \emph{humane}, I have sought to demonstrate
throughout the whole lecture course by bringing out moments in the distinctive view of life
of the new age, and now in what immediately precedes
by sketching the outline of the ethical-religious view of life
that expresses it. With the help of this it will be easy, in all the
\emph{mixed forms} in which the Christian and the humane seem to have
entered into a \emph{unity} \ins{(``Christian world'', ``Christian state'', ``Christian
ethics'', ``Christian science'', ``Christian art'', etc.)}, to
carry out the separation. It will show itself everywhere \emph{that the Christian here is only
a name, a reminiscence, a remnant for the imagination, but that the content belongs to the
humane}. It will prove to be \emph{an illusion} to believe that in these forms one dwells in
the Christian, which is only \emph{an aesthetic admixture}, an \emph{indeterminate}, \emph{featureless} something\ed{78.2},
which nowhere has validity \ins{and by which no \emph{actual}
relation to Christianity is constituted.} \emph{And yet the admixture
of it is enough to rule out that a wholly human life is lived, that one lives
really ethically in the ethical}.
% Koch prints "synes at have undgaaet en Enhed" ('avoided a unity'); the sense requires "indgaaet" ('entered into'); rendered as the sense
% "ubestemt Ei{en}dommelighedsløst": Koch's italic is broken by his supplied letters; kept as two spans

So there remains then only \emph{one choice}, and when the opposites
have been sharply and clearly distinguished from each other, then let each make the choice
honestly and self-consciously!

\sepline

With this I close these lectures. What I have wanted is consciousness,
honesty, clarity. Against illusions and confusion I have
fought; I have striven to distinguish for thought what was mixed together, unconsciously
or deliberately, and in this mixing was distorted.
In order to be able to do this, I have used the freedom of science\kpage{188},
which this university possesses perhaps in a higher
degree than any other university in Europe at this moment; I have
let thought have its say, and presented the attacks on the established order
with the same thoroughness and sharpness as the defense.
But in making full use of this freedom, I have
been conscious of the obligation that it laid upon me. I have
spoken earnestly about what is earnest, I have spoken the strict language of science and of
personal conviction; only against the frivolous
and the thoughtless have I used the weapon of mockery, and only then
when, by calm scientific consideration, it had proved to
have fallen forfeit to it. Perhaps \ed{78.3} at the decisive
point my way parts from that of most of you; but even if that is so, it is
my hope that those whose way is a different one from mine will also now be able to
walk it undisturbed by doubt and less exposed to delusion, and
that those who can follow the way I go will be able to do so with
consciousness and clarity, without discord in the mind, without dispersion in the
spirit.

For the interest with which you have followed these lectures right to the end,
I offer you my sincere thanks. It has
not only been personally dear and encouraging to me, but it has
been a sign to me that interest in the science that I have the
honor to lecture on has not died out, that the studies too which do not belong
to the official circle of the bread-and-butter studies have their devotees. What I have only
been able to intimate in these lectures, which were essentially of a historical-critical kind,
I hope one day, if time and strength are left me,
to find occasion to present in detailed connection.
Until then I ask you to be content with what I have been able to give here,
and to keep a friendly memory of these hours, and of the spiritual
fellowship that we have lived with one another in them: a fellowship during
which I feel that I myself have developed.

Receive, then, in parting, a friendly farewell!
\begin{flushright}(Last lecture, given on 31 May.)\end{flushright}

\sepline

25 May 1859.

\end{document}
